%%%%%%%%%%%%%%%%%%%%%%%%%%%%%%%%%%%%%%%%%
% Masters/Doctoral Thesis
% LaTeX Template
% Version 2.5 (27/8/17)
%
% This template was downloaded from:
% http://www.LaTeXTemplates.com
%
% Version 2.x major modifications by:
% Vel (vel@latextemplates.com)
%
% This template is based on a template by:
% Steve Gunn (http://users.ecs.soton.ac.uk/srg/softwaretools/document/templates/)
% Sunil Patel (http://www.sunilpatel.co.uk/thesis-template/)
%
% Template license:
% CC BY-NC-SA 3.0 (http://creativecommons.org/licenses/by-nc-sa/3.0/)
%
%%%%%%%%%%%%%%%%%%%%%%%%%%%%%%%%%%%%%%%%%

%----------------------------------------------------------------------------------------
%	PACKAGES AND OTHER DOCUMENT CONFIGURATIONS
%----------------------------------------------------------------------------------------

\documentclass[
11pt, % The default document font size, options: 10pt, 11pt, 12pt
%oneside, % Two side (alternating margins) for binding by default, uncomment to switch to one side
english, % ngerman for German
singlespacing, % Single line spacing, alternatives: onehalfspacing or doublespacing
%draft, % Uncomment to enable draft mode (no pictures, no links, overfull hboxes indicated)
%nolistspacing, % If the document is onehalfspacing or doublespacing, uncomment this to set spacing in lists to single
%liststotoc, % Uncomment to add the list of figures/tables/etc to the table of contents
%toctotoc, % Uncomment to add the main table of contents to the table of contents
%parskip, % Uncomment to add space between paragraphs
%nohyperref, % Uncomment to not load the hyperref package
headsepline, % Uncomment to get a line under the header
%chapterinoneline, % Uncomment to place the chapter title next to the number on one line
%consistentlayout, % Uncomment to change the layout of the declaration, abstract and acknowledgements pages to match the default layout
]{MastersDoctoralThesis} % The class file specifying the document structure
\usepackage[utf8]{inputenc} % Required for inputting international characters
\usepackage[T1]{fontenc} % Output font encoding for international characters
\usepackage{tabularx}
\usepackage{subcaption}
\usepackage{mathpazo} % Use the Palatino font by default
\usepackage[backend=biber, style=ieee]{biblatex}

%\usepackage[backend=biber, style=apa, natbib=true]{biblatex}
\DeclareLanguageMapping{english}{english-apa}

\addbibresource{references.bib}  % This adds your .bib file

\usepackage[autostyle=true]{csquotes} % Required to generate language-dependent quotes in the bibliography

%----------------------------------------------------------------------------------------
%	MARGIN SETTINGS
%----------------------------------------------------------------------------------------

\geometry{
	paper=a4paper, % Change to letterpaper for US letter
	inner=2.5cm, % Inner margin
	outer=3.8cm, % Outer margin
	bindingoffset=.5cm, % Binding offset
	top=1.5cm, % Top margin
	bottom=1.5cm, % Bottom margin
	%showframe, % Uncomment to show how the type block is set on the page
}

%----------------------------------------------------------------------------------------
%	THESIS INFORMATION
%----------------------------------------------------------------------------------------

\thesistitle{Optimizing Large Language Models for Low-Latency Inference under Resource Constraints} % Your thesis title, this is used in the title and abstract, print it elsewhere with \ttitle
\degree{Masters of Science in Artificial Intelligence} % Your degree name, this is used in the title page and abstract, print it elsewhere with \degreename
\author{Mohammad Lootah} % Your name, this is used in the title page and abstract, print it elsewhere with \authorname
\addresses{} % Your address, this is not currently used anywhere in the template, print it elsewhere with \addressname

\subject{Biological Sciences} % Your subject area, this is not currently used anywhere in the template, print it elsewhere with \subjectname
\keywords{} % Keywords for your thesis, this is not currently used anywhere in the template, print it elsewhere with \keywordnames
\university{\href{http://www.adu.ac.ae}{Abu Dhabi University}} % Your university's name and URL, this is used in the title page and abstract, print it elsewhere with \univname
\department{\href{http://department.university.com}{Department of Electrical and Computer Engineering}} % Your department's name and URL, this is used in the title page and abstract, print it elsewhere with \deptname
\faculty{\href{https://www.adu.ac.ae/study/colleges/college-of-engineering/faculty}{Faculty of Engineering}} % Your faculty's name and URL, this is used in the title page and abstract, print it elsewhere with \facname

\AtBeginDocument{
\hypersetup{pdftitle=\ttitle} % Set the PDF's title to your title
\hypersetup{pdfauthor=\authorname} % Set the PDF's author to your name
\hypersetup{pdfkeywords=\keywordnames} % Set the PDF's keywords to your keywords
}

\begin{document}

\frontmatter % Use roman page numbering style (i, ii, iii, iv...) for the pre-content pages

\pagestyle{plain} % Default to the plain heading style until the thesis style is called for the body content

%----------------------------------------------------------------------------------------
%	TITLE PAGE
%----------------------------------------------------------------------------------------

\begin{titlepage}
\begin{center}

\vspace*{.02\textheight}
\includegraphics[width=0.45\textwidth]{Figures/adu-logo.jpg}\\[1cm] % University logo
{\scshape\LARGE \univname\par}\vspace{1.5cm} % University name
\textsc{\Large Thesis Report}\\[0.5cm] % Thesis type

\HRule \\[0.4cm] % Horizontal line
{\huge \bfseries \ttitle\par}\vspace{0.4cm} % Thesis title
\HRule \\[1.5cm] % Horizontal line

\begin{minipage}[t]{0.4\textwidth}
\begin{flushleft} \large
\emph{Author:}\\
{\authorname} % Author name - remove the \href bracket to remove the link
\end{flushleft}
\end{minipage}
\begin{minipage}[t]{0.4\textwidth}
\begin{flushright} \large
\emph{Supervisor:} \\
{Dr. Sajid Khawaja} % Supervisor name - remove the \href bracket to remove the link
\end{flushright}
\end{minipage}\\[3cm]

\vfill

\large \textit{A thesis submitted in fulfillment of the requirements\\ for the degree of \degreename}\\[0.3cm] % University requirement text
\textit{in the}\\[0.4cm]
\deptname\\[2cm] % Research group name and department name

\vfill

{\large \today}\\[4cm] % Date


\vfill
\end{center}
\end{titlepage}

%----------------------------------------------------------------------------------------
%	DECLARATION PAGE
%----------------------------------------------------------------------------------------

\begin{declaration}
\addchaptertocentry{\authorshipname} % Add the declaration to the table of contents

\noindent I, \authorname, declare that this report titled, \enquote{\ttitle} and the work presented in it are my own. I confirm that:

\begin{itemize}
\item This work was done wholly or mainly while in candidature for a research degree at this University.
\item Where any part of this report has previously been submitted for a degree or any other qualification at this University or any other institution, this has been clearly stated.
\item Where I have consulted the published work of others, this is always clearly attributed.
\item Where I have quoted from the work of others, the source is always given. With the exception of such quotations, this thesis is entirely my own work.
\item I have acknowledged all main sources of help.
\item Where the report is based on work done by myself jointly with others, I have made clear exactly what was done by others and what I have contributed myself.
\end{itemize}

\noindent Signed:\\
\rule[0.5em]{25em}{0.5pt} % Signature line

\noindent Date:\\
\rule[0.5em]{25em}{0.5pt} % Date line
\end{declaration}


%----------------------------------------------------------------------------------------
%	ABSTRACT PAGE
%----------------------------------------------------------------------------------------

\begin{abstract}
\addchaptertocentry{\abstractname} % Add the abstract to the table of contents
Large language models are costly to serve because their weights and their key--value (KV) cache must fit in GPU memory and because token generation is limited by memory bandwidth. Post-training compression reduces this cost, but published methods are usually applied with one setting for every layer and are judged by simulated perplexity rather than by measurements in the engine that serves the model. This thesis presents a sensitivity-driven framework that addresses both points. Stage~0 measures each decoder layer by the rise in perplexity when the layer is bypassed and writes per-layer plans; Stages~1 to 3 apply published quantization, pruning and KV-cache techniques, each on a uniform plan (Standard) and on the Stage~0 plan (Our method) with the same code; Stage~4 measures the model in Hugging Face Transformers, llama.cpp and vLLM; and a constrained Bayesian search chooses the techniques and settings under limits on perplexity and decode latency. On TinyLlama (1.1 billion parameters), sensitivity is concentrated in five of the 22 layers. Giving 8 bits to the sensitive layers lowers the perplexity of all eleven weight-quantization techniques compared with 4 bits in every layer, for a file of 1.05~GB instead of 0.81~GB, and leaving them unpruned lowers the perplexity of all three pruning techniques at every level. The plan does not help where it spends fewer bits than the uniform setting or, in pruning, at equal size. Pruned models show no gain in size or latency in the engines, and models with rounded activations are 13 to 17 times slower than the uncompressed model in vLLM. The main result is a final model that uses round-to-nearest weight-only quantization (threshold 0.800, groups of 64), no pruning and SnapKV: in vLLM it has 0.898~GB and 2.00~ms per token against 2.200~GB and 2.54~ms per token for the uncompressed model, which is 59 percent smaller and 21 percent faster, and its vLLM perplexity of 11.10 against 10.18 is within both limits. A second final model, OmniQuant weight-only quantization (threshold 0.857) with unstructured pruning at a ratio of 0.082 and InfiniGen (0.838~GB, 1.98~ms per token), is within the perplexity limit in Hugging Face (11.50) but not in vLLM (13.00). Their mean accuracy on four tasks is 21 and 11 percent lower. A stage ablation shows that this loss comes from the KV-cache step, which runs only in Hugging Face, and not from quantization: with quantization alone the mean accuracy is 0.624 and 0.613 against 0.625.
\end{abstract}

\cleardoublepage

\chapter*{Executive Summary}
\addchaptertocentry{Executive Summary}

\noindent\textbf{Problem.} A large language model (LLM) is expensive to run: its weights and its key--value (KV) cache must fit in GPU memory, and the speed of token generation is limited by how fast the weights can be read from memory. Compression methods such as quantization, pruning and KV-cache compression reduce this cost. Two things are missing in the published work reviewed in Chapter~\ref{Chapter3}: the methods are rarely compared under identical conditions, and they are mostly judged by simulated perplexity, not by the size and speed of the model in the inference engine that would serve it.

\medskip\noindent\textbf{What was built.} The thesis implements a five-stage framework and applies it to TinyLlama, a 1.1-billion-parameter model of the Llama family. Stage~0 measures how sensitive each of the 22 decoder layers is and writes per-layer plans. Stages~1, 2 and 3 quantize the model, prune it and compress its KV cache with published techniques. Every technique is run twice with the same code: on a uniform plan (\emph{Standard}) and on the Stage~0 plan (\emph{Our method}). Stage~4 measures each model in Hugging Face Transformers, llama.cpp and vLLM. A constrained Bayesian search then chooses one technique per stage and the plan settings under two limits, a perplexity of at most 11.5 and a vLLM decode latency of at most 2.5~ms per token.

\medskip\noindent\textbf{Main results.}
\begin{itemize}
  \item Sensitivity is very uneven: five layers (0, 1, 2, 7 and 21) carry almost all of it. Measuring it takes about two minutes.
  \item Giving 8 bits to the sensitive layers and 4 bits to the rest lowered the perplexity of all eleven weight-quantization techniques compared with 4 bits everywhere, for a file of 1.05~GB instead of 0.81~GB (uncompressed: 2.20~GB). Leaving the sensitive layers unpruned lowered the perplexity of all three pruning techniques at every pruning level.
  \item The plan is not a gain at equal size: a pruning plan that removes the same total as uniform pruning was worse in 19 of 20 comparisons, and the activation plan, which spends fewer bits than the uniform setting, was worse for eight of twelve techniques.
  \item For weight-only quantization at equal size, the layers chosen by the score gave a lower perplexity than the least sensitive layers in all eight comparisons and than every random choice of layers in six of eight. A plan with 8 and 3 bits that has the size of 4 bits in every layer was much worse (12.55 against 10.97 for round-to-nearest).
  \item Pruning brought no gain in file size or latency in llama.cpp or vLLM, because removed weights are stored as zeros. Models with low-bit activations were small but 13 to 17 times slower than the uncompressed model in vLLM.
  \item The perplexity of a model can change when it is exported: for the same quantized model, without the KV-cache step, 11.25 in Hugging Face and 12.96 in vLLM with OmniQuant, and 10.44 and 11.10 with round-to-nearest.
  \item The search gave two final models. The main result uses round-to-nearest on the weights only (threshold 0.800, groups of 64), no pruning and SnapKV: 0.898~GB and 2.00~ms per token, against 2.200~GB and 2.54~ms per token for the uncompressed model. It is 59 percent smaller and 21 percent faster, and it is within both limits in vLLM (perplexity 11.10 against 10.18). The second uses OmniQuant on the weights only (threshold 0.857), unstructured pruning with a ratio of 0.082 and InfiniGen: 0.838~GB and 1.98~ms per token, with a perplexity of 11.50 in Hugging Face and 13.00 in vLLM. Their mean accuracy on four tasks is 21 and 11 percent lower. A stage ablation shows that this loss comes from the KV-cache step and not from quantization: with quantization alone the mean accuracy is 0.624 and 0.613 against 0.625.
  \item Compared with published work on TinyLlama in Hugging Face, the two final models raise perplexity by 9.2 and 13.0 percent, against 38.6 to 44.4 percent for pruning half of the weights and 6.9 percent for pruning 30 percent; the pruning studies lose less accuracy, on a different set of tasks.
\end{itemize}

\noindent\textbf{What it means.} For a model of this kind on a single GPU, weight-only quantization with more bits in the few sensitive layers is the most effective step; pruning is not worth its quality cost unless the engine can use sparse weights; and a model should be chosen on measurements made in the engine that will serve it, with a task accuracy among the requirements.

\medskip\noindent\textbf{Limits of the work.} The results are for one model and one seed, most techniques are simplified re-implementations, and no TensorRT-LLM measurement is reported. Chapter~\ref{Chapter6} lists the future work that follows from these limits.

\cleardoublepage




%----------------------------------------------------------------------------------------
%	LIST OF CONTENTS/FIGURES/TABLES PAGES
%----------------------------------------------------------------------------------------

\tableofcontents % Prints the main table of contents

\listoffigures % Prints the list of figures

\listoftables % Prints the list of tables



%----------------------------------------------------------------------------------------
%	THESIS CONTENT - CHAPTERS
%----------------------------------------------------------------------------------------

\mainmatter % Begin numeric (1,2,3...) page numbering

\pagestyle{thesis} % Return the page headers back to the "thesis" style

% Include the chapters of the thesis as separate files from the Chapters folder
% Uncomment the lines as you write the chapters

\clearpage
\chapter{Introduction}
\label{Chapter1}

\section{Motivation}
\label{sec:ch1-motivation}

With every iteration of large language model (LLM), the size of a model raises the hardware barrier to run them, because the weights must fit in memory and this leads to slow down in token generation due to the limitation of memory bandwidth \cite{lin2024awq}. In generative inference the output is produced one token at a time, and at small batch sizes and sequence lengths the time needed to load the weights from memory dominates the time of each step \cite{pope2023efficiently}. On top of the weights, the model keeps a key--value (KV) cache of intermediate attention state for every previous token; this cache is stored in GPU memory in addition to the parameters and grows linearly with the sequence length and the batch size \cite{zhang2023h2o}, and it is large enough that the way a serving system manages it limits how many requests can be batched together \cite{kwon2023vllm}.

The size and computational needs of LLMs make them hard to use in practice, especially where resources are limited, and model compression is the research area that addresses this; a recent survey covers quantization, pruning and knowledge distillation \cite{zhu2024survey}. Post-training quantization requires no retraining and no labelled data \cite{nagel2021whitepaper}. LLM.int8() showed that the feed-forward and attention projections of models with up to 175 billion parameters can be run with 8-bit matrix multiplication, halving inference memory, once the few emergent outlier features are handled separately in 16 bits \cite{dettmers2022llmint8}; GPTQ (published as OPTQ) quantizes weights to 3 or 4 bits using approximate second-order information and can process a 175-billion-parameter model in about four GPU hours \cite{frantar2023gptq}; AWQ protects the one percent of weights that matter most, as identified from the activations, and keeps low-bit weight-only quantization accurate without back-propagation \cite{lin2024awq}; SmoothQuant moves the difficulty of quantizing activation outliers into the weights so that both weights and activations can be stored in 8 bits \cite{xiao2023smoothquant}; Wanda prunes a model with a simple weight-times-activation criterion and no weight update \cite{sun2024wanda}; and H2O shrinks the KV cache by keeping only the recent tokens and the few tokens that receive most of the attention \cite{zhang2023h2o}.

This thesis is motivated by three questions: how these methods compare with one another, how much of their reported benefit survives in the inference engine that will actually serve the model, and whether a model should be compressed uniformly or differently in different parts.

\section{Problem Statement}
\label{sec:ch1-problem}

Existing compression methods for LLMs are judged mainly by perplexity, a metric that Jaiswal et al.\ describe as fairly simple and widely questioned \cite{jaiswal2024compressing}, and Li et al.\ state that a comprehensive evaluation of quantized LLMs is essential to guide the selection of a quantization method \cite{li2024evaluating}. The survey that preceded this thesis, written by the author and the supervisor and summarised in Section~\ref{sec:ch1-scope}, reviewed about fifty papers and identified five gaps that the present work takes as its starting point. First, there is no practical comparison of quantization precision levels under identical conditions. Second, there are no clear rules for choosing between quantization-only, pruning-only and combined pipelines. Third, pruning reduces the number of parameters but does not always reduce real latency. Fourth, there is little deployment-aware evaluation of the same compressed model across different inference tools such as Hugging Face Transformers, llama.cpp, vLLM and TensorRT-LLM. Fifth, there is little adaptive compression that chooses the bit width, pruning ratio or compression level according to a latency, memory or accuracy requirement.

A further, more specific problem underlies the fifth gap. Different layers of an LLM vary in their sensitivity to compression \cite{yuan2025asvd}, and a method that could measure how much each layer tolerates could compress the tolerant layers more and the sensitive layers less. Whether such a per-layer plan actually improves the trade-off between quality, size and speed, and whether a cheap sensitivity measurement is a good enough guide, has not been established under controlled conditions.

The problem addressed by this thesis is therefore to build a single framework in which published compression techniques for weights, activations and the KV cache can be run with and without a sensitivity-driven per-layer plan, under one protocol, and in which every resulting model is measured not only in a simulation of the compressed arithmetic but also in the inference engines that would deploy it, so that the five gaps above can be examined with directly comparable evidence.

\section{Aims and Objectives}
\label{sec:ch1-aims}

The aim of this thesis is to determine whether sensitivity-driven, per-layer compression planning improves the quality, size and latency trade-off of post-training compression for a decoder-only LLM, and to characterise how compression choices translate into real deployment metrics across inference engines. The aim is pursued through the following research questions.

\begin{enumerate}
  \item \textbf{RQ1.} How unevenly is sensitivity to compression distributed across the layers of a decoder-only LLM, and does a cheap measured score, namely the rise in perplexity when a layer is bypassed, rounded to fewer bits, or has its KV cache reduced, predict the damage done by compressing that layer?
  \item \textbf{RQ2.} Under identical calibration data, seed, evaluation data, hardware and software, does a sensitivity-driven per-layer plan that assigns 8 bits to sensitive layers and 4 bits to the rest improve the quality, size and latency trade-off of published weight, activation and KV-cache compression techniques compared with their usual uniform settings?
  \item \textbf{RQ3.} When is quantization-only, pruning-only or a combined pipeline preferable, given that pruned weights are stored as zeros and therefore bring no size or latency gain in the inference engines studied here?
  \item \textbf{RQ4.} How do simulated compression results obtained in Hugging Face Transformers carry over to real file size, memory and decode latency in llama.cpp and vLLM?
  \item \textbf{RQ5.} Can a constrained multi-objective search over techniques and their settings find configurations that satisfy a deployment requirement on perplexity and decode latency while minimising perplexity, latency and model size?
\end{enumerate}

To answer these questions the following objectives were set.

\begin{enumerate}
  \item \textbf{O1.} Build a per-layer sensitivity profiler and planner that measures the effect of bypassing, rounding and KV-cache reduction for each decoder layer, and that produces weight, activation, joint weight--activation, pruning and KV-cache plans together with a prediction of the memory each plan would need.
  \item \textbf{O2.} Implement or adapt published quantization, pruning, low-rank and KV-cache techniques so that each can be run with and without the plan on identical code.
  \item \textbf{O3.} Define one comparison protocol, consisting of the FP16 model, each technique at its usual uniform setting, and each technique following the plan, with the same calibration data, seed, evaluation data, hardware and software throughout.
  \item \textbf{O4.} Export the compressed models to llama.cpp and vLLM and measure its real file size, memory use and latency in each engine.
  \item \textbf{O5.} Implement a constrained multi-objective search that selects the technique and settings for each compression stage under a deployment requirement, and report its outcomes.
  \item \textbf{O6.} Analyse the results against the five gaps and derive guidance on precision levels, pipeline choice, pruning, engine behaviour and adaptive planning for practitioners.
\end{enumerate}

\section{Scope and Limitations}
\label{sec:ch1-scope}

The work builds on the survey by the author and the supervisor, which reviewed around fifty papers published from 2020 onwards in five groups (quantization, pruning and sparsity, hybrid quantization and pruning, knowledge distillation, and KV-cache compression) and which proposed a sensitivity-driven, adaptive framework. The thesis implements and evaluates that framework. The scope is deliberately bounded as follows.

Resource constraints in this thesis mean that the deployment needs a small model and a short time per token, which the search receives as limits and objectives; the GPUs of the experiments serve only as test machines.

The thesis considers post-training compression only. No full retraining, quantization-aware training or knowledge distillation is performed, and no custom hardware or new GPU kernels are developed. The techniques studied are published methods, most of them re-implemented in simplified form within the framework so that they can follow a per-layer plan; the aim is a controlled comparison of the planning idea, not a reproduction of each method's best published number. These simplified versions were not compared against the original implementations, so their results may differ from the published ones.

A single model is used throughout: TinyLlama, a 1.1-billion-parameter model that follows the architecture and tokenizer of Llama~2 \cite{zhang2024tinyllama}. It was chosen because it is small enough to be compressed, exported and measured many hundreds of times on the available hardware while still being a genuine decoder-only LLM of the Llama family. Findings may not transfer unchanged to other models, in particular to much larger ones.

Each configuration is run with one seed, so the variation between runs with different seeds is not measured. Quality is measured by perplexity on WikiText-2 and, for selected final models, by zero-shot accuracy on four common-sense benchmarks. Deployment metrics are measured in three engines: Hugging Face Transformers \cite{wolf2020transformers}, which evaluates the quality of the simulated compression, and llama.cpp \cite{ggml2025llamacpp} and vLLM \cite{kwon2023vllm}, which give real file size, memory and speed after re-rounding the weights onto their own grids. The framework can also export a model to TensorRT-LLM \cite{nvidia2025tensorrtllm}, but no TensorRT-LLM measurement is reported in this thesis. Perplexity measured in an engine can differ from perplexity measured in Hugging Face (for one quantized model in this work, 11.25 in Hugging Face and 12.96 in vLLM), so the last two searches judge perplexity in vLLM. The hardware is an NVIDIA A100 GPU on Google Colab and the host PC with an NVIDIA RTX~5070~Ti 16~GB GPU. Speed is measured on these two GPUs only, and the ranking of methods by speed may differ on other hardware.

Several techniques were attempted and dropped because they could not be made to run reliably in the framework, namely AQLM \cite{egiazarian2024aqlm}, QTIP \cite{tseng2024qtip} and ABQ-LLM \cite{zeng2025abqllm}. BiLLM \cite{huang2024billm} is included without a planned variant, and the planned variant of RPTQ \cite{yuan2023rptq} was dropped. Pruned weights are stored as zeros rather than in a sparse format, so the pruning experiments measure quality only; they cannot show size or latency gains in the engines, and the thesis does not claim any. Finally, the multi-objective search is reported as a set of outcomes rather than as a pass or fail, since there is no single correct answer to a constrained trade-off.

\section{Structure of the Thesis}
\label{sec:ch1-structure}

The remainder of the thesis is organised as follows. Chapter~\ref{Chapter2} gives the background needed to read the later chapters: how a language model is measured, the transformer and the decoder-only architecture, how inference works and why decoding is limited by memory bandwidth, the Llama family and in particular TinyLlama, the main compression techniques, the inference engines, the evaluation metrics, and the basics of multi-objective optimisation. Chapter~\ref{Chapter3} reviews the published literature in groups, compares the groups, and states the research gaps that this thesis responds to. Chapter~4 describes the framework and the experimental method, Chapter~5 reports and discusses the results, and Chapter~6 concludes and suggests future work.
\clearpage
\clearpage
\chapter{Background}
\label{Chapter2}

This chapter gives the technical background needed to read the rest of the thesis. It starts with how the quality of a language model is measured, then describes the transformer and the decoder-only architecture used by TinyLlama, explains how inference proceeds and why generating text is limited by memory bandwidth, introduces the Llama family and the particular model used in the experiments, and surveys the main families of compression techniques. It ends with the inference engines, the evaluation metrics and the multi-objective optimisation methods that the framework relies on.

\section{Language Models and Perplexity}
\label{sec:ch2-lm}

The LLMs considered in this thesis are autoregressive: the probability of a sequence of tokens $x_1, \dots, x_T$ is factorised as a product of conditional next-token probabilities,
\begin{equation}
  p(x_1, \dots, x_T) = \prod_{t=1}^{T} p(x_t \mid x_1, \dots, x_{t-1}),
  \label{eq:ch2-ar}
\end{equation}
each token being conditioned on the tokens that precede it.

Quality is measured by perplexity. Given a test sequence of $N$ tokens, perplexity is the exponential of the average negative log-likelihood that the model assigns to the tokens,
\begin{equation}
  \mathrm{PPL} = \exp\left( -\frac{1}{N} \sum_{t=1}^{N} \log p(x_t \mid x_{<t}) \right),
  \label{eq:ch2-ppl}
\end{equation}
so a lower value means that the model assigns a higher likelihood to the text. In this thesis the WikiText-2 test split is cut into 512-token windows and the perplexity is computed over them.

\section{The Transformer and Decoder-only LLMs}
\label{sec:ch2-transformer}

The transformer is a sequence-to-sequence model that consists of an encoder and a decoder, each of which is a stack of $L$ identical blocks; when only the decoder is used, the model is typically used for sequence generation such as language modelling \cite{lin2021survey}. TinyLlama has such a decoder-only architecture \cite{zhang2024tinyllama}. Each block contains a multi-head self-attention module and a position-wise feed-forward network, with a residual connection around each module and a normalisation layer \cite{lin2021survey}. The following paragraphs describe these components as they appear in the Llama family and in TinyLlama.

\subsection{Self-attention, Multi-head and Grouped-query Attention}
\label{sec:ch2-attention}

The transformer uses an attention mechanism with a query--key--value model \cite{lin2021survey}. In its scaled dot-product attention, with queries $\mathbf{Q}$, keys $\mathbf{K}$ and values $\mathbf{V}$, the output is
\begin{equation}
  \mathrm{Attention}(\mathbf{Q}, \mathbf{K}, \mathbf{V}) = \mathrm{softmax}\left( \frac{\mathbf{Q}\mathbf{K}^{\top}}{\sqrt{d_k}} \right) \mathbf{V},
  \label{eq:ch2-attn}
\end{equation}
where $d_k$ is the dimension of the keys and the division by $\sqrt{d_k}$ alleviates the vanishing-gradient problem of the softmax function \cite{lin2021survey}. In a decoder the self-attention is adapted so that each position cannot attend to subsequent positions \cite{lin2021survey}. Multi-head attention projects the queries, keys and values with $H$ different sets of learned projections, applies attention to each set, and concatenates and projects the results \cite{lin2021survey}.

Two later variants reduce the number of key and value heads. Multi-query attention (MQA) uses only a single key--value head and greatly speeds up decoder inference, but it can lead to a loss of quality \cite{ainslie2023gqa}. Grouped-query attention (GQA) is a generalisation that uses an intermediate number of key--value heads, more than one but fewer than the number of query heads; Ainslie et al.\ showed that a model uptrained to GQA reaches quality close to multi-head attention at a speed comparable to MQA \cite{ainslie2023gqa}. TinyLlama uses GQA with 32 query heads and 4 groups of key--value heads \cite{zhang2024tinyllama}, which matters for compression because the size of the KV cache depends on the number of key--value heads rather than the number of query heads (Section~\ref{sec:ch2-kvcache}).

\subsection{Positional Encoding and RoPE}
\label{sec:ch2-rope}

Because the transformer contains no recurrence or convolution, it has no information about token positions and needs an additional positional representation \cite{lin2021survey}. The Llama family uses rotary position embedding (RoPE) \cite{touvron2023llama}. RoPE encodes the absolute position with a rotation matrix and at the same time incorporates the relative position dependency into the self-attention formulation; its authors also note that it gives flexibility in sequence length and a decaying dependency between tokens as their distance grows \cite{su2023roformer}. Llama removes the absolute positional embeddings and adds rotary embeddings at each layer \cite{touvron2023llama}, and TinyLlama also uses RoPE \cite{zhang2024tinyllama}. RoPE also matters for KV-cache compression: KVQuant quantizes the keys before the rotary embedding is applied \cite{hooper2024kvquant}.

\subsection{Feed-forward Network and SwiGLU}
\label{sec:ch2-ffn}

The second sub-layer of each block is a position-wise feed-forward network, which operates separately and identically on each position \cite{lin2021survey}. Gated linear units (GLU) consist of the element-wise product of two linear projections, one of which first passes through a sigmoid function, and variants use other functions in place of the sigmoid; Shazeer tested such variants in the transformer feed-forward sub-layer and found that some of them improve quality over ReLU and GELU \cite{shazeer2020glu}. The variant that uses the Swish function is called SwiGLU \cite{shazeer2020glu}. Llama replaces the ReLU non-linearity with SwiGLU and uses a hidden dimension of $\frac{2}{3} \cdot 4d$ instead of $4d$ \cite{touvron2023llama}. TinyLlama adopts SwiGLU with an intermediate size of 5{,}632 for a hidden size of 2{,}048 \cite{zhang2024tinyllama}. In the framework of this thesis the structured pruning stage removes channels of this intermediate dimension.

\subsection{Normalisation and the Pre-norm Block}
\label{sec:ch2-norm}

In the original post-norm design the layer normalisation is placed between the residual blocks; in the pre-norm design it is placed inside the residual blocks \cite{xiong2020layernorm}. Xiong et al.\ showed that the gradients of a pre-norm transformer are well behaved at initialisation, which allows training without the learning-rate warm-up stage \cite{xiong2020layernorm}. Llama normalises the input of each sub-layer rather than its output, using RMSNorm, to improve training stability \cite{touvron2023llama}, and TinyLlama also uses pre-norm with RMSNorm \cite{zhang2024tinyllama}. A decoder layer of this kind therefore computes
\begin{equation}
  \mathbf{h} = \mathbf{x} + \mathrm{Attn}(\mathrm{RMSNorm}(\mathbf{x})), \qquad
  \mathbf{y} = \mathbf{h} + \mathrm{FFN}(\mathrm{RMSNorm}(\mathbf{h})),
  \label{eq:ch2-block}
\end{equation}
where $\mathbf{x}$ is the layer input and $\mathbf{y}$ its output. In the sensitivity profiling of this thesis a layer is bypassed by setting its output equal to its input.

\section{Inference}
\label{sec:ch2-inference}

\subsection{Prefill and Decode}
\label{sec:ch2-prefill}

Pope et al.\ divide inference into two phases \cite{pope2023efficiently}. In the first phase, which they call prefill, the input tokens present at the start of the request are processed in a single forward pass, all in parallel \cite{pope2023efficiently}. In the second phase, decode, the output tokens are generated autoregressively in a sequential loop, one forward pass per step \cite{pope2023efficiently}.

\subsection{The KV Cache}
\label{sec:ch2-kvcache}

The key--value (KV) cache stores the attention keys and values of earlier tokens to avoid recomputing them \cite{liu2024kivi}. It is held in GPU memory in addition to the model parameters, and it grows linearly with the sequence length and the batch size \cite{zhang2023h2o}. For a model with $L$ layers, $n_{kv}$ key--value heads of dimension $d_h$, a batch of $B$ sequences of length $S$ and $b$ bytes per element, the cache holds
\begin{equation}
  \mathrm{KV\;bytes} = 2 \cdot L \cdot n_{kv} \cdot d_h \cdot B \cdot S \cdot b,
  \label{eq:ch2-kvsize}
\end{equation}
the factor 2 counting keys and values. GQA lowers $n_{kv}$; TinyLlama uses it to reduce the memory-bandwidth overhead and speed up inference \cite{zhang2024tinyllama}. For TinyLlama, with $L = 22$, $n_{kv} = 4$ and $d_h = 64$, a single 2{,}048-token sequence in 16-bit precision needs about 46~MB of cache, small next to the 2.2~GB of weights. For long contexts and large batches the cache can exceed the weights: Pope et al.\ report that for a model of more than 500 billion parameters with multi-head attention, at batch size 512 and context length 2{,}048, the cache totals three times the size of the parameters \cite{pope2023efficiently}. How the cache is managed also matters: Kwon et al.\ observed that the cache memory of each request is large and grows and shrinks dynamically and that, when it is managed inefficiently, fragmentation and redundant duplication waste enough memory to limit the batch size, which motivated PagedAttention in vLLM \cite{kwon2023vllm}.

\subsection{Why Decoding is Memory-bandwidth Bound}
\label{sec:ch2-roofline}

Pope et al.\ state that the weights and the KV cache must be transferred from memory to the compute cores once per forward pass: at small batch sizes and sequence lengths the time to load the weights dominates, and the KV cache likewise has to be loaded from off-chip memory once for every token generated \cite{pope2023efficiently}. Liu et al.\ make the same observation for the cache, noting that loading the KV cache leaves the computational core idle, which limits the inference speed \cite{liu2024kivi}.

This matters for compression. Lin et al.\ identify the model size as what slows down token generation, because generation is limited by memory bandwidth, and weight quantization reduces that size \cite{lin2024awq}. Quantizing the KV cache likewise decreases the total bytes that it takes \cite{liu2024kivi}. Pruning individual weights, by contrast, is not hardware friendly and receives limited speed gains \cite{zhou2021nm}; this point returns in Section~\ref{sec:ch2-pruning} and in the pruning results of this thesis.

\section{Selected LLM Architectures}
\label{sec:ch2-archs}

Llama is a family of openly released foundation models from 7 to 65 billion parameters trained only on publicly available data \cite{touvron2023llama}. Its architecture combines the components described above: pre-normalisation with RMSNorm, SwiGLU feed-forward layers and rotary embeddings \cite{touvron2023llama}. Llama~2 ranges from 7 to 70 billion parameters and adds the Llama~2-Chat models fine-tuned for dialogue \cite{touvron2023llama2}.

\subsection{TinyLlama}
\label{sec:ch2-tinyllama}

TinyLlama is the model used throughout this thesis. Zhang et al.\ describe it as a compact 1.1-billion-parameter language model that builds on the architecture and tokenizer of Llama~2 and was pretrained on around one trillion tokens for up to three epochs \cite{zhang2024tinyllama}. Its configuration, taken from Table~1 of the technical report, is a hidden size of 2{,}048, an intermediate (feed-forward) size of 5{,}632, a context length of 2{,}048, 32 attention heads, 22 layers and a vocabulary of 32{,}000 tokens; the attention has 32 query heads and 4 groups of key--value heads, the positional information is injected with RoPE, pre-norm RMSNorm is used, the feed-forward network is SwiGLU, and training used FlashAttention-2 \cite{zhang2024tinyllama}. The chat version, TinyLlama-1.1B-Chat-v1.0, is the checkpoint used in the experiments. With about 1.1 billion parameters stored in 16-bit floating point the model occupies about 2.2~GB, which is the FP16 reference size against which compressed sizes are compared in Chapter~5.

TinyLlama was chosen for two reasons. It is a Llama-architecture decoder with GQA, RoPE, RMSNorm and SwiGLU \cite{zhang2024tinyllama}. And it is small enough that a full cycle of calibration, compression, export to three engines and evaluation can be repeated hundreds of times on one GPU, which the comparison protocol of this thesis requires.

\section{Model Compression}
\label{sec:ch2-compression}

Model compression transforms a large, resource-intensive model into a compact version suitable for deployment on resource-constrained devices \cite{zhu2024survey}. The survey by Zhu et al.\ covers quantization, pruning and knowledge distillation for LLMs and also discusses the benchmarks and metrics used to evaluate compressed models \cite{zhu2024survey}. This section explains quantization and pruning, adds low-rank factorisation, KV-cache compression and mixed precision, and defines the terms used later.

\subsection{Quantization}
\label{sec:ch2-quant}

Quantization stores the weights and activation tensors in a lower bit precision than the 16 or 32 bits in which they are usually trained \cite{nagel2021whitepaper}. Uniform affine (asymmetric) quantization is defined by three parameters, a scale factor $s$, an integer zero-point $z$ and a bit width $b$; a real value $x$ is mapped to the integer grid by dividing by the scale, rounding, adding the zero-point and clamping to $[0, 2^b - 1]$, and it is approximately recovered by de-quantization, which subtracts the zero-point and multiplies by the scale \cite{nagel2021whitepaper}:
\begin{equation}
  x_{\mathrm{int}} = \mathrm{clamp}\left( \left\lfloor \frac{x}{s} \right\rceil + z,\; 0,\; 2^b - 1 \right), \qquad
  \hat{x} = s \,(x_{\mathrm{int}} - z).
  \label{eq:ch2-quant}
\end{equation}
The symmetric quantizer restricts the zero-point to zero \cite{nagel2021whitepaper}. The parameters can be shared by a whole tensor (per-tensor quantization) or chosen separately for each output channel of a weight matrix (per-channel quantization) \cite{nagel2021whitepaper}. The parameters can also be set per group of weights: Frantar et al.\ describe standard grouping as applying independent quantization to groups of $g$ consecutive weights, and report that group sizes of 1{,}024 and 128 improve the perplexity of 4-bit models, with a group size of 128 costing about 0.15 extra bits per weight for the additional scale parameters \cite{frantar2023gptq}. The group size is one of the settings that the search in this thesis chooses for GPTQ.

Nagel et al.\ note that moving from 32 to 8 bits divides the memory needed to store the tensors by four \cite{nagel2021whitepaper}.

Two distinctions are used in this thesis. The first is between weight-only quantization, in which only the stored weights are integers and they are de-quantized to 16-bit floating point before the multiplication, and weight--activation quantization, in which the activations are also quantized so that the matrix multiplication itself runs in low precision. Weight-only quantization reduces model size and the bytes moved per token, which is what matters for memory-bound decoding, and methods such as GPTQ \cite{frantar2023gptq} and AWQ \cite{lin2024awq} belong to this class. In weight--activation quantization the activations are quantized as well: Dettmers et al.\ found that large transformers develop emergent outlier features that dominate attention and prediction, and LLM.int8() handles them by isolating the outlier dimensions into a separate 16-bit multiplication \cite{dettmers2022llmint8}; SmoothQuant instead migrates the quantization difficulty from the activations into the weights with a mathematically equivalent per-channel scaling so that both can be stored in 8 bits \cite{xiao2023smoothquant}; and QuaRot rotates the model so that the outliers disappear from the hidden state without changing the output, after which all matrix multiplications can run in 4 bits \cite{ashkboos2024quarot}. The notation W$x$A$y$ denotes $x$-bit weights and $y$-bit activations; W4A16 is 4-bit weight-only quantization and W8A8 is 8-bit weights and activations.

The second distinction is between post-training quantization (PTQ), which requires no retraining and no labelled data, and quantization-aware training (QAT), which requires fine-tuning and labelled training data and can approximate the gradient of the rounding operation with the straight-through estimator \cite{nagel2021whitepaper}. Nagel et al.\ observe that PTQ is in most cases sufficient for 8-bit quantization with close to floating-point accuracy, while QAT enables lower bit widths \cite{nagel2021whitepaper}. GPTQ quantizes a 175-billion-parameter model in about four GPU hours using approximate second-order information and 128 random 2{,}048-token segments of C4 as calibration data \cite{frantar2023gptq}, whereas EfficientQAT, a QAT method designed to be affordable, still needs 41 hours on an A100 for a 2-bit 70-billion-parameter model \cite{chen2025efficientqat}. This thesis uses PTQ only. Round-to-nearest (RTN) applies Equation~\ref{eq:ch2-quant} directly; the other methods used in this thesis are reviewed in Chapter~\ref{Chapter3}.

\subsection{Pruning and Sparsity}
\label{sec:ch2-pruning}

Pruning removes weights. Wanda removes the weights with the smallest product of weight magnitude and input activation norm, per output, with no weight update at all \cite{sun2024wanda}. The Wanda criterion is the unstructured pruning criterion implemented in this thesis.

Zhou et al.\ distinguish unstructured fine-grained sparsity, which zeroes individual weights and achieves high compression ratios but is not hardware friendly and therefore gains little speed, from structured coarse-grained sparsity, which removes whole blocks \cite{zhou2021nm}. Unstructured pruning with zeros left in place, as in the pruning stage of this thesis, reduces neither the stored size nor the decode time in engines that store dense tensors; it is used here to study quality and the effect of per-layer pruning ratios.

\subsection{Low-rank Factorisation}
\label{sec:ch2-lowrank}

A weight matrix $\mathbf{W}$ of size $m \times n$ can be approximated by a product $\mathbf{A}\mathbf{B}$ with $\mathbf{A}$ of size $m \times r$ and $\mathbf{B}$ of size $r \times n$; if the rank $r$ is small enough that $r(m + n) < mn$, the factored form has fewer parameters. Yuan et al.\ find that the difficulty of low-rank decomposition of LLM weights comes from the variance of the activation distribution and from the different sensitivity of different layers. Their method, ASVD, therefore transforms the weight matrix according to the activation distribution before decomposing it, so that activation outliers are absorbed into the transformed weights, and calibrates the decomposition of each layer separately; it needs no training and compresses a network by 10 to 30 percent \cite{yuan2025asvd}. The activation-aware decomposition in the pruning stage of this thesis follows the ASVD idea.

\subsection{KV-cache Compression}
\label{sec:ch2-kvcompress}

Two families of KV-cache compression methods are used in this thesis. Rounding methods quantize the cached keys and values: KVQuant quantizes the keys per channel and before RoPE, uses non-uniform data types weighted by per-layer sensitivity, and keeps outliers in a sparse format, reaching less than 0.1 perplexity degradation at 3 bits on Llama models \cite{hooper2024kvquant}; and QuaRot's rotation also applies to the cache, which it quantizes to 4 bits together with the weights and activations \cite{ashkboos2024quarot}. Eviction methods instead drop entries. H2O keeps a balance of recent tokens and the heavy-hitter tokens that contribute most of the attention value; with 20 percent heavy hitters it improves throughput over three inference systems by up to 29, 29 and 3 times on OPT-6.7B and OPT-30B and reduces latency by up to 1.9 times at the same batch size \cite{zhang2023h2o}; SnapKV selects, per attention head, the prompt positions that an observation window at the end of the prompt attends to, and reports 3.6 times faster generation and 8.2 times better memory efficiency for 16{,}000-token inputs \cite{li2024snapkv}; and InfiniGen, aimed at offloading systems where the cache lives in host memory, speculates which entries the next layer will need by a minimal rehearsal with part of its query weights and key cache, and prefetches only those \cite{lee2024infinigen}. The KV-cache stage of this thesis implements both families and plans, per layer, the bits for keys and values and the share of tokens to keep.

\subsection{Mixed Precision and Sensitivity}
\label{sec:ch2-mixed}

Several LLM compression methods treat parts of the model differently according to their sensitivity: SqueezeLLM assigns non-uniform quantization levels using second-order sensitivity and stores outliers and sensitive weights in a sparse format \cite{kim2024squeezellm}, KVQuant derives per-layer sensitivity-weighted data types for the cache \cite{hooper2024kvquant}, and Shao et al.\ allocate sparsity across layers according to Hessian-based sensitivity \cite{shao2024sensitivity}.

The framework of this thesis separates the question of where to compress from the question of how. Its sensitivity score is not a Hessian quantity but a direct measurement: the rise in calibration perplexity when a layer is bypassed, when its weights or activations are rounded to a given bit width, or when its KV cache is reduced. Gradient-based scores (gradient times weight, Fisher, Taylor, Hessian diagonal and movement scores) were also implemented, but the measured score is the default. The scores are rank-normalised to the interval $[0, 1]$ and turned into plans: 8-bit weights for sensitive layers and 4-bit for the others, an activation plan, a joint weight--activation plan chosen by dynamic programming under an average-bit budget, a pruning ratio per layer that never prunes the five most sensitive layers, and a KV-cache plan. Chapter~4 describes these plans in full.

\section{Inference Engines}
\label{sec:ch2-engines}

Four inference engines are used in this thesis.

Hugging Face Transformers is an open-source library that provides state-of-the-art transformer architectures under a unified API together with a collection of pretrained models \cite{wolf2020transformers}. In this thesis it is the reference environment: compressed models are evaluated in Transformers with simulated quantization, that is, with the quantized values de-quantized back to 16-bit floating point before the arithmetic, so the quality it measures does not depend on any kernel. Because the inference engines convert the weights into their own quantization types, the quality measured in an engine can differ from the quality measured in Transformers.

llama.cpp is a plain C/C++ implementation of LLM inference without dependencies whose stated goal is inference with minimal setup and high performance on a wide range of hardware; it supports integer quantization from 1.5 to 8 bits for faster inference and reduced memory, custom CUDA kernels for NVIDIA GPUs, and hybrid CPU--GPU execution for models that do not fit in GPU memory \cite{ggml2025llamacpp}. Its models are stored in the GGUF format, a binary single-file format designed for fast loading that contains all the information needed to load the model \cite{ggml2025gguf}. Exporting to GGUF converts the weights into one of llama.cpp's own quantization types, so the file size and speed measured in llama.cpp reflect that type, not the simulated one.

vLLM is the serving system built around PagedAttention, an attention algorithm that keeps the KV cache in non-contiguous paged memory in the manner of virtual memory, which removes almost all cache waste and lets the cache be shared within and across requests; its authors report two to four times the throughput of earlier systems at the same latency \cite{kwon2023vllm}. vLLM loads quantized checkpoints in several formats; in this thesis the GPTQ format and the compressed-tensors format are used. TensorRT-LLM is NVIDIA's inference library, which optimises LLM inference with specialised kernels, an efficient runtime and a Python framework for customisation \cite{nvidia2025tensorrtllm}; it builds a TensorRT engine for the model. Both vLLM and TensorRT-LLM re-round the exported weights onto their own quantization grids, so the per-layer plan survives export only to the extent that the engine's format can represent different bit widths in different layers. Measuring the same model in all four engines is what makes the deployment-aware evaluation of this thesis possible.

\section{Evaluation Metrics}
\label{sec:ch2-metrics}

Four kinds of measurement are used. Quality is measured first by perplexity on WikiText-2 (Section~\ref{sec:ch2-lm}), a dataset also used to evaluate GPTQ \cite{frantar2023gptq}, computed over 512-token windows; the first half of the windows is used for every decision in the framework and the second half is held out and never used for tuning. Quality is measured second, for selected final models, by zero-shot accuracy on four common-sense reasoning tasks: ARC-Easy, PIQA, HellaSwag and WinoGrande, which are among the tasks used to evaluate TinyLlama \cite{zhang2024tinyllama}. PIQA tests physical common sense, for example whether a cotton swab or a toothpick is the right tool for a task; humans answer its questions with 95 percent accuracy while large pretrained models reached about 77 percent when it was introduced \cite{bisk2020piqa}. WinoGrande is a large-scale dataset of 44{,}000 problems inspired by the Winograd Schema Challenge and built with systematic bias reduction \cite{sakaguchi2020winogrande}.

Size is measured as the file size of the exported model and as the GPU memory used when it is loaded, and the framework also predicts the memory of a plan before any model is built. Speed is measured as decode latency per token at batch size one, which is the quantity that the memory-bandwidth argument of Section~\ref{sec:ch2-roofline} governs, and, where the engine reports it, as throughput in tokens per second. This thesis reports latency, since its target is a single request on a constrained device.

\section{Multi-objective Optimisation}
\label{sec:ch2-moo}

Choosing a compression configuration means trading quality against size and speed, so the search stage of the framework is a constrained multi-objective optimisation problem. pymoo is a Python framework for multi-objective optimisation \cite{blank2020pymoo}, which the framework of this thesis uses for its NSGA-III option.

Each evaluation is expensive here, because every candidate requires compressing, exporting and measuring a model. Bayesian optimisation is a global optimisation technique for expensive black-box functions: it fits a surrogate model to the evaluations made so far and uses an acquisition function, which trades off exploitation and exploration, to generate the next candidate \cite{bischl2021hpo}. The framework of this thesis uses the Gaussian-process sampler of Optuna; the deployment requirement is expressed as limits on perplexity and on vLLM decode latency, and within those limits the search minimises perplexity, latency and model size. Successive halving \cite{bischl2021hpo} and NSGA-III were implemented in the framework, but the search results reported in Chapter~5 come from the Bayesian optimiser.

\section{Summary}
\label{sec:ch2-summary}

TinyLlama is a 22-layer, 1.1-billion-parameter decoder-only model built from pre-norm transformer blocks with grouped-query attention, rotary embeddings, SwiGLU feed-forward networks and RMSNorm. Text is generated one token at a time, and at small batch sizes each step is limited by the bytes of weights and KV cache that must be read from memory rather than by arithmetic, which is why reducing bit width speeds up decoding while leaving zeros in place does not. Quantization, pruning, low-rank factorisation and KV-cache compression each address a different part of this cost, and several methods treat parts of the model differently according to their sensitivity. Whether a compressed model is actually smaller or faster depends on the engine that runs it, so this thesis measures every model in Hugging Face Transformers, llama.cpp, vLLM and TensorRT-LLM. Finally, choosing a configuration under a deployment requirement is a constrained multi-objective problem, for which the framework uses Gaussian-process Bayesian optimisation. The next chapter reviews the published methods in each family and identifies the gaps that the thesis addresses.
\clearpage
\clearpage
% Chapter 3: Literature Review
% Text, tables and figure taken word for word from the survey paper
% "Survey: Optimizing Large Language Models for Low-Latency Inference under
% Resource Constraints" (M. Lootah and S. G. Khawaja), with its own 50 references.

\chapter{Literature Review}
\label{Chapter3}

\section{Introduction}
\label{sec:lr-intro}

Large Language Models (LLMs) have become a major direction in modern artificial intelligence due to their strong performance in language understanding, text generation, reasoning, summarization, and code-related tasks. Their success is mainly driven by transformer-based architectures, which allow models to learn complex contextual relationships from large-scale data. However, the same architectural growth that improves model capability also increases computational cost, memory usage, and inference latency. As a result, deploying LLMs in practical environments remains challenging, especially when the target system has limited hardware resources or strict response-time requirements \cite{survey2021transformers,survey2021ondevice}.

The deployment challenge becomes more significant in low-latency and resource-constrained environments such as edge devices, embedded platforms, mobile systems, and private enterprise servers. In these settings, large models often require high memory bandwidth, powerful GPUs, and considerable energy consumption, which may not always be available. Therefore, efficient LLM deployment requires optimization techniques that reduce model size and inference cost while maintaining acceptable model accuracy. This has motivated increasing research interest in quantization, pruning, sparsity, knowledge distillation, and hybrid compression strategies \cite{frantar2023sparsegpt,ma2023llmpruner,sanh2020distilbert,jiao2020tinybert}.

Quantization is one of the most widely studied approaches for improving LLM efficiency. It reduces the numerical precision of model weights and activations, allowing models to consume less memory and execute faster on supported hardware. Early and influential methods such as LLM.int8(), ZeroQuant, GPTQ, and SmoothQuant demonstrated that post-training quantization can significantly reduce inference cost while preserving model quality \cite{dettmers2022llmint8,yao2022zeroquant,frantar2023gptq,xiao2023smoothquant}. More recent methods, including AWQ, OmniQuant, SpQR, SqueezeLLM, and QuIP, further improve low-bit and weight-only quantization by addressing activation sensitivity, outlier effects, and dense-sparse compression \cite{lin2024awq,dettmers2024spqr,shao2024omniquant,kim2024squeezellm,chee2023quip}.

In addition to quantization, pruning-based compression has also been explored to reduce unnecessary parameters, layers, or structures within LLMs. Methods such as SparseGPT, Wanda, LLM-Pruner, SliceGPT, and ShortGPT show that large models often contain redundancy that can be removed with limited accuracy loss when pruning is carefully designed \cite{frantar2023sparsegpt,ma2023llmpruner,sun2024wanda,ashkboos2024slicegpt,men2024shortgpt}. Other recent works combine pruning with quantization or use mixed-precision strategies to better balance compression ratio, latency reduction, and model performance \cite{xu2026daypq,cho2025sub4bit,wang2025ofqllm}.

Despite these advances, efficient LLM deployment remains an open research problem. Many techniques report strong compression or accuracy results, but their real-world effectiveness depends heavily on hardware support, memory access patterns, batch size, sequence length, and inference framework implementation. In addition, reducing bit-width or removing model structures can introduce accuracy degradation, instability, or limited generalization across different model families. These issues highlight the need for clearer comparisons between compression ratio, latency, energy efficiency, memory usage, and model quality.

This survey reviews recent quantization and compression techniques for low-latency LLM deployment. It focuses on model-level and algorithm-level optimization methods, including post-training quantization, weight-only quantization, activation-aware quantization, mixed-precision quantization, pruning, sparsity, knowledge distillation, and hybrid compression. The main objective is to summarize current progress, compare important techniques, and identify research gaps related to accuracy preservation, hardware-aware evaluation, latency benchmarking, and practical deployment on real-world platforms. Building on these gaps, the survey further proposes a compression framework that coordinates pruning and quantization in a single adaptive pipeline.

The remainder of this survey is organized as follows. Section~\ref{sec:lr-method} describes the survey methodology, including the literature sources, search terms, and inclusion and exclusion criteria used to select the reviewed papers. Section~\ref{sec:lr-review} presents the literature review and comparative analysis, which is divided into five thematic subsections covering quantization-based optimization, pruning and sparsity-based compression, hybrid quantization and pruning, knowledge distillation, and KV-cache compression, with each subsection accompanied by a structured comparison table. Section~\ref{sec:lr-discussion} discusses the comparative findings across these technique families, addressing their prevalence, ease of implementation, latency and model-size benefits, hardware dependencies, and unresolved challenges. Section~\ref{sec:lr-gaps} identifies the main research gaps that remain in the field, and Section~\ref{sec:lr-future} proposes promising future research directions. Section~\ref{sec:lr-proposed} presents a proposed compression framework that builds on the identified gaps and directions, describing its design objectives, pipeline stages, sensitivity-profiling methods, constraint handling, and evaluation strategy.

\section{Survey Methodology}
\label{sec:lr-method}

This survey followed a structured literature review approach to identify and analyze recent research on quantization and compression techniques for low-latency LLMs deployment. The literature search was conducted using major academic databases and search platforms, including IEEE Xplore, Google Scholar, ScienceDirect, and relevant survey papers. These sources were selected because they provide access to peer-reviewed journal articles, conference papers, preprints, and recent surveys related to efficient deep learning and LLM optimization.

The search focused on papers published from 2020 onward to ensure that the reviewed work reflects recent developments in LLM compression and deployment. The main search terms included combinations of keywords such as ``large language model quantization,'' ``LLM compression,'' ``low-latency LLM deployment,'' ``post-training quantization,'' ``weight-only quantization,'' ``activation-aware quantization,'' ``mixed-precision quantization,'' ``LLM pruning,'' and ``efficient LLM inference.''

The inclusion criteria were based on publication year, relevance to the survey topic, technical contribution, and experimental validation. Papers were included if they focused on quantization, pruning, sparsity, model compression, or efficient inference for LLMs or transformer-based models. Priority was given to studies that provided experimental results, compared their methods with existing baselines, reported performance metrics such as accuracy, latency, memory reduction, compression ratio, or inference speed, and appeared in reputable journals, conferences, or widely cited preprint repositories.

Papers were excluded if they were published before 2020, were not directly related to quantization or compression, focused only on general deep learning without relevance to LLM deployment, lacked experimental evaluation, or did not provide enough technical detail for comparison. Duplicate papers, non-technical articles, and works with limited relevance to low-latency or resource-constrained deployment were also removed from consideration.

In total, approximately 50 papers were reviewed. These papers were analyzed to identify common optimization techniques, compare their advantages and limitations, and extract research gaps related to efficient, low-latency LLM deployment.

\section{Literature Review and Comparative Analysis}
\label{sec:lr-review}

This section reviews recent research on quantization and compression techniques for low-latency LLM deployment and it is organized according to the main optimization directions found in the reviewed papers, namely quantization-based optimization, pruning and sparsity, hybrid quantization and pruning, knowledge distillation, and KV-cache compression. Each group is analyzed in terms of the problem addressed, methodology, experimental focus, key results, hardware dependencies if mentioned, and remaining research gaps. To make the comparison systematic, every subsection concludes with a structured comparison table.

\emph{Standard hardware} means normal CPUs and GPUs that run regular dense matrix operations. These systems do not need special support. Methods such as structural pruning, layer removal, and knowledge distillation create smaller dense models, so they can usually run faster on standard hardware. \emph{Sparse kernels} are special computation methods that skip weights with zero values. These are important for unstructured pruning because removing individual weights does not automatically make the model faster. Without sparse-kernel support, unstructured pruning may reduce the number of parameters, but it may not reduce inference latency. This support is still limited on many common hardware platforms. \emph{Low-rank execution} means using compressed weight matrices directly as smaller matrix operations instead of rebuilding the original full matrix. Low-rank methods only reduce computation if the inference system can run these factorized matrices efficiently.

\subsection{Quantization-Based Optimization}

Quantization is one of the most common methods used to reduce the cost of running LLMs. In normal training and inference, model weights and activations are usually stored using high-precision formats such as FP16 or FP32. These formats require a large amount of memory and computation. Quantization reduces this cost by representing values using lower precision, such as INT8, INT4, or even lower bit-widths. This can reduce memory usage and improve inference speed, especially when the hardware supports low-bit computation.

LLM.int8() showed how large transformer models can be executed using 8-bit matrix multiplication while maintaining good accuracy \cite{dettmers2022llmint8}. The main problem addressed by this method is that simple 8-bit quantization can damage model quality because LLMs contain sensitive outlier values. To solve this, the method keeps important outlier features in higher precision while quantizing the rest of the computation. This makes the model smaller and faster while reducing accuracy loss, cutting inference memory by roughly half and supporting models up to 175B parameters. However, the method mainly focuses on INT8 quantization and does not fully solve the challenges of more aggressive quantization such as INT4 or 2-bit precision.

ZeroQuant also focuses on post-training quantization for large transformer models \cite{yao2022zeroquant}. Post-training quantization is useful because it does not require retraining the full model, which can be expensive. ZeroQuant makes large models smaller by quantizing their values and training the compressed model to imitate the larger model layer by layer, helping it keep good accuracy. Beyond the algorithm, it provides an optimized inference backend and reports speedups of up to 5.19$\times$ over baseline execution. This method is useful for practical deployment because it reduces the cost of optimization. However, its performance depends on calibration data and may vary depending on the model and hardware used.

GPTQ focuses mainly on weight-only quantization, where the model weights are compressed to lower precision while activations are kept in higher precision \cite{frantar2023gptq}. GPTQ uses approximate second-order information to decide how to quantize weights with less accuracy loss. The method showed that large models can be compressed to 3-bit and 4-bit precision with acceptable degradation, quantizing a 175B-parameter model in about four GPU hours and achieving roughly 3.25$\times$ inference speedup on an A100 GPU. However, because GPTQ mainly quantizes weights, it may not fully reduce activation memory or all runtime costs during inference.

SmoothQuant addresses another important problem in LLM quantization: activation outliers \cite{xiao2023smoothquant}. Activations in LLMs can contain very large values, which makes activation quantization difficult. SmoothQuant solves this by shifting some of the quantization difficulty from activations to weights through an equivalent smoothing transformation, enabling training-free W8A8 quantization with close to FP16 accuracy and around 2$\times$ memory reduction. The method is important because activation quantization can help reduce inference cost further. However, it still depends on calibration data and optimized hardware support.

AWQ, or Activation-aware Weight Quantization, improves low-bit weight quantization by identifying which weights are most important for model performance \cite{lin2024awq}. Instead of treating all weights equally, AWQ protects important weights using activation information and quantizes the rest, enabling INT3 and INT4 weight-only deployment without backpropagation or heavy retraining. AWQ is practical for 4-bit LLM deployment and is widely relevant to low-latency inference. However, it mainly focuses on weight compression and does not fully solve activation or KV-cache memory bottlenecks.

A second group of methods targets near-lossless and learnable post-training quantization. SpQR isolates outlier weights and stores them in higher precision while quantizing the remaining dense weights, achieving relative perplexity loss below one percent and allowing a 33B-parameter model to run on a single 24GB GPU \cite{dettmers2024spqr}. SqueezeLLM attacks the memory-bandwidth bottleneck of generative decoding using dense-and-sparse, sensitivity-aware non-uniform quantization, reaching 3-bit precision while improving quality over earlier methods \cite{kim2024squeezellm}. OmniQuant replaces hand-chosen quantization parameters with learnable weight clipping and learnable equivalent transformations, supporting a wide range of settings from W6A6 down to W2A16 \cite{shao2024omniquant}. These works show that careful outlier handling and calibration can preserve accuracy at low bit-widths, although their runtime gains still depend on optimized kernels.

A third group addresses activation quantization through reordering and rotation. RPTQ reorders activation channels and quantizes them in clusters, pushing activations toward 3-bit precision \cite{yuan2023rptq}. QuaRot applies Hadamard-style rotations that preserve the model function while removing outliers, enabling end-to-end 4-bit inference for weights, activations, and KV cache with about 99 percent zero-shot performance retention on LLaMA2-70B \cite{ashkboos2024quarot}. SpinQuant extends this idea by learning the rotation matrices rather than relying only on random rotations, further reducing outlier-driven error \cite{liu2025spinquant}. These methods are important because activation quantization is required for real latency and memory benefits beyond weight-only compression.

A fourth group pursues extreme and codebook-based compression. QuIP showed that quantization improves when weights and Hessian directions are made incoherent before rounding, providing the first viable 2-bit post-training quantization route with theoretical support \cite{chee2023quip}. QuIP\# builds on this with randomized Hadamard transforms and lattice codebooks for compression at or below 4 bits \cite{tseng2024quipsharp}, while AQLM adapts additive multi-codebook quantization to reach 2 to 3 bits per parameter \cite{egiazarian2024aqlm}. QTIP replaces vector codebooks with trellis-coded quantization to obtain a high effective quantization dimension without exponentially large codebooks \cite{tseng2024qtip}. Pushing precision further, BiLLM and PB-LLM binarize most weights while protecting salient ones \cite{huang2024billm,shang2023pbllm}, and PT-BitNet scales 1.58-bit ternary models through post-training quantization, reporting memory compression between 5.72$\times$ and 6.28$\times$\cite{survey2025ptbitnet}. These methods maximize model-size reduction but carry higher accuracy risk and depend strongly on specialized kernels.

Finally, several works connect quantization to training and deployment. EfficientQAT makes quantization-aware training feasible for large models, producing a 2-bit LLaMA-2-70B model with less than a three-point accuracy drop \cite{chen2025efficientqat}, and PV-Tuning provides a representation-agnostic fine-tuning framework for recovering accuracy after extreme 1 to 2-bit compression \cite{malinovskii2024pvtuning}. ABQ-LLM closes the gap between low-bit algorithms and actual GPU acceleration by supporting arbitrary bit-width combinations with distribution correction \cite{zeng2025abqllm}, and FRM-PTQ preserves internal feature relationships to improve low-bit post-training quantization quality for both weights and activations \cite{survey2026frmptq}. Table~\ref{tab:lr-quant} summarizes and compares these quantization techniques.

\begin{table}[htbp]
\centering
\scriptsize
\caption{Comparison of quantization-based optimization techniques.}
\label{tab:lr-quant}
\renewcommand{\arraystretch}{1.2}
\begin{tabular}{p{\dimexpr 0.186\linewidth-2\tabcolsep\relax}p{\dimexpr 0.307\linewidth-2\tabcolsep\relax}p{\dimexpr 0.307\linewidth-2\tabcolsep\relax}p{\dimexpr 0.200\linewidth-2\tabcolsep\relax}}
\hline
\textbf{Technique [Ref]} & \textbf{Approach / Contribution} & \textbf{Key Reported Results} & \textbf{Hardware \& Complexity} \\
\hline
LLM.int8() \cite{dettmers2022llmint8} & Mixed-precision INT8; keeps outliers in FP16 & $\sim$50\% memory cut; supports up to 175B & Moderate HW; low \\
ZeroQuant \cite{yao2022zeroquant} & PTQ with layer-wise distillation and backend & Up to 5.19$\times$ speedup; $\sim$3$\times$ memory cut & Needs backend; moderate \\
GPTQ \cite{frantar2023gptq} & One-shot weight-only PTQ (second-order) & 3--4-bit; $\sim$3.25$\times$ speedup on A100 & Low-bit kernels; low \\
SmoothQuant \cite{xiao2023smoothquant} & Activation-to-weight smoothing & W8A8; $\sim$2$\times$ memory; 1.56$\times$ speedup & INT8 kernels; low \\
AWQ \cite{lin2024awq} & Activation-aware salient-weight scaling & INT3/INT4 weight-only; accuracy preserved & Low-bit kernels; low \\
SpQR \cite{dettmers2024spqr} & Sparse-quantized outlier isolation & 3--4-bit; $>$4$\times$ cut; $<$1\% perplexity loss & Sparse kernels; moderate \\
OmniQuant \cite{shao2024omniquant} & Learnable clipping and transform for PTQ & W6A6 to W2A16; 70B calibrated in $\leq$16 h & Kernel-dependent; moderate \\
SqueezeLLM \cite{kim2024squeezellm} & Dense-and-sparse non-uniform quantization & Up to 3-bit; up to 2.3$\times$ speedup & Sparse kernels; moderate \\
RPTQ \cite{yuan2023rptq} & Reorder-based activation quantization & 3-bit activations; $\sim$80\% memory cut & HW-dependent; moderate \\
QuaRot \cite{ashkboos2024quarot} & Hadamard rotation removes outliers & W4A4; $\sim$99\% zero-shot retention & Rotated kernels; mod.--high \\
SpinQuant \cite{liu2025spinquant} & Learned rotations for weights, act., KV & 4-bit; $\sim$2.9-point zero-shot gap & Rotated kernels; high \\
QuIP \cite{chee2023quip} & Incoherence processing for 2-bit PTQ & First viable 2-bit weight quantization & Low-bit kernels; moderate \\
QuIP\# \cite{tseng2024quipsharp} & Hadamard incoherence and lattice codebooks & $\leq$4-bit; strong extreme compression & Codebook kernels; high \\
AQLM \cite{egiazarian2024aqlm} & Additive multi-codebook quantization & 2--3 bits per parameter; Pareto-strong & Impl.-dependent; high \\
QTIP \cite{tseng2024qtip} & Trellis-coded quantization and incoherence & State-of-the-art sub-4-bit quality/speed & Trellis kernels; high \\
BiLLM \cite{huang2024billm} & 1-bit PTQ with salient-weight protection & $\sim$1.08-bit; 7B binarized in $\sim$0.5 h & Binary kernels; high \\
PB-LLM \cite{shang2023pbllm} & Partial binarization, salient weights kept & Most weights 1-bit; beats naive binarization & Mixed kernels; high \\
PT-BitNet \cite{survey2025ptbitnet} & Post-training 1.58-bit ternary scaling & 5.7--6.3$\times$ memory cut; latency halved & Ternary kernels; high \\
EfficientQAT \cite{chen2025efficientqat} & Efficient quantization-aware training & 2-bit 70B in 41 h; $<$3-point accuracy drop & Post-QAT kernels; high \\
PV-Tuning \cite{malinovskii2024pvtuning} & Post-compression fine-tuning (1--2-bit) & Pareto-optimal 2-bit LLaMA-2 results & Pairs with codebooks; high \\
ABQ-LLM \cite{zeng2025abqllm} & Arbitrary-bit PTQ with GPU framework & W2A8; 1.6$\times$ accel., 2.7$\times$ memory cut & GPU framework; mod.--high \\
FRM-PTQ \cite{survey2026frmptq} & Feature-relationship matching for low-bit PTQ & W4A4/W3A3; $\sim$2$\times$ throughput gain & Low-bit kernels; moderate \\
\hline
\end{tabular}
\end{table}

\subsection{Pruning and Sparsity-Based Compression}

Pruning aims to remove less important weights, neurons, channels, layers, or structures from a model. Unlike quantization, which changes numerical precision, pruning reduces the number of active parameters or computations. This can theoretically reduce memory usage and inference cost. However, practical latency improvement depends strongly on whether the pruning pattern is supported by hardware and inference kernels.

SparseGPT demonstrated that large language models can be pruned in a one-shot manner without full retraining \cite{frantar2023sparsegpt}. The method uses approximate second-order information to remove weights while preserving model behavior, reaching at least 50 percent sparsity and processing OPT-175B and BLOOM-176B in under 4.5 hours. Its main contribution is showing that pruning can be applied to very large models efficiently. However, SparseGPT mainly focuses on unstructured and semi-structured sparsity, which may reduce parameter count but does not always lead to real inference speedup on general-purpose hardware.

Wanda proposed a simple pruning approach based on the product of weight magnitude and input activation magnitude \cite{sun2024wanda}. It showed that effective pruning can be achieved without complex retraining or heavy optimization, which makes it attractive for practical compression pipelines. However, similar to other unstructured pruning methods, the actual latency benefit depends on sparse kernel support and deployment framework compatibility.

LLM-Pruner focused on structural pruning of large language models \cite{ma2023llmpruner}. Instead of removing individual weights, it removes larger coupled structures such as channels or attention heads, then recovers accuracy with lightweight LoRA tuning in roughly three hours. This type of pruning is more likely to produce practical speedup because it changes the model structure in a way that dense hardware can exploit. However, structural pruning may cause larger accuracy degradation compared with unstructured pruning, especially when pruning ratios are high.

SliceGPT compresses LLMs by deleting rows and columns from weight matrices using transformations that preserve computation, removing up to 25 percent of parameters from LLaMA2-70B and OPT-66B while keeping about 99 percent of zero-shot performance \cite{ashkboos2024slicegpt}. ShortGPT investigates layer redundancy and removes less important layers using a Block Influence score, producing smaller dense models that are easier to accelerate than sparse masks \cite{men2024shortgpt}. LaCo follows a related direction by collapsing later layers into earlier ones to reduce depth while keeping average task performance above 80 percent at moderate pruning ratios \cite{yang2024laco}. These works show that LLMs contain substantial structural redundancy, but the best pruning strategy remains model-dependent, and a recurring limitation in pruning literature is that parameter reduction is often reported more clearly than actual latency improvement.

Several recent methods explore alternative pruning criteria and lower-cost pruning workflows. One-shot sensitivity-aware mixed sparsity pruning allocates sparsity adaptively using Hessian-based sensitivity, outperforming SparseGPT on several 13B models at equal sparsity \cite{shao2024sensitivity}. LoRAPrune couples structured pruning with low-rank parameter-efficient fine-tuning, guiding pruning with LoRA gradients and reporting about 52.6 percent memory reduction at 50 percent compression \cite{zhang2024loraprune}. SlimGPT applies an Optimal Brain Surgeon formulation with batched greedy selection and non-uniform pruning ratios for fast structured pruning \cite{ling2024slimgpt}, while bypass-backpropagation pruning and forward-pass-only pruning reduce optimization cost by avoiding full backpropagation through the model \cite{kolawole2026bonsai,gao2025bypass}. OOPS combines outlier awareness with quadratic-programming-based structured pruning and reports lower perplexity than several baselines at 20 to 40 percent pruning rates \cite{survey2026oops}. As a related model-size reduction route, TransAct prunes modules toward low-rank internal structures \cite{shen2024transact}, and SVD-LLM applies truncation-aware singular value decomposition to obtain low-rank compressed weights without sparse masks \cite{wang2025svdllm}. A key open challenge across these works is the gap between theoretical compression and real deployment performance. Table~\ref{tab:lr-prune} compares the main pruning and sparsity techniques.

\begin{table}[htbp]
\centering
\scriptsize
\caption{Comparison of pruning and sparsity-based compression techniques.}
\label{tab:lr-prune}
\renewcommand{\arraystretch}{1.2}
\begin{tabular}{p{\dimexpr 0.186\linewidth-2\tabcolsep\relax}p{\dimexpr 0.307\linewidth-2\tabcolsep\relax}p{\dimexpr 0.307\linewidth-2\tabcolsep\relax}p{\dimexpr 0.200\linewidth-2\tabcolsep\relax}}
\hline
\textbf{Technique [Ref]} & \textbf{Approach / Contribution} & \textbf{Key Reported Results} & \textbf{Hardware \& Complexity} \\
\hline
SparseGPT \cite{frantar2023sparsegpt} & One-shot unstructured/semi-structured pruning & $\geq$50\% sparsity; 175B pruned in $<$4.5 h & Sparse kernels; moderate \\
Wanda \cite{sun2024wanda} & Activation-aware magnitude pruning & 50\% sparsity; no retraining & Sparse kernels; very low \\
LLM-Pruner \cite{ma2023llmpruner} & Structural pruning with LoRA recovery & Recovery in $\sim$3 h; dense smaller model & Standard HW; moderate \\
SliceGPT \cite{ashkboos2024slicegpt} & Deletes rows/columns via PCA-style transform & Up to 25\% params cut; $\sim$99\% retention & Standard HW; moderate \\
ShortGPT \cite{men2024shortgpt} & Layer pruning via Block Influence score & Removes redundant layers; dense output & Standard HW; low \\
LaCo \cite{yang2024laco} & Layer collapse merging later layers & $>$80\% task perf. at 25--30\% pruning & Standard HW; low \\
Sensitivity Pruning \cite{shao2024sensitivity} & Hessian-aware mixed sparsity allocation & $\geq$50\% sparsity; beats SparseGPT on 13B & Sparse kernels; moderate \\
LoRAPrune \cite{zhang2024loraprune} & LoRA-guided structured pruning & $\sim$52.6\% memory cut at 50\% compression & Standard HW; moderate \\
SlimGPT \cite{ling2024slimgpt} & Optimal-Brain-Surgeon structured pruning & Local pruning in $\sim$1 h; SOTA on LLaMA & Standard HW; moderate \\
Bypass-Backprop \cite{gao2025bypass} & Policy-gradient structured pruning & Beats baselines at equal pruning ratios & Standard HW; moderate \\
Everybody Prune Now \cite{kolawole2026bonsai} & Forward-pass-only structured pruning & Competitive quality; no backpropagation & Standard HW; low \\
TransAct \cite{shen2024transact} & Intra-module low-rank structured pruning & Better compression-quality trade-off & Standard HW; high \\
OOPS \cite{survey2026oops} & Outlier-aware QP-based structured pruning & Beats baselines at 20--40\% pruning rates & Standard HW; high \\
SVD-LLM \cite{wang2025svdllm} & Truncation-aware SVD low-rank compression & Strong results at high compression ratios & Low-rank exec.; moderate \\
\hline
\end{tabular}
\end{table}

\subsection{Hybrid Quantization and Pruning}

Hybrid compression combines multiple optimization techniques, usually quantization and pruning, to achieve better efficiency than either method alone. The motivation is that quantization is effective for reducing memory usage, while pruning can reduce model complexity and computation. Combining both may produce better trade-offs between model size, latency, memory usage, and accuracy.

DayPQ proposes dynamic layer-wise pruning and quantization for LLM inference acceleration \cite{xu2026daypq}. The method addresses the problem that different layers have different sensitivity to compression. Instead of applying the same compression level to all layers, it dynamically selects pruning and quantization strategies across layers, and through an algorithm-hardware co-design with bias-shifting FP8 quantization it reports up to 3.44$\times$ speedup and more than 3$\times$ energy-efficiency improvement at less than one percent accuracy loss. However, the method also highlights a broader challenge: adaptive compression strategies require careful selection rules and may introduce extra design complexity.

Sub-4-bit power-of-two-based mixed-precision quantization focuses on efficient LLM compression and acceleration using very low precision \cite{cho2025sub4bit}. The work addresses the need for stronger compression beyond standard 4-bit quantization. By assigning bit-widths according to weight importance and using a logarithmic quantization scheme co-designed with an FPGA accelerator, it reaches an average of about 3.6 bits per weight and reports up to 11.8$\times$ model-size reduction. Its key strength is aggressive compression, but its limitation is that performance depends heavily on hardware support for low-bit arithmetic.

OFQ-LLM addresses outlier-related challenges in low-bit quantization \cite{wang2025ofqllm}. Since outliers can cause significant accuracy degradation during quantization, the method uses outlier-flexing encoding that reuses unused quantization space to represent outliers efficiently, handling both weights and activations and reporting prefill speedups of up to 3.83$\times$. This work is important because outlier handling is a recurring issue in LLM quantization. However, like many quantization methods, its performance must be evaluated across different models and deployment frameworks to confirm generality.

A further related direction combines pruning with knowledge distillation. The Minitron-style approach derives smaller dense models from a larger pretrained model through structured pruning followed by distillation-based recovery \cite{muralidharan2024minitron}, showing that hybridizing structural compression with a recovery stage can produce deployable compact models at a fraction of the original training cost. Hybrid compression remains promising, but the literature does not yet provide clear rules for choosing when to use quantization-only, pruning-only, or combined pipelines. This creates an important research gap for future work: testing these pipelines under the same model, dataset, hardware, and inference backend to determine which approach gives the best practical deployment trade-off. Table~\ref{tab:lr-hybrid} compares the main hybrid techniques.

\begin{table}[htbp]
\centering
\scriptsize
\caption{Comparison of hybrid quantization and pruning techniques.}
\label{tab:lr-hybrid}
\renewcommand{\arraystretch}{1.2}
\begin{tabular}{p{\dimexpr 0.186\linewidth-2\tabcolsep\relax}p{\dimexpr 0.307\linewidth-2\tabcolsep\relax}p{\dimexpr 0.307\linewidth-2\tabcolsep\relax}p{\dimexpr 0.200\linewidth-2\tabcolsep\relax}}
\hline
\textbf{Technique [Ref]} & \textbf{Approach / Contribution} & \textbf{Key Reported Results} & \textbf{Hardware \& Complexity} \\
\hline
DayPQ \cite{xu2026daypq} & Dynamic layer-wise pruning and FP8 quantization & Up to 3.44$\times$ speedup; $<$1\% accuracy loss & Custom accelerator; high \\
Sub-4-bit Mixed-Precision \cite{cho2025sub4bit} & Mixed-precision power-of-two quant. with FPGA & $\sim$3.6 bits avg.; up to 11.8$\times$ size cut & FPGA; moderate--high \\
OFQ-LLM \cite{wang2025ofqllm} & Outlier-flexing encoding for weights and act. & Prefill up to 3.83$\times$; decode up to 2.01$\times$ & Custom accelerator; high \\
Minitron Prune+Distill \cite{muralidharan2024minitron} & Structured pruning plus distillation recovery & 15B to 8B/4B using $<$3\% of original data & Standard HW; moderate \\
\hline
\end{tabular}
\end{table}

\subsection{Knowledge Distillation and Model Size Reduction}

Knowledge distillation is another compression technique in which a smaller student model learns from a larger teacher model. Although it is not always considered a direct low-latency inference optimization method, it is highly relevant because smaller models usually require less memory and computation.

DistilBERT showed that a smaller language model can retain much of the performance of BERT while being faster and lighter, achieving about 60 percent faster inference and roughly 97 percent of the original performance at about 40 percent smaller size \cite{sanh2020distilbert}. TinyBERT similarly demonstrated that transformer-based models can be compressed through a two-stage distillation process, with TinyBERT4 reaching over 96.8 percent of BERT-base performance while being about 7.5 times smaller and 9.4 times faster \cite{jiao2020tinybert}. More recent works extend distillation ideas to larger generative language models. Distilling Step-by-Step transfers not only labels but also teacher rationales, allowing a 770M-parameter student to outperform a much larger prompted model on selected tasks with less training data \cite{hsieh2023distilling}, and MiniLLM replaces the usual forward KL objective with reverse KL, which is better suited to open-ended generation \cite{gu2024minillm}. The Minitron-style pipeline, discussed above as a hybrid method, further combines structured pruning with distillation to produce compact model families \cite{muralidharan2024minitron}. These methods are useful for reducing model size, but they often require additional training data, teacher model access, and training cost. In addition, distilled models may not preserve all capabilities of the original model, especially for reasoning and long-context tasks. Table~\ref{tab:lr-kd} summarizes the main distillation techniques.

\begin{table}[htbp]
\centering
\scriptsize
\caption{Comparison of knowledge distillation techniques for model-size reduction.}
\label{tab:lr-kd}
\renewcommand{\arraystretch}{1.2}
\begin{tabular}{p{\dimexpr 0.186\linewidth-2\tabcolsep\relax}p{\dimexpr 0.307\linewidth-2\tabcolsep\relax}p{\dimexpr 0.307\linewidth-2\tabcolsep\relax}p{\dimexpr 0.200\linewidth-2\tabcolsep\relax}}
\hline
\textbf{Technique [Ref]} & \textbf{Approach / Contribution} & \textbf{Key Reported Results} & \textbf{Hardware \& Complexity} \\
\hline
DistilBERT \cite{sanh2020distilbert} & Knowledge distillation of BERT into a student & $\sim$40\% smaller, $\sim$60\% faster, $\sim$97\% perf. & Standard HW; low \\
TinyBERT \cite{jiao2020tinybert} & Two-stage pretraining and task distillation & $\sim$7.5$\times$ smaller; $\sim$9.4$\times$ faster inference & Standard HW; moderate \\
Distilling Step-by-Step \cite{hsieh2023distilling} & Rationale-based distillation, multi-task & 770M student beats a 540B prompted model & Standard HW; moderate \\
MiniLLM \cite{gu2024minillm} & Reverse-KL distillation for generative LLMs & Students 120M--13B; better preference scores & Standard HW; mod.--high \\
Minitron \cite{muralidharan2024minitron} & Pruning combined with distillation recovery & Up to 40$\times$ fewer training tokens per model & Standard HW; moderate \\
\hline
\end{tabular}
\end{table}

\subsection{KV-Cache Compression and Memory-Aware Inference}

A separate but increasingly important line of work targets the key-value (KV) cache rather than model weights. During autoregressive generation, LLMs store key and value tensors for every previous token, so the KV cache grows linearly with sequence length and quickly becomes the dominant memory and bandwidth cost in long-context inference. Because this cost appears at runtime, KV-cache compression directly affects latency and throughput in ways that weight compression alone cannot.

H2O addresses the KV-cache memory bottleneck through a heavy-hitter eviction policy that keeps the tokens contributing most to attention while discarding less useful cache entries \cite{zhang2023h2o}. Retaining only about 20 percent of tokens as heavy hitters, it reports throughput improvements of up to 29 times over baseline inference systems and latency reductions of up to 1.9 times, all without retraining. KVQuant instead applies low-bit quantization to the KV cache, using distribution-aware per-channel key quantization, pre-RoPE quantization, and dense-and-sparse outlier handling to reach sub-4-bit precision with less than 0.1 perplexity degradation, enabling context lengths up to one million tokens on a single A100-80GB GPU \cite{hooper2024kvquant}. SnapKV reduces KV-cache size in a fine-tuning-free manner by selecting important positions per attention head from attention patterns observed near the end of the prompt, reporting a 3.6 times generation speedup and 8.2 times memory-efficiency improvement at 16K tokens \cite{li2024snapkv}. InfiniGen takes a system-level view: in offloading-based serving it predicts and prefetches only the KV entries that will be important, improving a modern offloading system by up to 3 times over prior KV-cache management methods \cite{lee2024infinigen}.

Together, these works show that long-context, low-latency deployment depends on memory movement as much as on model compression. However, eviction- and selection-based methods risk accuracy loss when attention patterns shift unexpectedly, quantization-based methods depend on custom kernels, and system-level methods assume a particular offloading architecture. Table~\ref{tab:lr-kv} compares the main KV-cache techniques.

\begin{table}[htbp]
\centering
\scriptsize
\caption{Comparison of KV-cache compression and memory-aware inference techniques.}
\label{tab:lr-kv}
\renewcommand{\arraystretch}{1.2}
\begin{tabular}{p{\dimexpr 0.186\linewidth-2\tabcolsep\relax}p{\dimexpr 0.307\linewidth-2\tabcolsep\relax}p{\dimexpr 0.307\linewidth-2\tabcolsep\relax}p{\dimexpr 0.200\linewidth-2\tabcolsep\relax}}
\hline
\textbf{Technique [Ref]} & \textbf{Approach / Contribution} & \textbf{Key Reported Results} & \textbf{Hardware \& Complexity} \\
\hline
H2O \cite{zhang2023h2o} & Heavy-hitter KV-cache eviction policy & Up to 29$\times$ throughput; up to 1.9$\times$ latency & Existing systems; low \\
KVQuant \cite{hooper2024kvquant} & Sub-4-bit KV-cache quantization & $<$0.1 perplexity loss at 3-bit; 1M context & Custom CUDA; high \\
SnapKV \cite{li2024snapkv} & Per-head selection of important KV positions & 3.6$\times$ speedup, 8.2$\times$ memory at 16K tokens & Attention kernels; moderate \\
InfiniGen \cite{lee2024infinigen} & Speculative KV prefetch for offloading & Up to 3.0$\times$ over prior KV-cache methods & Offloading arch.; mod.--high \\
\hline
\end{tabular}
\end{table}

\section{Discussion}
\label{sec:lr-discussion}

The comparative analysis in the previous section reveals several consistent patterns across the reviewed literature, which are discussed here along six dimensions: prevalence, ease of implementation, low-latency performance, model-size reduction, hardware dependency, and unresolved challenges. Table~\ref{tab:lr-families} provides a consolidated qualitative comparison of the five technique families.

\begin{table}[htbp]
\centering
\scriptsize
\caption{Consolidated qualitative comparison of the five compression technique families.}
\label{tab:lr-families}
\renewcommand{\arraystretch}{1.2}
\begin{tabular}{p{\dimexpr 0.177\linewidth-2\tabcolsep\relax}p{\dimexpr 0.170\linewidth-2\tabcolsep\relax}p{\dimexpr 0.227\linewidth-2\tabcolsep\relax}p{\dimexpr 0.227\linewidth-2\tabcolsep\relax}p{\dimexpr 0.199\linewidth-2\tabcolsep\relax}}
\hline
\textbf{Technique Family} & \textbf{Prevalence in Literature} & \textbf{Implementation Ease} & \textbf{Latency / Size Benefit} & \textbf{Hardware Dependency} \\
\hline
Quantization (PTQ) & Very high; dominant category & High for INT8 and INT4 PTQ; lower for QAT and codebook methods & Strong size reduction; latency gains kernel dependent & Moderate; low-bit kernels needed \\
Pruning and sparsity & High & High for magnitude pruning; moderate for structured methods & Moderate size reduction; latency often unverified & High for unstructured; low for structured \\
Hybrid quantization and pruning & Growing rapidly & Low; co-design and adaptive rules add complexity & Among the highest reported speedups and size reduction & High; often needs custom accelerators \\
Knowledge distillation & Moderate & Moderate; needs teacher access and training & Strong via smaller dense models; latency hardware friendly & Low \\
KV-cache compression & Growing rapidly & High; several methods are fine-tuning-free & Strong for long-context latency and memory & Low to moderate \\
\hline
\end{tabular}
\end{table}

Regarding prevalence, quantization is clearly the most studied direction. From what has been gathered, there is substantially more quantization papers than any other category. Structural and unstructured pruning forms the second most active area, while hybrid compression, knowledge distillation, and KV-cache compression are smaller but rapidly growing categories. The recent emergence of dynamic and co-designed hybrid methods indicates a shift toward adaptive, deployment-aware optimization.

In terms of ease of implementation, post-training quantization methods are generally the most practical because they avoid full retraining and are already integrated into common serving libraries. Methods such as GPTQ, SmoothQuant, and AWQ require only calibration data and are relatively simple to apply \cite{frantar2023gptq,xiao2023smoothquant,lin2024awq}. Magnitude-based pruning methods such as Wanda are similarly lightweight \cite{sun2024wanda}, and several KV-cache methods are fine-tuning-free \cite{zhang2023h2o,li2024snapkv}. In contrast, quantization-aware training, learned-rotation methods, codebook quantization, and algorithm-hardware co-designs are considerably more complex, requiring training infrastructure, specialized kernels, or custom accelerators \cite{xu2026daypq,liu2025spinquant,tseng2024qtip,chen2025efficientqat}.

For low-latency performance, the strongest reported results come from methods that target the runtime bottleneck directly rather than only reducing stored size. Hybrid co-designed methods report some of the largest speedups, with DayPQ reporting up to 3.44 times speedup and OFQ-LLM reporting prefill speedups up to 3.83 times \cite{xu2026daypq,wang2025ofqllm}. KV-cache methods are especially effective for long-context workloads, where H2O reports throughput gains up to 29 times \cite{zhang2023h2o}. Among pure quantization methods, ZeroQuant and GPTQ report multi-fold speedups when paired with optimized backends \cite{yao2022zeroquant,frantar2023gptq}, but a recurring observation is that latency improvements depend heavily on kernel and hardware support.

For model-size reduction, extreme quantization achieves the most aggressive results. Sub-4-bit mixed-precision quantization reports up to 11.8 times model-size reduction \cite{cho2025sub4bit}, ternary PT-BitNet reports memory compression above 6 times \cite{survey2025ptbitnet}, and additive and trellis codebook methods reach 2 to 3 bits per parameter \cite{egiazarian2024aqlm,tseng2024qtip}. Pruning and low-rank methods typically achieve more moderate reductions, often in the 20 to 50 percent range \cite{ashkboos2024slicegpt,zhang2024loraprune,wang2025svdllm}, while distillation reduces size by producing entirely smaller architectures \cite{jiao2020tinybert,muralidharan2024minitron}.

Hardware dependency is one of the clearest dividing lines in the literature. Algorithm-hardware co-designs explicitly depend on FPGA or custom accelerator support \cite{xu2026daypq,cho2025sub4bit,wang2025ofqllm}, rotation-based and codebook quantization require specialized kernels \cite{ashkboos2024quarot,tseng2024qtip}, and unstructured pruning needs sparse kernels to convert parameter reduction into actual speedup \cite{frantar2023sparsegpt,sun2024wanda}. By contrast, structural pruning and distillation produce smaller dense models that run efficiently on standard hardware \cite{ashkboos2024slicegpt,sanh2020distilbert,men2024shortgpt}, which is a practical advantage for general-purpose deployment. illation techniques for model-size reduction. ssion and memory-aware inference techniques.

Finally, several challenges remain unresolved. Extreme low-bit quantization still carries accuracy risk and often requires post-compression fine-tuning to recover quality \cite{chen2025efficientqat,malinovskii2024pvtuning}. Pruning literature continues to report parameter reduction more clearly than measured latency. Eviction-based KV-cache methods can lose accuracy when attention patterns shift \cite{zhang2023h2o,li2024snapkv}. Across all categories, the field lacks unified evaluation under identical models, datasets, hardware, and inference backends, which makes direct comparison difficult.

\section{Research Gaps}
\label{sec:lr-gaps}

Although the reviewed techniques achieve strong compression and efficiency results, the comparative analysis exposes several research gaps that limit progress toward efficient, low-latency LLM deployment.

\subsection{Lack of a practical comparison between quantization precision levels}

Quantization methods are frequently evaluated using different models, datasets, hardware platforms, and metrics. Because of this inconsistency, it is difficult to determine which precision level provides the best balance between model size, latency, memory usage, and accuracy. A controlled comparison of precision levels such as INT8, INT4, and sub-4-bit settings \cite{dettmers2022llmint8,frantar2023gptq,huang2024billm} under identical experimental conditions would clarify these trade-offs and support more informed precision selection.

\subsection{Lack of clear rules for choosing the best compression method}

Quantization tends to be more effective at reducing memory usage, whereas pruning tends to be more effective at reducing computation, yet there is no clear guidance on when each approach should be preferred. A promising direction is to evaluate quantization-only, pruning-only, and combined quantization-plus-pruning pipelines \cite{xu2026daypq} on the same model and to directly compare their latency, model size, memory usage, and accuracy.

\subsection{Pruning reduces parameters but does not always reduce real latency}

Pruning can substantially reduce the number of parameters, but this theoretical compression does not always translate into faster inference, leaving a gap between reported compression and actual low-latency deployment. Future work could directly compare unstructured pruning \cite{frantar2023sparsegpt,sun2024wanda} and structured pruning \cite{ashkboos2024slicegpt,men2024shortgpt} while measuring real inference latency rather than parameter reduction alone.

\subsection{Limited deployment-aware compression across different inference tools}

Although many compression methods are tested on real hardware, particularly GPUs, they are usually evaluated within a narrow setup, such as a single GPU type, one model family, or one inference backend. This makes it difficult to know whether a compressed model that runs efficiently in one framework will achieve comparable speedups in another, since such gains often depend on the availability of optimized low-bit kernels. A useful direction is to evaluate the same compressed model across multiple inference tools, such as Hugging Face Transformers, llama.cpp, vLLM, and TensorRT-LLM, and to compare latency, memory usage, and accuracy across these environments.

\subsection{Limited adaptive compression based on latency and accuracy needs}

Most quantization and pruning methods apply fixed compression settings, such as a fixed bit-width or a fixed pruning ratio \cite{cho2025sub4bit}. Real deployment conditions, however, are not fixed: some applications require very low latency, others prioritize accuracy, and target devices differ in available memory. Current methods rarely adapt to these varying latency targets, memory limits, or accuracy requirements. A valuable direction is to develop and evaluate adaptive compression strategies that select the bit-width, pruning ratio, or overall compression level according to the specific deployment requirement. on of the five compression technique families.

\section{Future Directions}
\label{sec:lr-future}

Building on the gaps identified above, several research directions appear promising. Consistent with the model-level and algorithm-level focus of this survey, the directions below concern improvements to compression methods and inference algorithms themselves; hardware-level co-design and system-level deployment infrastructure are left outside the present scope.

The first direction is to design quantization and pruning together instead of applying them one after another. In many current methods, the model is first pruned and then quantized, or quantized and then pruned. This can cause problems because the first step does not consider the errors caused by the second step. A better approach would estimate how each layer reacts to both pruning and low-bit quantization at the same time. This would help choose which layers should be pruned, which layers should be quantized, and which layers need to be protected to avoid accuracy loss. This creates a more complete compression method instead of two separate compression steps.

A second direction is stronger low-bit quantization for extreme compression. While 8-bit and 4-bit quantization are now close to lossless, 1-bit, ternary, and 2-bit methods still degrade accuracy and depend on additional recovery steps \cite{survey2025ptbitnet,huang2024billm,shang2023pbllm}. Advances in outlier handling, codebook-and rotation-based quantization, and post-compression fine-tuning \cite{chen2025efficientqat,malinovskii2024pvtuning} could make extreme low-bit models reliable across diverse model families, instruction-tuned models, and reasoning tasks.

A third direction, this time at the algorithm level, is adaptive, input-dependent compression. Where most methods use a single fixed configuration, adaptive schemes could allocate bit-width and pruning ratio per layer, per input, or per accuracy target. Building on dynamic layer-wise and arbitrary-bit methods \cite{xu2026daypq,zeng2025abqllm} and on frameworks that already support several bit-width settings \cite{shao2024omniquant}, future work could train controllers that select a compression level directly from the model's own sensitivity signals.

Activations and the KV cache deserve more algorithm-level attention as well. Weight compression is comparatively mature, but many methods keep activations in higher precision, and the KV cache dominates memory and bandwidth cost in long-context inference. Progress on activation quantization \cite{xiao2023smoothquant,yuan2023rptq} and KV-cache compression \cite{zhang2023h2o,hooper2024kvquant,li2024snapkv} is therefore essential for reducing end-to-end latency rather than storage alone.

A further direction is combining model-level and algorithm-level compression, which are usually studied in isolation. Pairing a low-bit quantized model with an algorithm-level KV-cache method \cite{zhang2023h2o,hooper2024kvquant,ashkboos2024quarot} in a single pipeline, and characterizing how the two interact, could yield larger and more reliable latency gains than either level alone.

Finally, compression for multilingual and domain-specific models warrants study. Most methods are evaluated on general English benchmarks, so the effect of aggressive compression on multilingual capability, long-context reasoning, and specialized domains remains poorly understood. Combining compression with knowledge distillation and post-compression adaptation \cite{hsieh2023distilling,malinovskii2024pvtuning,muralidharan2024minitron} may help preserve these capabilities in compact, deployable models.

\section{Proposed Methodology}
\label{sec:lr-proposed}

Building on the research gaps and directions identified above, this section presents a proposed compression framework. The framework is a sensitivity-driven, adaptive pipeline that coordinates pruning and quantization per layer and extends compression to the KV cache and activations, with every configuration evaluated under a single measured protocol. It is described in terms of its design objectives, overall structure, stage-by-stage pipeline, the methods used to estimate weight importance, the handling of compute, memory, and energy constraints, and an evaluation strategy.

\subsection{Design Objectives}

\subsubsection{Objectives}

The framework is designed to satisfy the following objectives, each traceable to a survey finding:

\begin{itemize}

  \item \textbf{Joint, not sequential, compression:} Estimate how each layer responds to both pruning and low-bit quantization at the same time, so the choice to prune, quantize, protect, or do both is made per layer instead of as two independent passes whose errors do not account for each other.

  \item \textbf{Adaptive to the deployment requirement:} Decide how aggressively to compress each layer---how few bits to use, how much to prune, and how much to shrink the KV cache---based on how sensitive that layer is and on the target's speed, memory, and accuracy limits, instead of applying one fixed setting to the whole model.

  \item \textbf{Unified model-level and algorithm-level coverage:} Combine weight/activation quantization and pruning (model level) with KV-cache compression (algorithm level) in one pipeline, since the survey shows the KV cache dominates long-context memory and bandwidth and that activations remain comparatively neglected.

  \item \textbf{Deployment-aware, measured evaluation:} Report compression, latency, memory, energy, and accuracy under identical models, datasets, and inference backends---measuring real latency rather than parameter reduction---to close the inconsistent, single-setup evaluation the survey highlights.

  \item \textbf{Accuracy preservation under aggressive compression:} Include an explicit distillation/fine-tuning recovery stage so that extreme low-bit or high-sparsity operating points remain within the accuracy tolerance.

\end{itemize}

\subsubsection{Theoretical Justification of Design Choices}

Each of these choices rests on a specific finding from the survey:

\begin{itemize}

  \item \textbf{Joint over sequential:} Pruning and quantization each introduce error, and the two errors interact. When the operations are applied sequentially, the first commits to its decisions before the second is considered, in effect optimizing for a model whose numerical representation is about to change. Magnitude-based pruning, for instance, retains the largest-magnitude weights and discards the smallest; the retained large values widen the dynamic range that quantization must subsequently represent within a limited bit budget, coarsening the quantization grid and increasing error across all weights in the layer. The pruning step has no mechanism by which to account for the additional error it imposes on quantization. Estimating each layer's sensitivity to both operations jointly eliminates this blind spot: the planner can quantify, on a per-layer basis, the accuracy cost of pruning and of quantization, and allocate compression aggressively where it is inexpensive and conservatively where it is costly. The outcome is a single coordinated decision per layer rather than two mutually uninformed steps.

  \item \textbf{Sensitivity-driven adaptivity:} Within each layer, a small subset of weights and activations (outliers) accounts for a disproportionate share of model behavior. Uniform compression is consequently constrained by these worst-case elements, requiring conservative settings across the entire layer in order to preserve them. Allocating compression per layer according to measured sensitivity enables the framework to protect only the sensitive elements while compressing the remainder aggressively, thereby satisfying the deployment requirement without being limited by the worst case.

  \item \textbf{Distillation in the loop:} Compression and model capability are inherently in tension. Incorporating a lightweight recovery stage---a brief round of knowledge distillation or fine-tuning---restores part of the accuracy lost to compression, expanding the feasible operating region and rendering attainable the extreme low-bit configurations that would otherwise violate the accuracy tolerance.

  \item \textbf{Algorithm-level KV/activation focus:} In long-context inference, the key--value cache and the activations, rather than the weights, dominate runtime memory consumption and bandwidth. Reducing end-to-end latency therefore requires compressing these components as well, and not the weights alone.

  \item \textbf{Standard-hardware preference:} Many of the strongest results reported in the literature attain their speedups only on custom accelerators, which fall outside the scope of this work. Restricting the framework to methods that execute efficiently on commodity CPUs and GPUs---structured pruning, distillation, and INT4/INT8 post-training quantization---preserves deployability on widely available hardware while retaining compatibility with specialized kernels where present.

  \item \textbf{Unified measured evaluation:} For a given model and target device, the trade-off among model size, latency, memory, energy, and accuracy cannot be determined a priori and must be established empirically. Evaluating the same model across multiple inference backends and recording these metrics reframes deployment as a constrained search rather than an estimation, and yields the directly comparable evidence that the field currently lacks.

\end{itemize}

\subsection{Proposed Framework}

The framework is organized as a sensitivity-driven, adaptive pipeline, shown in Fig.~\ref{fig:lr-pipeline}. Inputs enter a joint planning stage (Stage 0) that profiles per-layer sensitivity and produces a compression plan conditioned on the deployment requirement. Four processing stages then apply joint quantization and pruning, distillation-based recovery, KV-cache and activation compression, and deployment-aware evaluation. A feedback path returns control to the planner so that the plan is revised until the deployment requirement is met. The lower band lists the unified metrics on which every configuration is measured.

\begin{figure}[htbp]
\centering
\includegraphics[width=\textwidth]{Figures/lr-fig1-pipeline.png}
\caption{Proposed sensitivity-driven joint compression pipeline. The dashed amber path is the adaptive feedback that re-plans per-layer compression until the deployment requirement is satisfied. Scope is model-level and algorithm-level only.}
\label{fig:lr-pipeline}
\end{figure}

\subsubsection{Key Components and Data Flow}

\begin{itemize}

  \item \textbf{Inputs.} A pretrained LLM; a deployment requirement expressed as a latency target, memory limit, and accuracy tolerance; and calibration/domain data used for post-training quantization calibration and for distillation.

  \item \textbf{Stage 0 --- Joint sensitivity profiling and planning.} Estimates each layer's sensitivity to both pruning and quantization and emits a per-layer plan. This is the only component with a global view and is what makes the pipeline adaptive.

  \item \textbf{Stages 1--4 --- Processing pipeline.} Joint quantization and pruning, distillation-based recovery, KV-cache and activation compression, and deployment-aware evaluation, detailed in the methodology pipeline below.

  \item \textbf{Adaptive feedback.} If evaluation shows the requirement is not met, control returns to Stage 0, which adjusts the per-layer bit-widths, pruning ratios, or KV policy and re-runs the affected stages---turning a one-shot script into a constrained search over the size--latency--memory--accuracy trade-off.

  \item \textbf{Output.} A deployable low-latency model for edge, enterprise, or offline serving, optionally embedded in a retrieval-augmented pipeline.

\end{itemize}

\subsection{Methodology Pipeline}

\subsubsection{Stage 0 --- Joint Sensitivity Profiling and Compression Planning}

The planner characterizes each layer's tolerance to compression using second-order (Hessian-based) and activation-based importance signals---the same signals that underpin the strongest individual methods, e.g. second-order weight quantization \cite{frantar2023gptq}, Hessian-aware mixed-sparsity pruning \cite{shao2024sensitivity}, and activation-aware weight protection \cite{lin2024awq}. Crucially, sensitivity is estimated with respect to \emph{both} pruning and quantization jointly, rather than for one in isolation, so that a layer fragile to one but robust to the other is treated correctly. The output is a per-layer plan---a target bit-width, a pruning ratio, and a ``protect'' flag for the most sensitive layers---chosen to satisfy the deployment requirement. This directly implements the survey's recommendation to estimate how each layer reacts to pruning and low-bit quantization at the same time, instead of running two separate compression steps.

\emph{Techniques that can be used, and how each works (sensitivity signals):}

\begin{itemize}

  \item \textbf{Magnitude scoring} --- ranks weights by absolute value and compresses the smallest first, assuming small weights contribute least; cheapest signal, but blind to how a weight interacts with its inputs.

  \item \textbf{Activation-aware scoring} \cite{lin2024awq,sun2024wanda} --- judges a weight by its magnitude times the size of the input it multiplies. Wanda scores each weight as $|w|\times$ input-feature norm and prunes the lowest per neuron; AWQ scales up the salient channels (those tied to large activations) before rounding so they survive low-bit quantization.

  \item \textbf{First-order / gradient (Taylor) scoring} \cite{ma2023llmpruner} --- estimates the loss increase from removing a weight as roughly $|$weight $\times$ gradient$|$, applied to coupled groups of weights that must be removed together, then recovers with LoRA (LLM-Pruner).

  \item \textbf{Second-order / Hessian scoring} \cite{frantar2023gptq,frantar2023sparsegpt,ling2024slimgpt} --- uses loss curvature: saliency is about $w^2$ over the inverse-Hessian diagonal, and the inverse Hessian also prescribes how to adjust the surviving weights to compensate. GPTQ applies this to rounding, SparseGPT to pruning, SlimGPT to structured pruning.

  \item \textbf{Fisher (squared-gradient) proxy} --- approximates the Hessian with averaged squared gradients, giving curvature-like sensitivity without computing true second derivatives.

  \item \textbf{Representational-change scoring} \cite{yang2024laco,men2024shortgpt} --- scores whole layers by how little they change the hidden state. ShortGPT drops layers whose input and output are nearly identical (low block influence); LaCo merges consecutive near-redundant layers to cut depth.

\end{itemize}

\subsubsection{Stage 1 --- Joint Quantization and Pruning}

Guided by the Stage 0 plan, each layer is compressed by the assigned combination of quantization and pruning in a single coordinated pass. The framework handles the two together rather than one after the other, which is the defining feature distinguishing it from the sequential prune-then-quantize or quantize-then-prune pipelines criticized in the survey.

\textbf{Quantization.} Mixed-precision post-training quantization (PTQ) is the default because it avoids full retraining and is already integrated into common serving libraries; representative operators include one-shot weight-only PTQ \cite{frantar2023gptq}, activation-aware scaling \cite{lin2024awq}, activation-to-weight smoothing \cite{xiao2023smoothquant}, and learnable clipping/transforms \cite{shao2024omniquant}. For extreme low-bit targets the planner may escalate to quantization-aware training \cite{chen2025efficientqat} or post-compression fine-tuning \cite{malinovskii2024pvtuning}.

\begin{itemize}

  \item \textbf{GPTQ} \cite{frantar2023gptq} --- quantizes weights one column at a time and, after each, updates the not-yet-quantized weights via the layer-input inverse Hessian to absorb the rounding error; weight-only 3--4-bit from a small calibration set.

  \item \textbf{AWQ} \cite{lin2024awq} --- identifies the salient weight channels from activation magnitudes and protects them by per-channel scaling before low-bit rounding; weight-only INT3/INT4 with no backpropagation.

  \item \textbf{SpQR} \cite{dettmers2024spqr} --- isolates the few outlier weights with the largest quantization error into a sparse high-precision store and quantizes the dense remainder to 3--4-bit; near-lossless.

  \item \textbf{SmoothQuant} \cite{xiao2023smoothquant} --- divides each activation channel by a scale and multiplies the matching weights by it (identical output), shifting the hard-to-quantize outliers from activations onto weights so both reach INT8 (W8A8).

  \item \textbf{OmniQuant} \cite{shao2024omniquant} --- learns, by gradient descent on each block's reconstruction error, a weight clipping range and an activation-weight smoothing transform; covers W6A6 down to W2A16.

  \item \textbf{QuaRot} \cite{ashkboos2024quarot} --- multiplies weights and activations by Hadamard rotations that cancel across the network, spreading outliers evenly so uniform 4-bit works for weights, activations, and KV cache (W4A4).

  \item \textbf{SpinQuant} \cite{liu2025spinquant} --- uses the same rotation idea but learns the rotation matrices on calibration data instead of using random ones, reducing the residual quantization error.

  \item \textbf{QuIP / QuIP\#} \cite{chee2023quip,tseng2024quipsharp} --- ``incoherence processing'' rotates weights and their Hessian so neither has concentrated large entries, making low-bit rounding provably easier; QuIP\# adds lattice codebooks for vector quantization at or below 4 bits.

  \item \textbf{AQLM} \cite{egiazarian2024aqlm} --- represents each weight group as a sum of vectors drawn from several learned codebooks, jointly optimized to minimize error; reaches 2--3 bits per parameter.

  \item \textbf{QTIP} \cite{tseng2024qtip} --- uses a trellis (stateful sequence) code to obtain a high effective quantization dimension without an exponentially large codebook, combined with incoherence processing.

  \item \textbf{SqueezeLLM} \cite{kim2024squeezellm} --- places non-uniform quantization levels by sensitivity-weighted k-means and splits outlier weights into a sparse high-precision part, targeting the decode memory-bandwidth bottleneck.

  \item \textbf{EfficientQAT} \cite{chen2025efficientqat} --- makes quantization-aware training affordable by first training each block's parameters, then end-to-end training only the quantization scales; 2-bit 70B with a small accuracy drop.

  \item \textbf{PV-Tuning} \cite{malinovskii2024pvtuning} --- recovers accuracy after extreme 1--2-bit compression by going beyond the straight-through estimator, alternately updating the continuous weights and the discrete codeword each maps to.

\end{itemize}

\textbf{Pruning.} Structured pruning is preferred because it yields smaller dense models that accelerate on standard hardware \cite{ma2023llmpruner,ashkboos2024slicegpt,men2024shortgpt}; unstructured or semi-structured pruning \cite{frantar2023sparsegpt,sun2024wanda} is selected only when the target provides sparse-kernel support, since otherwise parameter reduction does not translate into measured speedup. Low-rank factorization \cite{wang2025svdllm} is available as a further dense reduction route.

\begin{itemize}

  \item \textbf{LLM-Pruner} \cite{ma2023llmpruner} --- structural pruning that uses gradient/Taylor importance on coupled groups (heads, channels removed together), then restores accuracy with lightweight LoRA tuning.

  \item \textbf{SliceGPT} \cite{ashkboos2024slicegpt} --- applies a PCA-style orthogonal transform so signal concentrates in the top components, then deletes the least-important rows and columns of the weight matrices, shrinking the hidden dimension into a smaller dense model.

  \item \textbf{ShortGPT} \cite{men2024shortgpt} --- removes whole layers whose input and output hidden states are nearly identical (low block influence).

  \item \textbf{LaCo} \cite{yang2024laco} --- merges consecutive near-redundant layers into one to reduce depth while preserving behaviour.

  \item \textbf{SparseGPT} \cite{frantar2023sparsegpt} --- one-shot unstructured/semi-structured pruning that removes weights and compensates the rest using the layer-input Hessian; reaches 50\%+ sparsity without retraining.

  \item \textbf{Wanda} \cite{sun2024wanda} --- prunes by weight magnitude times input-activation norm, per output neuron, with no retraining or weight update.

  \item \textbf{Sensitivity-aware mixed sparsity} \cite{shao2024sensitivity} --- allocates a different sparsity ratio to each layer from its Hessian sensitivity, pruning robust layers harder and fragile ones less.

  \item \textbf{SVD-LLM} \cite{wang2025svdllm} --- low-rank factorization via truncation-aware SVD that accounts for how the truncation error propagates to the layer output, so the kept components minimize actual output error.

\end{itemize}

\subsubsection{Stage 2 --- Distillation-Based Recovery}

Aggressive joint compression removes capability that must be restored to stay within the accuracy tolerance. Rather than full retraining, the framework applies lightweight recovery: short-horizon knowledge distillation from the original model as teacher, and/or parameter-efficient (LoRA) fine-tuning on the domain corpus. The pruning-then-distillation pattern is well supported, with compact models derived from a larger pretrained model recovering most of the lost quality using a small fraction of the original training data \cite{muralidharan2024minitron}; reverse-KL distillation is used where open-ended generation quality matters \cite{gu2024minillm}. When the target is fixed and small, distillation can be the primary route to a low-latency model, with quantization layered on top.

\begin{itemize}

  \item \textbf{Minitron (prune + distill)} \cite{muralidharan2024minitron} --- prunes a large model in width and depth, then distills the original (teacher) into the pruned student on a small dataset, producing compact models at a fraction of the original training tokens.

  \item \textbf{MiniLLM (reverse-KL distillation)} \cite{gu2024minillm} --- distills with reverse KL divergence instead of forward KL, which stops the student from placing probability on text the teacher would never produce---better suited to open-ended generation.

  \item \textbf{LoRA recovery / LoRAPrune} \cite{ma2023llmpruner,zhang2024loraprune} --- trains small low-rank adapter matrices on the frozen compressed model to recover accuracy cheaply; LoRAPrune additionally uses the LoRA gradients to decide which structures to prune.

  \item \textbf{PV-Tuning} \cite{malinovskii2024pvtuning} --- post-compression fine-tuning for extreme 1--2-bit models that optimizes the discrete codeword assignments rather than relying on straight-through gradients.

  \item \textbf{Distilling Step-by-Step} \cite{hsieh2023distilling} --- extracts natural-language rationales from the teacher and trains the student to predict both the answer and the rationale, so a small model learns the reasoning and needs less data.

\end{itemize}

\subsubsection{Stage 3 --- KV-Cache and Activation Compression}

This stage addresses the runtime, algorithm-level bottleneck that weight compression alone cannot reach. Because the KV cache grows linearly with sequence length and dominates long-context memory and bandwidth, the framework compresses it through low-bit KV quantization \cite{hooper2024kvquant}, heavy-hitter eviction \cite{zhang2023h2o}, per-head selection of important positions \cite{li2024snapkv}, and speculative prefetching in offloading settings \cite{lee2024infinigen}. Activation quantization---which the survey flags as comparatively neglected---is handled with reordering \cite{yuan2023rptq} or rotation-based outlier removal \cite{ashkboos2024quarot}, enabling genuine end-to-end (e.g. W4A4) savings rather than weight-only compression.

\emph{KV cache:}

\begin{itemize}

  \item \textbf{KVQuant} \cite{hooper2024kvquant} --- quantizes the cache to sub-4-bit using per-channel key quantization (keys have channel-aligned outliers), quantizing keys before the rotary position embedding mixes channels, per-token value quantization, and an outlier split---enabling million-token context.

  \item \textbf{H2O} \cite{zhang2023h2o} --- keeps the few ``heavy-hitter'' tokens that carry most of the attention mass, plus recent tokens, and evicts the rest by accumulated attention score; no retraining.

  \item \textbf{SnapKV} \cite{li2024snapkv} --- reads the attention pattern over the end of the prompt to predict which positions each head will keep attending to, and stores only those KV entries per head.

  \item \textbf{InfiniGen} \cite{lee2024infinigen} --- for offloaded caches (KV in CPU memory), speculatively predicts which entries will receive high attention next step and prefetches only those to the GPU, hiding the transfer latency.

\end{itemize}

\emph{Activations:}

\begin{itemize}

  \item \textbf{RPTQ} \cite{yuan2023rptq} --- reorders and clusters activation channels into groups of similar range and quantizes each group with its own scale, pushing activations toward 3-bit.

  \item \textbf{QuaRot / SpinQuant} \cite{ashkboos2024quarot,liu2025spinquant} --- rotate activations (and the KV cache) so outliers are spread out and vanish, enabling clean 4-bit activation quantization (the same rotation methods used in Stage 1).

\end{itemize}

\subsubsection{Stage 4 --- Deployment-Aware Evaluation and Adaptive Re-Planning}

The compressed artifact is evaluated under one consistent protocol across multiple inference backends---Hugging Face Transformers, llama.cpp, vLLM, and TensorRT-LLM---measuring compression ratio, \emph{real} prefill and decode latency, weight and KV memory, energy per token, and accuracy (perplexity plus a zero-shot task suite). Reporting measured latency across backends, rather than parameter reduction in a single setup, is the framework's response to the survey's central methodological concern about inconsistent, single-setup evaluation. If a metric misses the deployment requirement, the adaptive feedback path returns to Stage 0 and the per-layer plan is revised, realizing the adaptive compression the survey calls for.

\begin{itemize}

  \item \textbf{Inference backends} --- measure the same artifact across Hugging Face Transformers, llama.cpp, vLLM, and TensorRT-LLM.

  \item \textbf{Reported metrics} --- compression ratio, measured prefill and decode latency, weight + KV memory, energy per token, accuracy (perplexity + zero-shot suite).

  \item \textbf{Precision-level sweep} --- compare INT8, INT4, and sub-4-bit settings under one identical protocol.

  \item \textbf{Adaptive re-planning} --- feed measured shortfalls back to Stage 0 to adjust bit-width, pruning ratio, or KV policy.

\end{itemize}

\subsection{Methods for Estimating Weight Importance and Allocating Compression}

Before any layer can be compressed, the framework must decide which parts of it matter and how hard each part can be squeezed, because compressing an important weight too aggressively is what makes accuracy fall. This subsection reviews the ways to make that decision. They fall into three broad approaches, plus one technique that recurs across all of them, and are compared in Table~\ref{tab:lr-importance}. The approaches trade how much they cost to compute against how accurate they are, and the strongest pipelines often combine more than one.

\begin{table}[htbp]
\centering
\scriptsize
\caption{Approaches to estimating weight importance and allocating compression: mechanism, representative methods, computational cost, and typical use.}
\label{tab:lr-importance}
\renewcommand{\arraystretch}{1.2}
\begin{tabular}{p{\dimexpr 0.164\linewidth-2\tabcolsep\relax}p{\dimexpr 0.243\linewidth-2\tabcolsep\relax}p{\dimexpr 0.229\linewidth-2\tabcolsep\relax}p{\dimexpr 0.136\linewidth-2\tabcolsep\relax}p{\dimexpr 0.229\linewidth-2\tabcolsep\relax}}
\hline
\textbf{Approach} & \textbf{Mechanism} & \textbf{Representative methods} & \textbf{Computational cost} & \textbf{Best suited to} \\
\hline
Magnitude & Rank weights by size & Baseline & Very low & Quick baselines; mild pruning \\
Activation-aware & Weight size, scaled by its input size & Wanda \cite{sun2024wanda}, AWQ \cite{lin2024awq} & Low & Fast weight-only low-bit plans \\
First-order (gradient) & Estimated error rise if a weight is removed & LLM-Pruner \cite{ma2023llmpruner} & Low--moderate & Structured pruning of coupled blocks \\
Second-order (Hessian) & Error curvature; adjusts the surviving weights & GPTQ \cite{frantar2023gptq}, SparseGPT \cite{frantar2023sparsegpt}, SlimGPT \cite{ling2024slimgpt} & Moderate--high & Accurate low-bit; one-shot pruning \\
Representational change & How much a layer changes its input & ShortGPT \cite{men2024shortgpt}, LaCo \cite{yang2024laco} & Low & Depth pruning / layer dropping \\
Distribution transform & Spread outliers out before quantizing & SmoothQuant \cite{xiao2023smoothquant}, QuaRot \cite{ashkboos2024quarot}, QuIP\# \cite{tseng2024quipsharp} & Moderate & Activation and KV quantization (W4A4) \\
Learned / search & Search or optimize the settings directly & OmniQuant \cite{shao2024omniquant}, AQLM \cite{egiazarian2024aqlm}, DayPQ \cite{xu2026daypq} & High & Extreme compression; adaptive plans \\
\hline
\end{tabular}
\end{table}

\subsubsection{Importance Scoring and Selection}

The most direct approach gives each weight, channel, or layer an importance score, then compresses the low-scoring parts hard and protects the high-scoring ones. The scoring methods range from cheap and rough to expensive and accurate. \emph{Magnitude} scoring ranks weights by size, assuming the largest matter most; it is the cheapest option but ignores how a weight is actually used. \emph{Activation-aware} scoring also looks at the size of the inputs a weight multiplies, so a small weight that handles large inputs is still treated as important---used by Wanda \cite{sun2024wanda} and AWQ \cite{lin2024awq}. \emph{Gradient-based} scoring estimates how much the model's error would rise if a weight were removed; LLM-Pruner \cite{ma2023llmpruner} applies this to groups such as attention heads and channels. \emph{Hessian-based} scoring uses the curvature of the error: near a trained model the slope is flat, so it is the curvature that says how much a small change will hurt---weights in flat directions can be compressed freely, while weights in sharply curved directions must be protected. The same curvature also shows how to adjust the remaining weights to make up for the ones that change, which is the basis of GPTQ \cite{frantar2023gptq}, SparseGPT \cite{frantar2023sparsegpt}, and SlimGPT \cite{ling2024slimgpt}. Because the exact curvature is far too large to compute for billions of weights, it is approximated one layer at a time from a small calibration set, with a cheaper squared-gradient (Fisher) version as a middle ground. Finally, \emph{layer-level} scoring rates whole layers by how much they change the data passing through them: ShortGPT \cite{men2024shortgpt} removes layers that barely change it, and LaCo \cite{yang2024laco} merges layers that do nearly the same work, which is how the model's depth is reduced.

\subsubsection{Distribution Transformation}

A second approach does not score importance at all. Instead, it reshapes the spread of the numbers so that no part is hard to compress. The usual obstacle is a few outlier values, far larger than the rest, that stretch the range which must fit into only a few bits. SmoothQuant \cite{xiao2023smoothquant} moves this difficulty from the activations onto the weights by rescaling, while QuaRot \cite{ashkboos2024quarot}, SpinQuant \cite{liu2025spinquant}, and the related QuIP and QuIP\# \cite{chee2023quip,tseng2024quipsharp} rotate the values in a way that leaves the model's output unchanged but spreads the outliers across many dimensions. Once the outliers are spread out, ordinary low-bit quantization works without having to single anything out, which is what makes compressing weights, activations, and the KV cache together (for example to 4-bit, W4A4) practical.

\subsubsection{Learned and Search-Based Allocation}

A third approach replaces a fixed scoring rule with optimization, letting the method work out the best settings itself. OmniQuant \cite{shao2024omniquant} learns the best rounding ranges and rescaling automatically. Codebook methods fit the set of allowed values to the model's actual weights instead of using evenly spaced levels---AQLM \cite{egiazarian2024aqlm}, QTIP \cite{tseng2024qtip}, and the non-uniform levels in SqueezeLLM \cite{kim2024squeezellm}. Search-based methods choose the configuration for each layer directly: OOPS \cite{survey2026oops} solves an optimization problem for structured pruning, bypass-backpropagation pruning \cite{gao2025bypass} avoids full backpropagation by using a search-based update, and DayPQ \cite{xu2026daypq} picks the pruning and quantization for each layer on the fly. This group is the closest in spirit to the adaptive planner in Stage 0.

\subsubsection{Outlier Detection and Protection}

Running through all three approaches is one recurring technique: find the small set of outlier values and keep only those in high precision, while compressing everything else hard. LLM.int8() \cite{dettmers2022llmint8} picks out outlier activation values by size, SpQR \cite{dettmers2024spqr} isolates the weights that quantize worst, and SqueezeLLM \cite{kim2024squeezellm} splits the weights into a compressed bulk and a small high-precision remainder. Because the same few values tend to drive both pruning and quantization difficulty, finding and protecting them once is a cheap way to recover most of the accuracy that heavy compression would otherwise cost.

In the proposed framework these are not competing winners but a set of options the Stage 0 planner chooses from, depending on how aggressive the target is and how much planning compute is available. A cheap activation-aware score is enough for moderate weight-only quantization; a curvature-based (or Fisher) score is the natural choice when the per-layer plan must stay accurate at very low bit-widths; and a distribution transform is preferred when activations and the KV cache must also be compressed. Outlier protection can be added on top of any of them. TABLE VII: Approaches to estimating weight importance and computational cost, and typical use.

\subsection{Handling Compute, Memory, and Energy Constraints}

Autoregressive inference has two phases with different cost profiles: the prefill phase, which processes the prompt in parallel, is compute-bound, whereas the decode phase, which generates one token at a time, is memory-bandwidth-bound and governs end-to-end latency. The framework therefore treats compute, memory, and energy as explicit constraints, and reducing the per-token cost along these three axes is what yields low-latency inference.

\begin{itemize}

  \item \textbf{Compute.} Compute demand is reduced both structurally and numerically. Structural pruning and low-rank factorization (Stage 1) cut the dimensions of the weight matrices, lowering the number of floating-point operations directly, while layer pruning shortens the sequential depth each token must traverse. Low-bit (INT4/INT8) arithmetic further reduces the cost of each operation on supporting hardware. These reductions are most consequential in the compute-bound prefill phase.

  \item \textbf{Memory.} Memory is managed as two distinct budgets that scale differently. Static weight memory is bounded by weight quantization and structural pruning (Stage 1); the dynamic KV cache, which grows linearly with sequence length and dominates long-context inference, is bounded by KV quantization, eviction, and per-head selection (Stage 3). Because quantization reduces both the bytes stored and the bytes transferred per token, it relieves memory capacity and the decode-phase bandwidth bottleneck at the same time.

  \item \textbf{Energy.} In inference, energy is dominated by data movement rather than arithmetic, since fetching a value from off-chip memory costs substantially more than the operation that consumes it. The same measures that lower memory traffic---lower-precision weights and KV cache, and fewer parameters---therefore reduce energy in approximate proportion to the traffic removed, with reduced-precision arithmetic contributing a further, smaller saving. Energy is consequently handled as a measured budget rather than optimized in isolation.

\end{itemize}

Because Stage 0 plans pruning and quantization together per layer and the feedback loop searches configurations, the framework selects the most aggressive setting that still satisfies the accuracy tolerance; Stage 4 then confirms the resulting compute, memory, and energy savings as measured values on the target backend rather than as theoretical reductions.

\subsection{Evaluation Strategies}

A minimal study on one open model (e.g. a 7B-class LLaMA/Mistral) would mirror the stage order and, in line with the deployment-aware evaluation gap, run the same artifact across several backends:

\begin{itemize}

  \item \textbf{Baseline.} FP16 perplexity (WikiText-2, C4), a zero-shot suite (e.g. PIQA, WinoGrande, ARC), and measured decode latency, memory, and energy on the target.

  \item \textbf{Plan.} Compute joint per-layer pruning/quantization sensitivity; derive a plan for the declared requirement.

  \item \textbf{Compress.} Apply the plan: e.g. structured slicing/pruning \cite{ma2023llmpruner,ashkboos2024slicegpt} plus 4-bit weight PTQ \cite{frantar2023gptq,lin2024awq}, with the most sensitive layers protected.

  \item \textbf{Recover.} Short LoRA/distillation recovery \cite{muralidharan2024minitron}; confirm accuracy returns toward baseline.

  \item \textbf{KV + activations.} Add KV quantization \cite{hooper2024kvquant} and, if needed, activation rotation \cite{ashkboos2024quarot}.

  \item \textbf{Evaluate across backends.} Load the artifact in HF Transformers, llama.cpp, vLLM, and TensorRT-LLM; report the unified metrics and compare. Iterate via the feedback loop until the requirement is met.

\end{itemize}


\chapter{System Design and Methodology}
\label{Chapter4}

This chapter describes the framework that was built for this thesis and the method used to evaluate it. Every statement about the framework in this chapter describes its code as it was used for the experiments. Section~\ref{sec:ch4-overview} gives the overall design. Sections~\ref{sec:ch4-stage0} to \ref{sec:ch4-stage4} describe the five stages. Section~\ref{sec:ch4-protocol} states the comparison protocol, Section~\ref{sec:ch4-search} describes the constrained search, and Section~\ref{sec:ch4-setup} gives the experimental setup.

\section{Overview of the Framework}
\label{sec:ch4-overview}

The framework is organised as five stages, shown in Figure~\ref{fig:ch4-overview}. Stage~0 measures how sensitive each decoder layer of the model is and turns the measurements into per-layer plans. Stages~1, 2 and 3 compress the model by following those plans: Stage~1 quantizes the weights and the activations, Stage~2 prunes weights, and Stage~3 compresses the KV cache. Stage~4 evaluates the resulting model in the inference engines. A search sits on top of the stages: it proposes which technique each stage should use and which settings Stage~0 should plan with, receives the Stage~4 measurements, and proposes the next configuration.

\begin{figure}[htbp]
\centering
\includegraphics[width=\textwidth]{Figures/ch4-framework-overview.pdf}
\caption{The framework. Stage~0 measures the model and writes per-layer plans; Stages~1 to 3 compress the model following the plans; Stage~4 evaluates the result. The search proposes the techniques and the plan settings and receives the measurements.}
\label{fig:ch4-overview}
\end{figure}

The design separates two questions. Where to compress, that is, which layers keep more bits, which layers are pruned and how much of the cache each layer keeps, is decided once by Stage~0 and written into the plans. How to compress, that is, which published technique does the rounding or the pruning, is decided in Stages~1 to 3. Because of this separation, every technique can be run twice with the same code: once on a uniform plan that treats every layer the same, called \emph{Standard} in this thesis, and once on the Stage~0 plan, called \emph{Our method}. The only difference between the two runs is the plan.

The stages have been tested in two ways, listed in Table~\ref{tab:ch4-paths}. In the single-stage paths was mainly used for unit testing a stage, one of Stages~1 to 3 is applied to a fresh copy of the FP16 model, so that the effect of each stage can be studied on its own. In the sequential path, Stages~1, 2 and 3 are applied one after another to the same model, and only the final model is evaluated.

\begin{table}[htbp]
\centering
\small
\caption{The paths through the framework that are evaluated in this thesis.}
\label{tab:ch4-paths}
\renewcommand{\arraystretch}{1.2}
\begin{tabularx}{\textwidth}{lX}
\hline
\textbf{Path} & \textbf{What is built and measured} \\
\hline
0 $\rightarrow$ 1 $\rightarrow$ 4 & Each quantization technique on the FP16 model, as Standard and as Our method. \\
0 $\rightarrow$ 2 $\rightarrow$ 4 & Each pruning technique on the FP16 model at each pruning level, as Standard and as Our method. \\
0 $\rightarrow$ 3 $\rightarrow$ 4 & Each KV-cache technique on the FP16 model, as Standard and as Our method. \\
0 $\rightarrow$ 1 $\rightarrow$ 2 $\rightarrow$ 3 $\rightarrow$ 4 & One model that is quantized, then pruned, then given a KV-cache technique. This is the path that the search uses. \\
\hline
\end{tabularx}
\end{table}

\section{Stage 0: Sensitivity Profiling and Planning}
\label{sec:ch4-stage0}

Stage~0 takes the FP16 model and a sample of calibration text and produces four things: a sensitivity score for each decoder layer, a weight plan, an activation plan and a KV-cache plan, each with a prediction of the memory it will need. The measurements depend only on the model and the calibration data, so they are stored in a cache on disk and reused by every later run that has the same model, data and seed.

\subsection{Calibration and Evaluation Data}
\label{sec:ch4-data}

The calibration data are 64 windows of 512 tokens, drawn at random positions from the shuffled documents of the training split of WikiText-2 with a fixed seed. The evaluation data are the test split of WikiText-2, joined into one text and cut into windows of 512 tokens that do not overlap. The first half of the windows is the validation set and the second half is the held-out set. Every decision in the framework, including every decision of the search, uses the validation set only. The held-out perplexity is measured and reported next to it, but it is never used to choose anything.

\subsection{Sensitivity Score}
\label{sec:ch4-score}

The default sensitivity score is a direct measurement. For each decoder layer $l$, the layer is bypassed, so that its input is passed straight to the next layer, and the perplexity on the calibration text is measured again. The raw score is the rise in perplexity:
\begin{equation}
  r_l = \mathrm{PPL}_{\mathrm{calib}}(\text{model with layer } l \text{ bypassed}) - \mathrm{PPL}_{\mathrm{calib}}(\text{model}).
  \label{eq:ch4-score}
\end{equation}
This needs one pass over the calibration text for the unchanged model and one pass for each layer. The raw scores are then rank-normalised: the layers are sorted by $r_l$ and layer $l$ receives $s_l = \mathrm{rank}(r_l) / (n - 1)$, where $n$ is the number of layers and tied layers share the average rank. The most sensitive layer therefore has $s_l = 1$ and the most robust has $s_l = 0$. Six other scores are implemented and can be selected in place of the default: the perplexity rise when only layer $l$ is rounded, gradient times weight, Fisher, Taylor, Hessian diagonal and movement scores. The experiments in this thesis use the default score. Perplexity decides where the model is compressed; it is not a guarantee of task accuracy, which is measured separately on the final models.

\subsection{Weight Plan}
\label{sec:ch4-weightplan}

The weight plan gives each decoder layer a bit width $b_l$ and a pruning ratio $p_l$. It is made from the normalised scores with one threshold $\tau$ and one pruning ratio $\rho$:
\begin{equation}
  (b_l,\; p_l) = \left\{
  \begin{array}{ll}
    (8,\; 0) & \text{if } s_l \geq \tau \quad \text{(protected layer)} \\
    (4,\; \rho) & \text{if } s_l < \tau \quad \text{(compressed layer)}
  \end{array}
  \right.
  \label{eq:ch4-plan}
\end{equation}
A second rule, the guard, is applied on top: the five layers with the largest raw score $r_l$ are never pruned, whatever the threshold is. The default values are $\tau = 0.5$ and $\rho = 0.3$; both are settings that the search can change (Section~\ref{sec:ch4-search}). Only the linear weight matrices inside the decoder layers are compressed. The embeddings, the normalisation layers and the output head stay in 16-bit floating point in every plan.

From this rule Stage~0 writes separate plan files for the later stages: a quantization plan with the bit widths only, used by Stage~1, and a pruning plan with the pruning ratios only, used by Stage~2. For the pruning experiments it also writes one pruning plan for each pruning level from 30 to 90 percent, and for each level a second plan that removes exactly as many weights in total as the Standard plan does but places them by sensitivity. In that plan, layer $l$ is pruned at $\min(1,\; c\,(1 - s_l + 1/(n-1)))$, guarded layers are not pruned, and the constant $c$ is chosen so that the total number of removed weights equals the level times the number of decoder weights. Since the bits and the number of weights kept are the same as in the Standard plan, any difference in quality between the two comes only from where the pruning is placed.

\subsection{Activation Plan}
\label{sec:ch4-actplan}

The activation plan gives each layer a bit width for its activations. Stage~0 measures it in the same direct way as the weights: for each layer and for each candidate bit width (4 and 8 bits), the inputs of every linear module of that layer are rounded to that bit width while all other layers are left unchanged, and the rise in calibration perplexity is recorded. The plan then spends an average of 6 bits per layer. Every layer starts at 4 bits, and the layers are raised to 8 bits one at a time, taking first the layer for which the upgrade removes the most measured damage per extra bit, until the budget is used. The Standard activation plan gives every layer 8 bits.

\subsection{KV-cache Plan}
\label{sec:ch4-kvplan}

The KV-cache plan gives each layer a bit width for its cached keys, a bit width for its cached values, and the share of past tokens that the layer keeps. For the bit widths, Stage~0 rounds the keys of one layer (per channel, before the rotary position encoding is applied) or the values of one layer (per token) to 2, 4 and 8 bits in turn and records the rise in calibration perplexity. The bits are then allocated with the same greedy rule as in the activation plan, under an average of 4 bits per cached number, which is the size of the Standard plan (every key and value at 4 bits). For the share of tokens, Stage~0 records, for each layer and for keep ratios of 10, 20, 30, 50 and 75 percent, how much of the layer's attention falls on the most-attended tokens. The layer is given the smallest keep ratio that still covers 95 percent of its attention. The memory of a cache plan is predicted for a context of 2{,}048 tokens and a batch of one.


\section{Stage 1: Quantization}
\label{sec:ch4-stage1}

Stage~1 rounds the weights and the activations of the model to the bit widths in the plan. Its techniques fall into three kinds, listed in Table~\ref{tab:ch4-stage1}. Weight techniques round the weight matrices. Activation techniques round or transform the inputs of the linear modules. A combined technique, written as a weight technique with an activation technique, applies both to one model: the activation technique first rewrites the model, the weight technique then rounds the weights, and the activations are rounded while the model is measured. Each layer of a combined model therefore has its own pair of weight bits and activation bits.

The basic rounding operation is the uniform affine quantizer of Equation~\ref{eq:ch2-quant}, applied per group of $g$ consecutive weights with the scale taken from the minimum and maximum of the group and an integer zero-point. Round-to-nearest applies this operation directly. GPTQ rounds a weight matrix column by column and compensates the error of each column in the columns that are not yet rounded, using second-order information from the calibration inputs of the layer \cite{frantar2023gptq}; in the framework the layers are processed in order, and each layer is calibrated on inputs that have passed through the layers already rounded. AWQ scales the weight channels according to the size of their activations before rounding \cite{lin2024awq}. The other techniques are published methods that were re-implemented in simplified form inside the framework so that they can follow a per-layer plan, as stated in Section~\ref{sec:ch1-scope}; the simplifications of each technique are recorded with its results.

\begin{table}[htbp]
\centering
\small
\caption{Stage~1 techniques. Each weight technique can be combined with round-to-nearest activations or with SmoothQuant.}
\label{tab:ch4-stage1}
\renewcommand{\arraystretch}{1.2}
\begin{tabularx}{\textwidth}{lX}
\hline
\textbf{Kind} & \textbf{Techniques} \\
\hline
Weights & Round-to-nearest; GPTQ \cite{frantar2023gptq}; AWQ \cite{lin2024awq}; OmniQuant \cite{shao2024omniquant} (learnable weight clipping only); SqueezeLLM \cite{kim2024squeezellm}; SpQR \cite{dettmers2024spqr}; EfficientQAT \cite{chen2025efficientqat} (block-wise step only); QuIP \cite{chee2023quip}; QuIP\# \cite{tseng2024quipsharp}; PB-LLM \cite{shang2023pbllm}; BiLLM \cite{huang2024billm} (Standard only); the 4-bit and 8-bit storage of the bitsandbytes library \\
Activations & Round-to-nearest activations; SmoothQuant \cite{xiao2023smoothquant}; QuaRot \cite{ashkboos2024quarot}; SpinQuant \cite{liu2025spinquant}; RPTQ \cite{yuan2023rptq} (Standard only) \\
\hline
\end{tabularx}
\end{table}

Two remarks apply to this table. First, EfficientQAT is included only with its block-wise step, which trains the quantization parameters of one block at a time on the calibration data; its end-to-end training step is not run. Second, in the Standard run a weight technique uses 4 bits in every layer, and a combined technique uses 4-bit weights and 8-bit activations in every layer; in the Our method run the same technique uses the Stage~0 quantization plan and, for a combined technique, the Stage~0 activation plan.

In Stages~1 to 3 the compression is simulated in Hugging Face Transformers: the rounded values are written back into 16-bit floating-point tensors. The quality measured on this model is the quality of the rounding itself. Its size and speed are not real, which is why the model is exported to the inference engines in Stage~4.

\section{Stage 2: Pruning}
\label{sec:ch4-stage2}

Stage~2 removes weights and rounds nothing. It has three techniques.

\begin{itemize}
  \item \textbf{Unstructured pruning} uses the Wanda criterion: within each output row of a weight matrix, the weights with the smallest product of weight magnitude and input activation norm are set to zero, with no weight update \cite{sun2024wanda}.
  \item \textbf{Structured pruning} removes whole inner channels of the feed-forward network of a layer, that is, a row of each input matrix and the matching column of the output matrix. Channels are removed in increasing order of the product of the channel's activation norm and the norm of its output column, until the planned share of the layer's weights is removed. At most 75 percent of a layer is removed, and attention heads are never removed.
  \item \textbf{Low-rank factorisation} replaces each weight matrix by a product of two smaller matrices obtained from an activation-aware singular value decomposition, following the ASVD idea \cite{yuan2025asvd}. The rank is chosen so that the layer keeps the planned share of its numbers. This technique cannot follow Stage~1 and is therefore used only in the path 0 $\rightarrow$ 2 $\rightarrow$ 4.
\end{itemize}

Each technique is run at every pruning level from 30 to 90 percent. In the Standard run the level is removed from every layer. In the Our method run the level is removed from the layers that Stage~0 found robust, and nothing is removed from the protected and guarded layers. A third run uses the same-size plan of Section~\ref{sec:ch4-weightplan}, which removes the same total as the Standard run. In addition, Stage~0 runs a pruning-levels test in which each technique prunes the FP16 model at each level on three sets of layers (all layers, all layers except the guarded ones, and the guarded layers only), which shows how much of the damage of pruning comes from the guarded layers.

Removed weights are stored as zeros. As stated in Section~\ref{sec:ch1-scope}, the pruning experiments therefore measure quality only.

\section{Stage 3: KV-cache Compression}
\label{sec:ch4-stage3}

Stage~3 compresses the KV cache and leaves the weights as they are. Its techniques, listed in Table~\ref{tab:ch4-stage3}, are of two kinds: rounding techniques store the cached keys and values in fewer bits, and eviction techniques keep only part of the past tokens. All of them are implemented inside the attention computation of the Hugging Face model: keys and values are rounded in place, and evicted tokens are masked out of the attention. The cache memory reported for a technique is the prediction of its plan, for a context of 2{,}048 tokens and a batch of one.

\begin{table}[htbp]
\centering
\small
\caption{Stage~3 techniques as implemented in the framework.}
\label{tab:ch4-stage3}
\renewcommand{\arraystretch}{1.2}
\begin{tabularx}{\textwidth}{lX}
\hline
\textbf{Technique} & \textbf{Implementation in the framework} \\
\hline
Round-to-nearest KV & Keys rounded per channel before the position encoding, values per token, with the same rounding as the Stage~0 measurement. \\
QuaRot KV \cite{ashkboos2024quarot} & Keys and values rotated by a Hadamard matrix and then rounded. \\
KVQuant \cite{hooper2024kvquant} & Keys per channel before the position encoding and values per token, with non-uniform levels found by k-means and one percent of outliers left unrounded. The levels are not weighted by sensitivity as in the paper. \\
H2O \cite{zhang2023h2o} & A cache of fixed size per layer, half for the most recent tokens and half for the heavy-hitter tokens. \\
SnapKV \cite{li2024snapkv} & An observation window of the latest tokens and the earlier tokens that the window attends to most, in a cache of the same fixed size as H2O. \\
InfiniGen \cite{lee2024infinigen} & The tokens to keep are predicted from 30 percent of the channels, in a cache of the same fixed size as H2O. The offloading of the full cache to CPU memory described in the paper is not run. \\
\hline
\end{tabularx}
\end{table}

In the Standard run a rounding technique stores every key and value at 4 bits and keeps every token, and an eviction technique keeps 20 percent of the past tokens at full precision. In the Our method run each technique follows the Stage~0 KV-cache plan: the planned key bits, value bits and share of tokens for each layer.

\section{Stage 4: Evaluation}
\label{sec:ch4-stage4}

Stage~4 measures a model. It always measures the simulated model in Hugging Face Transformers \cite{wolf2020transformers}, and it exports the model to llama.cpp \cite{ggml2025llamacpp} and vLLM \cite{kwon2023vllm} to measure real file size, memory and speed. An export to TensorRT-LLM \cite{nvidia2025tensorrtllm} is also implemented, but no TensorRT-LLM measurement is reported in this thesis.

\subsection{Measurements in Hugging Face Transformers}
\label{sec:ch4-hf}

Perplexity is computed with Equation~\ref{eq:ch2-ppl} over the validation windows and, separately, over the held-out windows. Latency is measured at batch size one with a random prompt of 128 tokens: the prefill time is the time to process the prompt, and the decode latency is the time per token over 64 greedy decoding steps with a pre-allocated KV cache. On a GPU the prefill and one decoding step are each captured once as a CUDA graph and replayed, so that the measured time is the work of the GPU and not the overhead of Python; when a technique prevents the capture, the same loop is run without it and the result records which mode ran. Each measurement is preceded by 3 warm-up runs and repeated 10 times, and the mean and the standard deviation are reported. The model size and the peak GPU memory are also recorded. Zero-shot accuracy on ARC-Easy, PIQA, HellaSwag and WinoGrande (Section~\ref{sec:ch2-metrics}) is measured with the lm-evaluation-harness library on the model as compressed, so that the tasks see exactly what a technique changed.

\subsection{Export to the Inference Engines}
\label{sec:ch4-export}

Table~\ref{tab:ch4-engines} summarises how a model reaches each engine. In every case the per-layer bit widths of the plan are kept in the exported file, and the embeddings and the output head stay in 16-bit floating point.

\begin{table}[htbp]
\centering
\small
\caption{How a compressed model is exported to and measured in each inference engine.}
\label{tab:ch4-engines}
\renewcommand{\arraystretch}{1.2}
\begin{tabularx}{\textwidth}{lX}
\hline
\textbf{Engine} & \textbf{Export and measurement} \\
\hline
llama.cpp & The model is converted to GGUF \cite{ggml2025gguf} and quantized with llama.cpp's own quantization types, each decoder layer at its planned bit width. Speed is measured with the llama-bench program and perplexity with the llama-perplexity program, which scores the second half of each window. \\
vLLM, weights only & The model is saved as a GPTQ-format checkpoint with per-layer bit widths of 2, 3, 4 or 8 bits. The weights are re-rounded onto symmetric per-group grids, because vLLM loads only symmetric grids in this format. \\
vLLM, weights and activations & The model is saved as a compressed-tensors checkpoint in which each decoder layer has its own weight bits and activation bits; vLLM rounds the activations while the model runs. \\
TensorRT-LLM (implemented, not measured in this thesis) & The same GPTQ-format checkpoint is built into a TensorRT engine, with one bit width or with a per-layer mixed-precision configuration. \\
\hline
\end{tabularx}
\end{table}

vLLM and TensorRT-LLM run in Python environments of their own, separate from the framework, and are started by the framework as worker processes. In both, the latency is measured with the same prompt length, number of decoded tokens, warm-up runs and repeats as in Hugging Face: the prefill time is the time of a request that generates one token after the 128-token prompt, and the decode latency is the time of a request that generates 64 tokens minus the prefill time, divided by 63. Prefix caching is switched off so that every prefill is a real one.

Three properties of the exports matter for reading the results in Chapter~5. First, each engine re-rounds the weights onto its own grid, so a perplexity measured in an engine is comparable with other perplexities measured in the same engine, and not with the Hugging Face perplexity. Second, in vLLM the weight-only models and the models with rounded activations run on different kernels: in the experiments of this thesis a weight-only model decodes in about 2~ms per token, while a model with rounded activations decodes in about 35 to 42~ms per token. Decode latencies of the two kinds are therefore comparable only within each kind. Third, some models cannot be exported: models rewritten by QuaRot or SpinQuant need a rotation at run time that the export formats do not carry, and the KV-cache techniques of Stage~3 exist only in the Hugging Face model. Models of the path 0 $\rightarrow$ 3 $\rightarrow$ 4 are therefore measured in Hugging Face only.

\section{Comparison Protocol}
\label{sec:ch4-protocol}

Every experiment in the single-stage paths reports three kinds of row for each technique:

\begin{itemize}
  \item \textbf{FP16}: the uncompressed model.
  \item \textbf{Standard}: the technique on the uniform plan (for example 4-bit weights in every layer, or 4-bit weights and 8-bit activations in every layer).
  \item \textbf{Our method}: the same technique but with Stage~0's plan output.
\end{itemize}

The three rows are produced under identical conditions: the same calibration windows, the same seed, the same evaluation windows, the same GPU and the same engine versions. Every row starts from a fresh copy of the FP16 model, so that no row is affected by an earlier one. Results that are expensive to compute, such as the FP16 measurements and the sensitivity measurements, are stored in a cache whose key contains every setting that determines them, so that a result is reused only when all of those settings are equal. A row that fails is kept in the results with its error message; it is not removed.

\section{Constrained Multi-objective Search}
\label{sec:ch4-search}

The search answers the question of which combination of techniques and settings to use for a given deployment requirement. Figure~\ref{fig:ch4-search} shows one search run.

\begin{figure}[htbp]
\centering
\includegraphics[width=0.95\textwidth]{Figures/ch4-search-loop.pdf}
\caption{One run of the search. Each try builds one model through Stages~1, 2 and 3 in sequence and measures it. After the last try the final model is evaluated in the engines and on the accuracy tasks.}
\label{fig:ch4-search}
\end{figure}

\subsection{Deployment Requirement and Objectives}
\label{sec:ch4-requirement}

A deployment requirement is a set of upper limits. The framework supports limits on validation perplexity, on decode latency, on weight memory and on KV-cache memory. The requirement used in this thesis has two limits: a validation perplexity of at most 11.5 and a vLLM decode latency of at most 2.5~ms per token. The perplexity limit can be checked on the Hugging Face model or on the model as loaded in vLLM; Chapter~5 states which was used for each search. For a model that misses a limit, the framework records the shortfall, that is, by how much each limit was missed.

Within the limits the search minimises up to three objectives: validation perplexity, vLLM decode latency and predicted weight memory. The objectives of each search run are given with its results in Chapter~5. The held-out perplexity cannot be selected as an objective or as a limit.

Written as a formula, let $x = (t_1, t_2, t_3, \theta)$ be a point of the search space $\mathcal{X}$ (Section~\ref{sec:ch4-space}), with $t_i$ the technique of Stage~$i$ and $\theta$ the numeric settings, and let $P(x)$ be the validation perplexity, $L(x)$ the vLLM decode latency in ms per token and $S(x)$ the predicted weight memory in GB of the model built from $x$. For the requirement of this thesis the search solves
\begin{equation}
  \min_{x \in \mathcal{X}} \; \bigl(f_1(x), \ldots, f_k(x)\bigr) \quad \text{subject to} \quad P(x) \le P_{\max}, \;\; L(x) \le L_{\max},
  \label{eq:ch4-search}
\end{equation}
where the objectives $f_1, \ldots, f_k$ ($k \le 3$) are taken from $P$, $L$ and $S$ in a stated order, $P_{\max} = 11.5$ and $L_{\max} = 2.5$~ms per token.

\subsection{Search Space}
\label{sec:ch4-space}

A point of the search space consists of one technique for each of Stages~1, 2 and 3 and of the numeric settings in Table~\ref{tab:ch4-space}. The Stage~1 choices are the weight techniques round-to-nearest, GPTQ, AWQ, OmniQuant, SqueezeLLM, SpQR, EfficientQAT, QuIP, QuIP\#, PB-LLM and bitsandbytes, each combined with round-to-nearest activations or with SmoothQuant, and in addition round-to-nearest, GPTQ, AWQ, OmniQuant and EfficientQAT as weight-only techniques. The Stage~2 choices are unstructured and structured pruning. The Stage~3 choices are the six techniques of Table~\ref{tab:ch4-stage3}.

\begin{table}[htbp]
\centering
\small
\caption{Numeric settings of the search space.}
\label{tab:ch4-space}
\renewcommand{\arraystretch}{1.2}
\begin{tabularx}{\textwidth}{Xll}
\hline
\textbf{Setting} & \textbf{Full range} & \textbf{Range set for the search} \\
\hline
Sensitivity threshold $\tau$ & 0.1 to 0.9 & 0.8 to 0.9 \\
Pruning ratio $\rho$ of compressed layers & 0.0 to 0.6 & 0.0 to 0.3 \\
Calibration dataset & WikiText-2, C4, Pile (10k sample) & as full range \\
Number of calibration windows & 16, 32, 64, 128 & as full range \\
Quantization group size & 32, 64, 128, per channel & as full range \\
SmoothQuant migration strength & 0 to 1 & as full range \\
\hline
\end{tabularx}
\end{table}

\subsection{Evaluating One Point}
\label{sec:ch4-point}

For each point, Stage~0 makes the plans for the point's settings. The weight memory that the plan would need is estimated first; a point that is over a memory limit is not built. Otherwise one model is built in sequence: the Stage~1 technique quantizes the FP16 model following the quantization and activation plans, the Stage~2 technique prunes that model following the pruning plan, and the Stage~3 technique is applied to the result following the KV-cache plan. Only the Our method side is built. The model is then measured in Hugging Face, and also in vLLM when an objective or a limit is a vLLM measurement. llama.cpp and TensorRT-LLM are not run during the tries.

Two mechanisms prevent the same model from being built twice. Within a run, a model is built once per plan: two thresholds that fall between the same two layers give the same plan and therefore share one model. Across runs, every combination that has been tried is written to a log file that is kept with the cache. When a later run proposes a combination that is already in the log, the logged measurements are used, the requirement is checked again, and nothing is built. A search can therefore be extended with more tries without repeating earlier ones.

\subsection{Optimiser}
\label{sec:ch4-optimiser}

The optimiser is multi-objective Bayesian optimisation (Section~\ref{sec:ch2-moo}), for which the framework uses the Gaussian-process sampler of the Optuna library. All objectives are minimised, and the requirement is passed to the sampler as a constraint whose value is the sum of the shortfalls of the point. For the two limits of Equation~\ref{eq:ch4-search} it is
\begin{equation}
  c(x) = \max\bigl(0, P(x) - P_{\max}\bigr) + \max\bigl(0, L(x) - L_{\max}\bigr),
  \label{eq:ch4-constraint}
\end{equation}
which is zero for a point that meets every limit. A point that was not built because its plan is over a limit, or that failed while being built or measured, is reported to the sampler with the worst objective values seen so far and a positive constraint value, so that the sampler learns to avoid that region. The run ends when the budget of tries is used. A point that is not built because its plan is over a limit does not use up the budget.

Two other optimisers are implemented in the framework: successive halving \cite{bischl2021hpo}, which evaluates many points on a fraction of the evaluation windows and keeps the best third at each of three rounds, and NSGA-III through the pymoo library \cite{blank2020pymoo}. The budget is counted in full evaluations, that is, models measured on all validation windows, and a try on a fraction of the windows uses that fraction of one evaluation. For the same budget, successive halving therefore builds more models than the Bayesian optimiser: it first measures random points on one ninth of the windows, then the best third of them on one third of the windows, and then the best third of those on all windows. The search results in Chapter~5 come from the Bayesian optimiser.

\subsection{Final Model}
\label{sec:ch4-final}

After the last try, the final model is selected among all tries, of this run and of earlier runs in the log, that were measured on the full validation set and meet every limit. The selected try is the one with the lowest value of the first objective; ties are broken by the second objective and then by the third. This model is built again and its exported files are kept. It is then evaluated in the engines that are installed on the machine of the run (Hugging Face and vLLM for the searches of Chapter~5) and on the four accuracy tasks. The FP16 model is evaluated on the same tasks in the same run, so that the two accuracies are measured under the same conditions. When no try meets every limit, no final model is selected, and the report lists the tries with their shortfalls.

\section{Experimental Setup}
\label{sec:ch4-setup}

All experiments use TinyLlama-1.1B-Chat-v1.0 (Section~\ref{sec:ch2-tinyllama}) and the Stage~0 settings given in Section~\ref{sec:ch4-stage0}. Table~\ref{tab:ch4-setup} lists the hardware and software. The seed fixes the random number generators of Python, NumPy and PyTorch, and the deterministic mode of PyTorch is switched on where a deterministic implementation exists. Each configuration is run with this one seed (Section~\ref{sec:ch1-scope}).

\begin{table}[htbp]
\centering
\small
\caption{Hardware and software of the experiments.}
\label{tab:ch4-setup}
\renewcommand{\arraystretch}{1.2}
\begin{tabularx}{\textwidth}{lX}
\hline
\textbf{Item} & \textbf{Value} \\
\hline
GPUs & NVIDIA A100 on Google Colab; NVIDIA RTX~5070~Ti 16~GB in the host PC \\
Framework & Python 3.10 or later, PyTorch 2.1 or later, Transformers 4.56 or later \\
Search & Optuna 5 or later (Gaussian-process sampler); pymoo 0.6 or later (NSGA-III option) \\
Accuracy tasks & lm-evaluation-harness 0.4 or later \\
vLLM & version 0.31.0, in its own Python environment (Linux only) \\
TensorRT-LLM & version 1.2, in its own Python environment (Linux only); not measured in this thesis \\
llama.cpp & release programs with CUDA support \\
\hline
\end{tabularx}
\end{table}

Because vLLM runs on Linux only, every measurement that involves it, including all search runs, was made on the A100. Every run writes its configuration, a log, a table of all rows and a report to its own folder, and the tables and figures of Chapter~5 are taken from these files.

\section{Summary}
\label{sec:ch4-summary}

The framework separates where to compress from how to compress. Stage~0 measures each decoder layer directly, by the rise in calibration perplexity, and writes per-layer plans for the weights, the activations, pruning and the KV cache together with a prediction of their memory. Stages~1 to 3 apply published techniques that follow those plans, and each technique is also run on a uniform plan so that Standard and Our method differ only in the plan. Stage~4 measures the models in Hugging Face Transformers and, after export, in llama.cpp and vLLM. A constrained Bayesian search chooses one technique per stage and the plan settings under limits on perplexity and decode latency, builds each candidate as one model through Stages~1 to 3 in sequence, never builds the same combination twice, and selects a final model for full evaluation. Chapter~5 reports the results obtained with this method.
\clearpage
\chapter{Results and Discussion}
\label{Chapter5}

This chapter reports what was measured with the framework of Chapter~\ref{Chapter4} and discusses it. Section~\ref{sec:ch5-experiments} lists the experiments and gives the measurements of the uncompressed model. Section~\ref{sec:ch5-stage0} reports the Stage~0 profile. Sections~\ref{sec:ch5-quant}, \ref{sec:ch5-prune} and \ref{sec:ch5-kv} report the three single-stage paths, and Section~\ref{sec:ch5-search} reports the sequential path and the search. Section~\ref{sec:ch5-checks} tests the score, the per-layer plan and the contribution of each stage, and Section~\ref{sec:ch5-published} compares the final models with published work. Section~\ref{sec:ch5-discussion} answers the research questions of Chapter~1 and Section~\ref{sec:ch5-limits} states the limitations of the evidence. In every table and figure a lower perplexity, a smaller size and a lower latency are better, and a higher accuracy and a higher number of tokens per second are better.

\section{Experiments and the FP16 Reference}
\label{sec:ch5-experiments}

Table~\ref{tab:ch5-experiments} lists the experiments. All of them use TinyLlama, the seed 0 and the data of Section~\ref{sec:ch4-data}. The single-stage experiments use the default Stage~0 settings, a sensitivity threshold $\tau = 0.5$ and, for quantization, a group size of 128 weights.

\begin{table}[htbp]
\centering
\small
\caption{The experiments reported in this chapter.}
\label{tab:ch5-experiments}
\renewcommand{\arraystretch}{1.2}
\begin{tabularx}{\textwidth}{>{\raggedright\arraybackslash}p{2.9cm}X>{\raggedright\arraybackslash}p{1.7cm}>{\raggedright\arraybackslash}p{2.7cm}}
\hline
\textbf{Path} & \textbf{What was run} & \textbf{GPU} & \textbf{Engines measured} \\
\hline
0 $\rightarrow$ 1 $\rightarrow$ 4 & 11 weight techniques and 12 combined techniques, as Standard and as Our method & RTX~5070~Ti & Hugging Face, llama.cpp \\
0 $\rightarrow$ 2 $\rightarrow$ 4 & 3 pruning techniques at levels of 30 to 90 percent, as Standard, as Our method and with the same-size plan (61 rows; structured pruning up to 80 percent) & A100 & Hugging Face, llama.cpp, vLLM \\
0 $\rightarrow$ 3 $\rightarrow$ 4 & 6 KV-cache techniques, as Standard and as Our method & A100 & Hugging Face \\
0 $\rightarrow$ 1 $\rightarrow$ 2 $\rightarrow$ 3 $\rightarrow$ 4 & Search runs of 30 or 40 tries; each try is one model built through Stages~1, 2 and 3 & A100 & Hugging Face, vLLM \\
Checks (Section~\ref{sec:ch5-checks}) & Each layer changed alone; bit plans of equal size for 4 weight techniques; stage ablation of the two final models found under both limits; the compressed caches of the engines & A100 & Hugging Face, llama.cpp, vLLM \\
\hline
\end{tabularx}
\end{table}

The RTX~5070~Ti is the GPU of the host PC. The row count includes the FP16 row. No TensorRT-LLM measurement is reported in this chapter; the measured engines are Hugging Face Transformers, llama.cpp and vLLM.

Table~\ref{tab:ch5-fp16} gives the measurements of the uncompressed FP16 model, which is the reference for everything that follows. As explained in Section~\ref{sec:ch4-export}, llama.cpp scores only the second half of each window, so its perplexity (8.12) is comparable with other llama.cpp perplexities and not with the Hugging Face and vLLM value (10.18).

\begin{table}[htbp]
\centering
\small
\caption{The uncompressed FP16 model. Perplexity is given on the validation set and on the held-out set.}
\label{tab:ch5-fp16}
\renewcommand{\arraystretch}{1.2}
\begin{tabularx}{\textwidth}{lX}
\hline
\textbf{Engine} & \textbf{Measurement} \\
\hline
Hugging Face & Perplexity 10.18 / 10.62; decode 4.32~ms per token and prefill 5.87~ms on the RTX~5070~Ti; decode 4.24~ms per token on the A100; peak GPU memory 2.41~GB \\
llama.cpp & File 2.20~GB; perplexity 8.12 / 8.56; 324 tokens per second on the RTX~5070~Ti and 421 on the A100 \\
vLLM (A100) & File 2.20~GB; perplexity 10.18 / 10.62; decode 2.54~ms per token; prefill 9.16~ms; peak GPU memory 3.77~GB \\
Accuracy & ARC-Easy 0.548, PIQA 0.745, HellaSwag 0.604, WinoGrande 0.603; mean 0.625 \\
\hline
\end{tabularx}
\end{table}

\section{Stage 0: Layer Sensitivity}
\label{sec:ch5-stage0}

Figure~\ref{fig:ch5-sensitivity} shows the raw sensitivity score of Equation~\ref{eq:ch4-score} for the 22 decoder layers. The perplexity of the unchanged model on the calibration text is 14.12. Bypassing layer~0 raises it by 1{,}176 and bypassing layer~2 by 368. Three more layers have a rise above 10: layer~7 (36.5), layer~21 (23.6) and layer~1 (11.6). For the other 17 layers the rise is between 1.4 and 5.7. The sensitivity is therefore very uneven: the largest score is about 800 times the smallest, and the five largest scores belong to the first three layers, one middle layer and the last layer. These five layers are the guarded layers, which are never pruned.

\begin{figure}[htbp]
\centering
\includegraphics[width=0.95\textwidth]{Figures/ch5-sensitivity.pdf}
\caption{Stage~0 sensitivity score of each decoder layer (log scale; a higher bar is a more sensitive layer) and the weight plan at the default threshold $\tau = 0.5$.}
\label{fig:ch5-sensitivity}
\end{figure}

At the default threshold the weight plan protects 11 of the 22 layers (layers 0, 1, 2, 3, 7, 15, 16, 18, 19, 20 and 21), which keep 8 bits and are not pruned; the other 11 layers get 4 bits and the pruning ratio. The activation plan gives 8 bits to 11 layers and 4 bits to the other 11, an average of 6 bits. The KV-cache plan keeps the keys and values of every layer at 4 bits, except the values of layer~0 at 2 bits, and lets each layer keep between 10 and 100 percent of the past tokens; over the whole cache it keeps 25.9 percent of the tokens at an average of 3.82 bits.

The profile is cheap. The 23 passes over the calibration text that give the scores took 96~s on the RTX~5070~Ti and 121~s on the A100. The activation measurements took 188~s and building the KV-cache plan took 541~s on the RTX~5070~Ti. All of them are computed once and reused from the cache by every later run.

\section{Quantization (Path 0 $\rightarrow$ 1 $\rightarrow$ 4)}
\label{sec:ch5-quant}

\subsection{Weights Only}

Table~\ref{tab:ch5-quant-wo} and Figure~\ref{fig:ch5-quant} compare the eleven weight techniques on the Standard plan (4 bits in every layer) and on the Stage~0 plan (8 bits in the 11 protected layers, 4 bits in the others).

\begin{table}[htbp]
\centering
\small
\caption{Weight-only quantization, Standard (Std) against Our method (Ours), measured on the RTX~5070~Ti. Perplexity: lower is better. Tokens per second: higher is better. The llama.cpp file is 0.81~GB for every Standard row and 1.05~GB for every Our method row; the FP16 file is 2.20~GB.}
\label{tab:ch5-quant-wo}
\renewcommand{\arraystretch}{1.15}
\begin{tabularx}{\textwidth}{Xrrrrrrrr}
\hline
 & \multicolumn{2}{c}{\textbf{Hugging Face}} & \multicolumn{2}{c}{\textbf{Hugging Face}} & \multicolumn{2}{c}{\textbf{llama.cpp}} & \multicolumn{2}{c}{\textbf{llama.cpp}} \\
 & \multicolumn{2}{c}{\textbf{validation}} & \multicolumn{2}{c}{\textbf{held-out}} & \multicolumn{2}{c}{\textbf{validation}} & \multicolumn{2}{c}{\textbf{tokens/s}} \\
\textbf{Technique} & Std & Ours & Std & Ours & Std & Ours & Std & Ours \\
\hline
FP16 & \multicolumn{2}{c}{10.18} & \multicolumn{2}{c}{10.62} & \multicolumn{2}{c}{8.12} & \multicolumn{2}{c}{324} \\
\hline
RTN & 10.97 & 10.41 & 11.39 & 10.84 & 9.02 & 8.38 & 651 & 551 \\
GPTQ & 10.68 & 10.31 & 11.12 & 10.73 & 8.76 & 8.27 & 647 & 559 \\
AWQ & 10.66 & 10.37 & 11.08 & 10.80 & 8.84 & 8.38 & 655 & 559 \\
OmniQuant & 10.90 & 10.39 & 11.34 & 10.82 & 8.95 & 8.37 & 654 & 532 \\
SqueezeLLM & 10.60 & 10.33 & 11.03 & 10.76 & 8.82 & 8.30 & 654 & 552 \\
SpQR & 10.27 & 10.21 & 10.71 & 10.65 & 8.63 & 8.23 & 632 & 539 \\
EfficientQAT & 10.51 & 10.25 & 10.94 & 10.70 & 8.58 & 8.24 & 654 & 553 \\
QuIP & 10.42 & 10.30 & 10.89 & 10.74 & 8.74 & 8.31 & 654 & 554 \\
QuIP\# & 10.37 & 10.27 & 10.82 & 10.70 & 8.73 & 8.31 & 655 & 553 \\
PB-LLM & 11.70 & 10.48 & 12.11 & 10.92 & 10.02 & 8.50 & 652 & 556 \\
bitsandbytes & 10.74 & 10.30 & 11.15 & 10.74 & 8.80 & 8.32 & 624 & 550 \\
\hline
\end{tabularx}
\end{table}

\begin{figure}[htbp]
\centering
\includegraphics[width=0.9\textwidth]{Figures/ch5-quant-dumbbell.pdf}
\caption{Weight-only quantization: Hugging Face validation perplexity of each technique on the Standard plan (circle) and with Our method (square). Further left is better. The dashed line is the FP16 model.}
\label{fig:ch5-quant}
\end{figure}

\begin{figure}[htbp]
\centering
\includegraphics[width=0.9\textwidth]{Figures/ch5-quant-gain.pdf}
\caption{How much lower the validation perplexity of each weight technique is with Our method than on the Standard plan, in percent of the Standard value (a longer bar is a larger gain).}
\label{fig:ch5-quant-gain}
\end{figure}

\textbf{What this shows.} Our method has a lower perplexity than Standard for all eleven techniques, on the validation set and on the held-out set, in Hugging Face and in llama.cpp. The gain is largest for the techniques that lose most at 4 bits: PB-LLM improves from 11.70 to 10.48 and round-to-nearest from 10.97 to 10.41. It is smallest for the technique that is already close to FP16: SpQR improves from 10.27 to 10.21. With Our method every technique is between 10.21 and 10.48, so the choice of technique matters much less than on the Standard plan, where the range is 10.27 to 11.70. Figure~\ref{fig:ch5-quant-gain} gives the same gain in percent: 10.4 percent for PB-LLM, 5.1 percent for round-to-nearest and 0.6 percent for SpQR.

\textbf{What it costs.} The Our method file is larger, 1.05~GB against 0.81~GB, because half of the layers keep 8 bits. It is also slower in llama.cpp: 532 to 559 tokens per second against 624 to 655. Both are still well ahead of the FP16 model, which needs 2.20~GB and reaches 324 tokens per second on the same GPU. The comparison is therefore not at equal size. Our method buys back a part of the quality lost at 4 bits (for round-to-nearest, 71 percent of the gap between the 4-bit model and FP16) for 0.24~GB of extra file size and about 15 percent of the llama.cpp speed. Figure~\ref{fig:ch5-quant-size} shows this trade of quality against size for every technique.

\begin{figure}[htbp]
\centering
\includegraphics[width=0.9\textwidth]{Figures/ch5-quant-size.pdf}
\caption{Quality against size for weight-only quantization. Each line joins one technique on the Standard plan (circle, 0.81~GB) to the same technique with Our method (square, 1.05~GB). Lower is better for perplexity and smaller is better for size. The dashed line is the FP16 model.}
\label{fig:ch5-quant-size}
\end{figure}

\textbf{What Hugging Face does not show.} The decode latency of the simulated model in Hugging Face is between 4.09 and 4.22~ms per token for every weight-only row, against 4.32~ms for FP16. The simulated model stores rounded values in 16-bit tensors, so it cannot be faster or smaller; the real gain appears only in the exported file.

\subsection{Weights and Activations}

Table~\ref{tab:ch5-quant-wa} and Figure~\ref{fig:ch5-quant-wa} report the combined techniques.

\begin{figure}[htbp]
\centering
\includegraphics[width=0.9\textwidth]{Figures/ch5-quant-wa-dumbbell.pdf}
\caption{Weights and activations: Hugging Face validation perplexity of each combined technique on the Standard plan (circle) and with Our method (square). Further left is better. The dashed line is the FP16 model.}
\label{fig:ch5-quant-wa}
\end{figure}
 The Standard row uses 4-bit weights and 8-bit activations in every layer. The Our method row uses the Stage~0 quantization plan and the Stage~0 activation plan, which averages 6 bits. A second Our method row uses the joint weight and activation plan, and a second Standard row uses 8-bit weights and 8-bit activations in every layer.

\begin{table}[htbp]
\centering
\small
\caption{Weights and activations: Hugging Face validation perplexity (lower is better), measured on the RTX~5070~Ti. Each weight technique is combined with round-to-nearest activations unless stated.}
\label{tab:ch5-quant-wa}
\renewcommand{\arraystretch}{1.15}
\begin{tabularx}{\textwidth}{Xrrrr}
\hline
\textbf{Weight technique} & \textbf{Standard} & \textbf{Our method} & \textbf{Our method,} & \textbf{Standard,} \\
 & \textbf{(4-bit weights,} & \textbf{(weight plan,} & \textbf{joint plan} & \textbf{8-bit weights,} \\
 & \textbf{8-bit activations)} & \textbf{activation plan)} & & \textbf{8-bit activations} \\
\hline
RTN & 10.97 & 10.89 & 10.93 & 10.19 \\
GPTQ & 10.69 & 10.76 & 10.79 & 10.19 \\
AWQ & 10.67 & 10.85 & 10.84 & 10.20 \\
OmniQuant & 10.91 & 10.89 & 10.90 & 10.20 \\
SqueezeLLM & 10.61 & 10.76 & 10.83 & 10.19 \\
SpQR & 10.28 & 10.64 & 10.64 & 10.14 \\
EfficientQAT & 10.52 & 10.73 & 10.76 & 10.20 \\
QuIP & 10.43 & 10.76 & 10.71 & 10.19 \\
QuIP\# & 10.38 & 10.70 & 10.67 & 10.19 \\
PB-LLM & 11.71 & 11.00 & 11.17 & 10.19 \\
bitsandbytes & 10.75 & 10.77 & 10.84 & 10.19 \\
RTN with SmoothQuant & 10.84 & 10.75 & 10.78 & 10.19 \\
\hline
\end{tabularx}
\end{table}

\textbf{What this shows.} With rounded activations the picture changes. Our method is better than Standard for four of the twelve techniques (round-to-nearest, OmniQuant, PB-LLM and round-to-nearest with SmoothQuant) and worse for the other eight; for SpQR the perplexity rises from 10.28 to 10.64. The reason is in the plans: Our method gives more bits to the weights of the protected layers, but its activation plan gives 4 bits to half of the layers, where the Standard plan gives 8 bits to all of them. For the techniques whose 4-bit weights are already good, the loss from the 4-bit activations is larger than the gain from the 8-bit weights. The joint plan gives almost the same result as the separate plans. With 8 bits everywhere, every technique is within 0.04 of the FP16 perplexity, at a llama.cpp file of 1.29~GB.

\textbf{What it means.} A per-layer plan helps where the uniform setting is short of bits (4-bit weights) and does not help where the uniform setting is already generous (8-bit activations) and the plan spends fewer bits than it.

\section{Pruning (Path 0 $\rightarrow$ 2 $\rightarrow$ 4)}
\label{sec:ch5-prune}

Figure~\ref{fig:ch5-prune} and Table~\ref{tab:ch5-prune} report the three pruning techniques. At a pruning level of 30 percent the Standard plan removes 30 percent of every decoder layer, which is 26 percent of all parameters of the model. The Our method plan removes 30 percent of the 11 compressed layers only, which is 13 percent of all parameters. The same-size plan removes the same total as the Standard plan but places it by sensitivity: nothing in the five guarded layers and up to 61 percent in the most robust layer.

\begin{figure}[htbp]
\centering
\includegraphics[width=\textwidth]{Figures/ch5-pruning.pdf}
\caption{Hugging Face validation perplexity after pruning, against the pruning level (log scale; lower is better). The dashed line is the FP16 model.}
\label{fig:ch5-prune}
\end{figure}

\begin{table}[htbp]
\centering
\small
\caption{Hugging Face validation perplexity after pruning (lower is better), measured on the A100. FP16: 10.18.}
\label{tab:ch5-prune}
\renewcommand{\arraystretch}{1.1}
\begin{tabularx}{0.9\textwidth}{Xrrrr}
\hline
\textbf{Technique} & \textbf{Level (\%)} & \textbf{Standard} & \textbf{Our method} & \textbf{Our method,} \\
 & & & & \textbf{same total removed} \\
\hline
Unstructured (Wanda) & 30 & 10.66 & 10.37 & 11.60 \\
 & 40 & 11.54 & 10.66 & 22.88 \\
 & 50 & 14.12 & 11.51 & 59.30 \\
 & 60 & 26.17 & 14.97 & 1{,}233 \\
 & 70 & 94.55 & 28.85 & 13{,}185 \\
\hline
Structured & 30 & 17.47 & 12.00 & 18.65 \\
 & 40 & 26.54 & 13.64 & 47.04 \\
 & 50 & 48.18 & 17.02 & 791 \\
 & 60 & 111 & 25.84 & 2{,}343 \\
 & 70 & 1{,}187 & 67.86 & 2{,}343 \\
\hline
Low-rank & 30 & 163 & 26.84 & 175 \\
 & 40 & 425 & 51.08 & 573 \\
 & 50 & 1{,}168 & 114 & 2{,}009 \\
 & 60 & 6{,}428 & 267 & 8{,}252 \\
 & 70 & 11{,}186 & 647 & 160{,}107 \\
\hline
\end{tabularx}
\end{table}

\textbf{What this shows.} Our method has a lower perplexity than Standard at every level for all three techniques (20 of 20 comparisons). With unstructured pruning at 50 percent the perplexity is 11.51 against 14.12, and at 70 percent 28.85 against 94.55. As in quantization, the comparison is not at equal size, because Our method removes about half as many weights. The same-size plan gives the comparison at equal size, and it is worse than Standard in 19 of the 20 comparisons: at 30 percent unstructured pruning 11.60 against 10.66. Moving the pruning away from the sensitive layers and concentrating it on the robust ones does more damage than spreading it evenly. Of the three techniques, unstructured pruning loses least, structured pruning is clearly worse (17.47 at 30 percent on the Standard plan) and the low-rank factorisation is the worst (163).

\textbf{What the engines show.} Pruned weights are stored as zeros, and the engines confirm that this brings nothing: the llama.cpp file and the vLLM file are 2.20~GB in all 61 rows, the vLLM decode latency is between 2.51 and 2.55~ms per token in all rows (FP16: 2.54), and llama.cpp runs at 412 to 426 tokens per second (FP16: 421). Because the weights of a pruned model are exported at 16 bits without re-rounding, the vLLM perplexity agrees with the Hugging Face perplexity to within 0.05 percent in every row.

\section{KV-cache Compression (Path 0 $\rightarrow$ 3 $\rightarrow$ 4)}
\label{sec:ch5-kv}

Table~\ref{tab:ch5-kv} reports the six KV-cache techniques. The cache memory is the memory that the setting needs for a context of 2{,}048 tokens and a batch of one; the FP16 cache needs 46.1~MB. The Standard rounding techniques use 4 bits for every key and value and keep every token (13.0~MB). The Standard eviction techniques keep 20 percent of the tokens at 16 bits (9.2~MB). Our method applies the Stage~0 KV-cache plan, which both rounds and evicts, and needs 3.2~MB.

\begin{table}[htbp]
\centering
\small
\caption{KV-cache compression, measured in Hugging Face on the A100. Perplexity on the validation set and on the held-out set (lower is better); cache memory for 2{,}048 tokens (smaller is better). FP16: 10.18 / 10.62 with 46.1~MB.}
\label{tab:ch5-kv}
\renewcommand{\arraystretch}{1.15}
\begin{tabularx}{\textwidth}{Xrrrrrr}
\hline
 & \multicolumn{3}{c}{\textbf{Standard}} & \multicolumn{3}{c}{\textbf{Our method}} \\
\textbf{Technique} & Val. & Held-out & MB & Val. & Held-out & MB \\
\hline
Round-to-nearest KV & 10.25 & 10.68 & 13.0 & 11.68 & 12.01 & 3.2 \\
QuaRot KV & 10.33 & 10.78 & 13.0 & 11.79 & 12.12 & 3.2 \\
KVQuant & 10.27 & 10.71 & 13.0 & 11.74 & 12.07 & 3.2 \\
H2O & 11.10 & 11.49 & 9.2 & 11.68 & 12.01 & 3.2 \\
SnapKV & 10.56 & 10.99 & 9.2 & 11.11 & 11.51 & 3.2 \\
InfiniGen & 11.00 & 11.41 & 9.2 & 10.38 & 10.80 & 3.2 \\
\hline
\end{tabularx}
\end{table}

\textbf{What this shows.} For five of the six techniques the Our method row has a higher perplexity than the Standard row, and it uses a cache that is 4.0 times smaller than the Standard rounding cache and 2.9 times smaller than the Standard eviction cache. The exception is InfiniGen, where Our method is better on both counts: 10.38 against 11.00 with about a third of the memory. The three Standard rounding techniques stay close to FP16 (10.25 to 10.33), which shows that 4-bit keys and values cost little on this model; most of the loss in the Our method rows comes from the tokens that the plan evicts. The Our method rows of round-to-nearest KV and H2O are equal because under the plan the two techniques apply the same rounding and keep the same tokens.

\textbf{What it means.} The KV-cache plan trades quality for memory rather than improving both. Whether it is a good trade depends on the technique that carries it out: with InfiniGen the plan gives the lowest perplexity of all KV-cache rows at the smallest cache. The KV-cache techniques exist only in the Hugging Face model, so these results describe quality and cache size, not speed.

\section{Sequential Path and Search}
\label{sec:ch5-search}

In the sequential path one model is quantized, then pruned, then given a KV-cache technique, and the search of Section~\ref{sec:ch4-search} chooses the three techniques and the plan settings. All searches in this section use the Bayesian optimiser and the search ranges of Table~\ref{tab:ch4-space} unless stated. Following Section~\ref{sec:ch1-scope}, the searches are reported as outcomes.

\subsection{Searches under the Two Limits}

Table~\ref{tab:ch5-searches} lists the searches that were run under the requirement of Section~\ref{sec:ch4-requirement}: a validation perplexity of at most 11.5 and a vLLM decode latency of at most 2.5~ms per token.

\begin{table}[htbp]
\centering
\footnotesize
\caption{Searches under the limits of perplexity at most 11.5 and vLLM decode latency at most 2.5~ms per token. HF: Hugging Face. The FP16 model has a perplexity of 10.18 and a vLLM decode latency of 2.54~ms per token.}
\label{tab:ch5-searches}
\renewcommand{\arraystretch}{1.2}
\begin{tabularx}{\textwidth}{>{\raggedright\arraybackslash}p{4.3cm}lrr>{\raggedright\arraybackslash}X}
\hline
\textbf{Search (objectives in order)} & \textbf{Perplexity} & \textbf{Tries} & \textbf{Within} & \textbf{Outcome} \\
 & \textbf{judged in} & & \textbf{both limits} & \\
\hline
Lowest perplexity, then lowest latency, then smallest size & HF & 30 & 2 & Final model: 0.84~GB, 1.98~ms (Table~\ref{tab:ch5-finals}) \\
The same, with weight-only quantization only & HF & 30 & 17 & Final model: 0.90~GB, 2.00~ms (Table~\ref{tab:ch5-finals}) \\
The same as the first, with a perplexity limit of 11.0 & HF & 30 & 0 & Closest try: perplexity 11.50 at 2.0~ms \\
The same as the first, with a latency limit of 2.0~ms & HF & 30 & 1 & In effect the final model of the first search (same technique and calibration, perplexity 11.50) \\
The same as the first, with structured pruning only & HF & 30 & 0 & Closest try: perplexity 11.52 at 1.99~ms, with nothing pruned \\
Smallest size, then lowest perplexity, then lowest latency & HF & 30 & 0 & Fast tries at perplexity 11.75 or higher; the tries below 11.5 decode in about 42~ms \\
The same as the first, repeated with 40 tries & HF & 40 & 0 & Closest fast try: perplexity 11.75 at 2.14~ms \\
Lowest perplexity, then smallest size & vLLM & 40 & 0 & Figure~\ref{fig:ch5-tries}, left \\
Fastest, then smallest size & vLLM & 40 & 0 & Figure~\ref{fig:ch5-tries}, right \\
\hline
\end{tabularx}
\end{table}

Two searches gave a final model. Their settings and measurements are in Table~\ref{tab:ch5-finals}, next to the FP16 model and to two models from earlier searches that are described in Section~\ref{sec:ch5-smallest}.

\begin{table}[htbp]
\centering
\footnotesize
\caption{The final models, measured in vLLM on the A100, and their accuracy. Perplexity is given on the validation set and on the held-out set. Lower is better for size, perplexity, latency and memory; higher is better for accuracy.}
\label{tab:ch5-finals}
\renewcommand{\arraystretch}{1.2}
\begin{tabularx}{\textwidth}{>{\raggedright\arraybackslash}p{2.25cm}>{\raggedright\arraybackslash}p{1.15cm}>{\raggedright\arraybackslash}X>{\raggedright\arraybackslash}X>{\raggedright\arraybackslash}X>{\raggedright\arraybackslash}X}
\hline
 & \textbf{FP16} & \textbf{Lowest-perplexity search} & \textbf{Weight-only search} & \textbf{Smallest-size search} & \textbf{Smallest-size search, narrower ranges} \\
\hline
Quantization & none & OmniQuant, weights only; threshold 0.857; per channel; C4, 32 windows & Round-to-nearest, weights only; threshold 0.800; group 64; C4, 128 windows & GPTQ with round-to-nearest activations; threshold 0.759; group 128; WikiText-2, 64 windows & SqueezeLLM with SmoothQuant (strength 0.638); threshold 0.876; group 32; Pile, 128 windows \\
Pruning & none & Unstructured, ratio 0.082 & Unstructured, ratio 0 (nothing removed) & Unstructured, ratio 0 (nothing removed) & Unstructured, ratio 0.262 \\
KV cache & none & InfiniGen & SnapKV & InfiniGen & InfiniGen \\
\hline
vLLM file (GB) & 2.200 & 0.838 & 0.898 & 0.894 & 0.873 \\
HF perplexity & 10.18 / 10.62 & 11.50 / 11.91 & 11.12 / 11.52 & 11.23 / 11.62 & 11.44 / 11.87 \\
vLLM perplexity & 10.18 / 10.62 & 13.00 / 13.41 & 11.10 / 11.54 & 12.09 / 12.52 & 11.46 / 11.89 \\
vLLM decode (ms per token) & 2.54 & 1.98 & 2.00 & 34.42 & 42.45 \\
vLLM peak memory (GB) & 3.77 & 2.40 & 2.45 & 2.46 & 2.44 \\
\hline
ARC-Easy & 0.548 & 0.443 & 0.325 & 0.402 & 0.415 \\
PIQA & 0.745 & 0.691 & 0.630 & 0.676 & 0.676 \\
HellaSwag & 0.604 & 0.566 & 0.516 & 0.551 & 0.555 \\
WinoGrande & 0.603 & 0.523 & 0.510 & 0.521 & 0.509 \\
Mean accuracy & 0.625 & 0.556 & 0.495 & 0.538 & 0.539 \\
\hline
\end{tabularx}
\end{table}

\textbf{The lowest-perplexity search.} Two of its 30 tries were within both limits. The final model uses OmniQuant on the weights only (sensitivity threshold 0.857, per-channel groups, calibrated on 32 windows of C4), unstructured pruning with a ratio of 0.082 and InfiniGen for the KV cache. Its vLLM file is 0.838~GB, 62 percent smaller than the FP16 file, and it decodes in 1.98~ms per token, 22 percent faster than FP16. Its perplexity is 11.50 in Hugging Face but 13.00 in vLLM, and its mean accuracy is 0.556 against 0.625 for FP16, a relative drop of 11.1 percent.

\textbf{The weight-only search.} When Stage~1 was restricted to weight-only techniques, 17 of the 30 tries were within both limits. The final model uses round-to-nearest on the weights only (threshold 0.800, which gives 8 bits to the five layers 0, 1, 2, 7 and 21; groups of 64 weights; calibrated on 128 windows of C4), unstructured pruning with a ratio of 0, so that nothing is removed, and SnapKV for the KV cache. Its vLLM file is 0.898~GB (59 percent smaller than FP16) and it decodes in 2.00~ms per token (21 percent faster). Its perplexity is 11.12 in Hugging Face and 11.10 in vLLM, so this model is within both limits also when the perplexity is judged in vLLM. The two values are close by coincidence: the Hugging Face value includes the KV-cache step and the vLLM value does not (Section~\ref{sec:ch5-checks}). Its mean accuracy is 0.495, a relative drop of 20.8 percent, the largest of the final models, although its perplexity is the lowest of them.

\textbf{Tighter limits and other variants.} Lowering the perplexity limit to 11.0 left no try within the limits, and lowering the latency limit to 2.0~ms left one, which is in effect the final model of the first search. With structured pruning as the only Stage~2 technique, the closest try pruned nothing and reached 11.52. With size as the first objective, no try was within the limits. A repeat of the first search with 40 tries, made with a later version of the code, also had no try within both limits; its closest fast try had a perplexity of 11.75.

\textbf{Perplexity judged in vLLM.} Because the perplexity of an exported model can differ from its Hugging Face perplexity, the last two searches checked the perplexity limit on the model as loaded in vLLM. Figure~\ref{fig:ch5-tries} shows all 40 tries of each. The tries fall into two groups. The 30 weight-only tries, all of which used OmniQuant, decode in 1.97 to 2.14~ms per token, faster than FP16, with a vLLM file of 0.82 to 0.86~GB, but their vLLM perplexity is 12.97 or higher. The 10 tries with rounded activations decode in 33 to 43~ms per token; one of them reaches a vLLM perplexity of 10.98. No try is within both limits. Round-to-nearest as a weight-only technique, which gave the final model of the weight-only search, was not proposed in these tries.

\begin{figure}[htbp]
\centering
\includegraphics[width=\textwidth]{Figures/ch5-search-tries.pdf}
\caption{All 40 tries of the two searches that judge perplexity in vLLM. Lower and further left is better. The dashed lines are the two limits; a try within both limits would lie in the lower left corner.}
\label{fig:ch5-tries}
\end{figure}

\subsection{Searches for the Smallest Model}
\label{sec:ch5-smallest}

Before the latency limit was introduced, searches were run with the perplexity limit only (Hugging Face perplexity at most 11.5) and with the predicted weight memory as the first objective. The first of the two final models in Table~\ref{tab:ch5-finals} comes from a search over the full ranges of Table~\ref{tab:ch4-space} and uses GPTQ with round-to-nearest activations; the second comes from a search over the narrower ranges, in which a layer could also be given 3 bits, and uses SqueezeLLM with SmoothQuant; no layer of this final model has 3 bits. Both are small, 0.894 and 0.873~GB in vLLM, and both are within the perplexity limit in Hugging Face. Both keep the activations at low bit widths, and in vLLM they decode in 34.4 and 42.5~ms per token, 13.5 and 16.7 times slower than FP16. A small file is therefore not a fast model, and this is what led to the latency limit.

These two searches were also run with successive halving on the same search space and seed and with the same budget of 30 full evaluations (Section~\ref{sec:ch4-optimiser}), which gives the comparison of Table~\ref{tab:ch5-optimisers}. In both, the Bayesian optimiser found a smaller model within the limit, with fewer models built and in less time.

\begin{table}[htbp]
\centering
\small
\caption{Bayesian optimisation against successive halving in the two searches for the smallest model with a Hugging Face perplexity of at most 11.5.}
\label{tab:ch5-optimisers}
\renewcommand{\arraystretch}{1.15}
\begin{tabularx}{\textwidth}{llrrr>{\raggedright\arraybackslash}X}
\hline
\textbf{Search} & \textbf{Optimiser} & \textbf{Models} & \textbf{Build time} & \textbf{Within} & \textbf{Smallest model within the limit} \\
 & & \textbf{built} & & \textbf{the limit} & \textbf{(predicted size, perplexity)} \\
\hline
First & Bayesian & 30 & 245 min & 4 & 0.924~GB at 11.39 \\
 & Successive halving & 129 & 495 min & 10 & 0.939~GB at 11.44 \\
Second & Bayesian & 30 & about 190 min & 8 & 0.868~GB at 11.42 \\
 & Successive halving & 130 & about 500 min & 0 & none (closest: 0.916~GB at 11.66) \\
\hline
\end{tabularx}
\end{table}

\section{Checks of the Score, the Plan and the Stages}
\label{sec:ch5-checks}

Three further experiments test the parts of the framework one at a time: whether the sensitivity score predicts the damage of compressing a layer, whether the layers chosen by the score are a better choice than other layers at the same size, and what each stage contributes to the two final models found under both limits. A fourth measures a compressed KV cache in the engines. They were run on the A100; the latencies of the stage ablation come from a separate session, so they are compared only with each other.

\subsection{Does the Score Predict Compression Damage?}

Each of the 22 layers was changed alone in four ways: its weights were rounded to 8 bits and to 4 bits (round-to-nearest, groups of 128), and 25 percent of its weights were pruned by magnitude and by Wanda. The rise in perplexity was compared with the rise when the layer is bypassed, by Spearman rank correlation. On the calibration text, which the score uses, the correlation with bypassing is 0.04 for 8 bits, 0.18 for 4 bits, 0.15 for magnitude pruning and 0.41 for Wanda; on the validation set it is 0.46, 0.51, 0.54 and 0.42. Rounding one layer to 8 bits changes the perplexity by less than 0.002 in every layer, and the largest rise from 4 bits in one layer is 0.13. Layer~0, by far the most sensitive to bypassing, is hardly affected by 4-bit rounding alone (0.004 on the validation set). The score of Equation~\ref{eq:ch4-score} is therefore only a weak to moderate predictor of the damage of compressing a single layer.

\subsection{Per-layer Plan at Equal Size and at Equal Perplexity}

Table~\ref{tab:ch5-equal} compares bit plans of the same size for weight-only quantization alone, for four of the weight techniques. With 8 bits in 11 layers and 4 bits in the rest, the plan chosen by the score has a lower perplexity than the inverse plan and than each of five random plans, for all four techniques. With 8 bits in 5 layers it is better than the inverse plan for all four techniques and better than every random plan for round-to-nearest and OmniQuant; for GPTQ and SqueezeLLM one of the five random plans is better (10.46 against 10.49 and 10.47 against 10.49). So the plan is a better choice of layers than the alternatives at the same size. The plan is not a gain over 4 bits in every layer at equal size: giving 8 bits to the five most sensitive layers and 3 bits to the rest has about the size of uniform 4 bits but a much higher perplexity (12.55 against 10.97 for round-to-nearest), and every plan with 8 bits in some layers and 4 bits in the rest is larger than uniform 4 bits (0.84~GB or more against 0.78~GB). The plan trades size for quality.

\begin{table}[htbp]
\centering
\footnotesize
\caption{Weight-only quantization in Hugging Face with bit plans of equal size: validation perplexity (FP16: 10.18; lower is better). Plan: 8 bits in the most sensitive layers; inverse: 8 bits in the least sensitive layers; random: range of five random choices of layers. Sizes are the predicted weight memory for round-to-nearest, GPTQ and OmniQuant; for SqueezeLLM they are 0.79, 0.84, 0.94 and 1.12~GB.}
\label{tab:ch5-equal}
\scriptsize
\setlength{\tabcolsep}{3pt}
\renewcommand{\arraystretch}{1.15}
\begin{tabularx}{\textwidth}{>{\raggedright\arraybackslash}p{1.7cm}>{\raggedright\arraybackslash}p{1.1cm}>{\raggedright\arraybackslash}p{1.5cm}>{\raggedright\arraybackslash}p{0.75cm}>{\raggedright\arraybackslash}p{0.9cm}>{\raggedright\arraybackslash}X>{\raggedright\arraybackslash}p{0.75cm}>{\raggedright\arraybackslash}p{0.9cm}>{\raggedright\arraybackslash}X}
\hline
 & \textbf{4 bits in every layer} & \textbf{8 bits in 5 layers, 3 bits in the rest} & \multicolumn{3}{>{\raggedright\arraybackslash}p{4.2cm}}{\textbf{8 bits in 5 layers, 4 bits in the rest (0.89~GB)}} & \multicolumn{3}{>{\raggedright\arraybackslash}p{4.2cm}}{\textbf{8 bits in 11 layers, 4 bits in the rest (1.02~GB)}} \\
\textbf{Technique} & (0.78~GB) & (0.79~GB) & Plan & Inverse & Random & Plan & Inverse & Random \\
\hline
Round-to-nearest & 10.97 & 12.55 & 10.61 & 10.84 & \mbox{10.68--10.80} & 10.41 & 10.71 & \mbox{10.49--10.63} \\
GPTQ & 10.63 & 11.66 & 10.49 & 10.58 & \mbox{10.46--10.57} & 10.32 & 10.51 & \mbox{10.36--10.41} \\
OmniQuant & 10.89 & 12.37 & 10.56 & 10.79 & \mbox{10.62--10.75} & 10.39 & 10.66 & \mbox{10.45--10.58} \\
SqueezeLLM & 10.60 & 11.29 & 10.49 & 10.52 & \mbox{10.47--10.51} & 10.33 & 10.43 & \mbox{10.36--10.45} \\
\hline
\end{tabularx}
\end{table}

Table~\ref{tab:ch5-equalppl} turns the comparison around: for a target perplexity, it gives the smallest plan of each kind, among the plans of Table~\ref{tab:ch5-equal}, that reaches the target. The targets are the perplexities of the two score plans. To reach the perplexity of the score plan with 8 bits in 5 layers (0.89~GB), a random plan needs 11 layers at 8 bits (1.02~GB) for round-to-nearest and OmniQuant, and no inverse plan reaches it; for GPTQ and SqueezeLLM one random plan of the same size reaches it, and for SqueezeLLM the inverse plan needs 11 layers (1.12~GB). No inverse or random plan reaches the perplexity of the score plan with 8 bits in 11 layers, for any of the four techniques. At equal perplexity the plan chosen by the score is therefore as small as or smaller than the other choices of layers.

\begin{table}[htbp]
\centering
\footnotesize
\caption{Weight-only quantization in Hugging Face at equal perplexity: predicted weight memory in GB of the smallest plan of each kind that reaches the target validation perplexity (smaller is better). The targets are the perplexities of the score plans with 8 bits in 5 and in 11 layers. Inverse and random plans were built with 5 and with 11 layers at 8 bits; none: no such plan reaches the target.}
\label{tab:ch5-equalppl}
\setlength{\tabcolsep}{3pt}
\renewcommand{\arraystretch}{1.15}
\begin{tabularx}{\textwidth}{>{\raggedright\arraybackslash}p{2.4cm}>{\raggedright\arraybackslash}X>{\raggedright\arraybackslash}X>{\raggedright\arraybackslash}X>{\raggedright\arraybackslash}X>{\raggedright\arraybackslash}X>{\raggedright\arraybackslash}X>{\raggedright\arraybackslash}X>{\raggedright\arraybackslash}X}
\hline
 & \multicolumn{4}{>{\raggedright\arraybackslash}p{5.2cm}}{\textbf{Target: score plan, 5 layers at 8 bits}} & \multicolumn{4}{>{\raggedright\arraybackslash}p{5.2cm}}{\textbf{Target: score plan, 11 layers at 8 bits}} \\
\textbf{Technique} & Target & Plan & Inverse & Random & Target & Plan & Inverse & Random \\
\hline
Round-to-nearest & 10.61 & 0.89 & none & 1.02 & 10.41 & 1.02 & none & none \\
GPTQ & 10.49 & 0.89 & none & 0.89 & 10.32 & 1.02 & none & none \\
OmniQuant & 10.56 & 0.89 & none & 1.02 & 10.39 & 1.02 & none & none \\
SqueezeLLM & 10.49 & 0.94 & 1.12 & 0.94 & 10.33 & 1.12 & none & none \\
\hline
\end{tabularx}
\end{table}

\subsection{Contribution of Each Stage}

Table~\ref{tab:ch5-ablation} repeats the two final models found under both limits with stages left out. The first uses round-to-nearest on the weights only (threshold 0.800, groups of 64) and SnapKV; the second uses OmniQuant on the weights only (threshold 0.857, per-channel groups) and InfiniGen. In the rows with pruning, both are pruned with unstructured pruning at a ratio of 0.082 in the layers below their threshold, which removes 5.6 percent of the weights.

\begin{table}[htbp]
\centering
\scriptsize
\caption{Stage ablation with the settings of the two final models found under both limits. Q: quantization; P: unstructured pruning at a ratio of 0.082; KV: KV-cache technique. Perplexity on the validation set; accuracy measured in Hugging Face. vLLM was measured on the model before the KV-cache step, so rows that differ only in KV share their vLLM values. Lower is better except for accuracy.}
\label{tab:ch5-ablation}
\renewcommand{\arraystretch}{1.1}
\begin{tabularx}{\textwidth}{>{\raggedright\arraybackslash}p{1.1cm}>{\raggedright\arraybackslash}p{1.5cm}>{\raggedright\arraybackslash}X>{\raggedright\arraybackslash}X>{\raggedright\arraybackslash}X>{\raggedright\arraybackslash}X>{\raggedright\arraybackslash}X>{\raggedright\arraybackslash}X>{\raggedright\arraybackslash}X>{\raggedright\arraybackslash}X>{\raggedright\arraybackslash}X}
\hline
\textbf{Row} & \textbf{Stages} & HF perplexity & ARC-Easy & PIQA & Hella\-Swag & Wino\-Grande & Mean accuracy & vLLM file (GB) & Peak memory (GB) & Decode (ms per token) \\
\hline
\multicolumn{11}{l}{\textbf{Round-to-nearest (threshold 0.800), SnapKV}} \\
Baseline & none & 10.18 & 0.548 & 0.745 & 0.604 & 0.603 & 0.625 & 2.200 & 3.77 & 2.80 \\
E1 & Q & 10.44 & 0.548 & 0.739 & 0.598 & 0.609 & 0.624 & 0.898 & 2.45 & 1.98 \\
E2 & P & 10.18 & 0.548 & 0.746 & 0.606 & 0.605 & 0.626 & 2.200 & 3.77 & 2.82 \\
E3 & KV & 10.82 & 0.332 & 0.632 & 0.521 & 0.504 & 0.497 & 2.200 & 3.77 & 2.80 \\
E4 & Q + P & 10.44 & 0.548 & 0.739 & 0.598 & 0.609 & 0.624 & 0.898 & 2.45 & 2.03 \\
E5 & Q + KV & 11.11 & 0.325 & 0.630 & 0.516 & 0.510 & 0.495 & 0.898 & 2.45 & 1.98 \\
E6 & Q + P + KV & 11.11 & 0.325 & 0.630 & 0.516 & 0.510 & 0.495 & 0.898 & 2.45 & 2.03 \\
\hline
\multicolumn{11}{l}{\textbf{OmniQuant (threshold 0.857), InfiniGen}} \\
Baseline & none & 10.18 & 0.548 & 0.745 & 0.604 & 0.603 & 0.625 & 2.200 & 3.77 & 2.81 \\
E1 & Q & 11.25 & 0.529 & 0.742 & 0.586 & 0.594 & 0.613 & 0.838 & 2.40 & 2.03 \\
E2 & P & 10.18 & 0.548 & 0.745 & 0.606 & 0.605 & 0.626 & 2.200 & 3.77 & 2.81 \\
E3 & KV & 10.36 & 0.464 & 0.694 & 0.584 & 0.505 & 0.562 & 2.200 & 3.77 & 2.81 \\
E4 & Q + P & 11.25 & 0.529 & 0.742 & 0.586 & 0.594 & 0.613 & 0.838 & 2.40 & 2.09 \\
E5 & Q + KV & 11.47 & 0.432 & 0.689 & 0.565 & 0.526 & 0.553 & 0.838 & 2.40 & 2.03 \\
E6 & Q + P + KV & 11.47 & 0.432 & 0.689 & 0.565 & 0.526 & 0.553 & 0.838 & 2.40 & 2.09 \\
\hline
\end{tabularx}
\end{table}

Quantization gives all of the gain in size, memory and speed: for round-to-nearest the vLLM file falls from 2.200 to 0.898~GB, the peak memory from 3.77 to 2.45~GB and the decode latency from 2.80 to 1.98~ms per token, for a Hugging Face perplexity of 10.44 and an unchanged mean accuracy (0.624 against 0.625). Pruning alone at this ratio leaves perplexity, file size and latency unchanged, and after 4-bit rounding it changes nothing at all: rows E4 and E6 equal rows E1 and E5 in every quality number. The KV-cache step causes the loss of accuracy: SnapKV alone raises the perplexity to 10.82 and lowers the mean accuracy from 0.625 to 0.497, and InfiniGen alone gives 10.36 and 0.562. ARC-Easy loses the most: SnapKV alone lowers it from 0.548 to 0.332 and the other three tasks by 0.08 to 0.11, and InfiniGen alone lowers it to 0.464. The loss of 11 to 21 percent in Table~\ref{tab:ch5-finals} therefore comes from the KV-cache technique and not from quantization. Because that step does not run in vLLM, the model that vLLM serves is row E1.

The ablation also separates the two effects in the comparison of the engines. The Hugging Face values of Table~\ref{tab:ch5-finals} include the KV-cache step and the vLLM values do not. For the quantized model alone (row E1) the perplexity is 10.44 in Hugging Face and 11.10 in vLLM for round-to-nearest, and 11.25 and 12.96 for OmniQuant: the export raises it in both cases. The agreement of 11.12 and 11.10 for the round-to-nearest final model is therefore a coincidence of two different effects.

The ratio of 0.082 is the value chosen by the search, and it is too small to show what pruning does; larger ratios are covered by Section~\ref{sec:ch5-prune}. In the searches, 14 tries combined OmniQuant on the weights only with unstructured pruning at ratios of 0.20 to 0.29 and thresholds of 0.857 to 0.891. Their vLLM perplexity is 13.34 to 13.88, against 12.96 for OmniQuant alone in row E1, with files of 0.815 to 0.838~GB and 1.99 to 2.14~ms per token. These tries also differ from row E1 in the threshold and in the other settings, so the difference cannot be assigned to pruning alone.

\subsection{KV-cache Compression in the Engines}

To measure a compressed cache in a real engine, the KV-cache technique of the pipeline was replaced by the method that each engine provides: an 8-bit cache in vLLM and a 4-bit cache in llama.cpp. The two final models, built without their KV-cache step, were measured with the 16-bit and with the compressed cache. In vLLM the 8-bit cache holds twice as many tokens in the same 1~GiB (95{,}312 against 47{,}648). For the round-to-nearest model the perplexity rises from 11.10 to 11.14, the decode latency from 1.97 to 2.50~ms per token and the peak memory from 2.45 to 3.17~GB; for the OmniQuant model the values are 12.96 to 13.02, 2.01 to 2.48~ms and 2.40 to 3.12~GB. In llama.cpp the 4-bit cache is 3.6 times smaller (0.013 against 0.046~GB); the perplexity of the round-to-nearest model rises from 8.40 to 8.53 and its decode speed falls from 507 to 463 tokens per second. At a batch of one and short texts a compressed cache saves little memory and costs time; its gain, more text in the same memory, was not measured with longer contexts or larger batches.

\section{Comparison with Published Work on TinyLlama}
\label{sec:ch5-published}

The final models are compared with published studies that compress the same model, TinyLlama, and only results measured in the same engine are compared. Because the studies use different evaluation settings, the comparison is made on the change against each study's own uncompressed model. Table~\ref{tab:ch5-published} covers Hugging Face, where two pruning studies report perplexity and accuracy \cite{ilin2026qkwanda, ilin2025thanos}. The two final models found under both limits raise the Hugging Face perplexity by 9.2 and 13.0 percent while the vLLM file becomes 59 and 62 percent smaller; pruning half of the weights raises it by 38.6 to 44.4 percent, and pruning 30 percent by 6.9 percent. The pruning studies lose less accuracy than the two final models, on a different set of tasks, and their uncompressed perplexity (7.77 and 7.97 against 10.18) shows that the evaluation settings differ.

\begin{table}[htbp]
\centering
\small
\caption{Hugging Face results on TinyLlama: change against each study's own uncompressed model. Values for other studies are calculated from the numbers reported in the cited papers (perplexity on WikiText-2). The perplexity of this work is the Hugging Face value on the validation set. Lower is better for the change in perplexity; a smaller loss of accuracy is better.}
\label{tab:ch5-published}
\renewcommand{\arraystretch}{1.15}
\begin{tabularx}{\textwidth}{>{\raggedright\arraybackslash}X>{\raggedright\arraybackslash}p{2.6cm}>{\raggedright\arraybackslash}p{2.2cm}>{\raggedright\arraybackslash}p{1.7cm}>{\raggedright\arraybackslash}p{1.6cm}}
\hline
\textbf{Study} & \textbf{Methods used} & \textbf{Uncompressed perplexity} & \textbf{Perplexity} & \textbf{Mean accuracy} \\
\hline
This work: round-to-nearest, weights only (threshold 0.800, group 64); no pruning; SnapKV & Quantization, pruning, KV cache & 10.18 & $+9.2\%$ & $-21\%$ \\
This work: OmniQuant, weights only (threshold 0.857, per channel); pruning 0.082; InfiniGen & Quantization, pruning, KV cache & 10.18 & $+13.0\%$ & $-11\%$ \\
QK-Wanda \cite{ilin2026qkwanda}, 30\% pruned & Unstructured pruning & 7.77 & $+6.9\%$ & unchanged \\
QK-Wanda \cite{ilin2026qkwanda}, 50\% pruned & Unstructured pruning & 7.77 & $+44.4\%$ & $-2.7\%$ \\
Thanos \cite{ilin2025thanos}, 50\% pruned & Unstructured pruning & 7.97 & $+38.6\%$ & $-7.3\%$ \\
\hline
\end{tabularx}
\end{table}

For the other engines there is no result on both sides. In llama.cpp, Das et al.\ prune 7 percent of TinyLlama, recover with LoRA and quantize to 4-bit GGUF; the FP16 file becomes 4.9 percent smaller and the SQuAD F1 score falls by 52 percent \cite{das2026largemodels}; the final models of this thesis were not measured in llama.cpp. On a phone processor, MobileQuant quantizes weights and activations to 8 bits with 4.7 percent higher perplexity and 15 percent lower decode time than with 16-bit activations \cite{tan2024mobilequant}. EdgeProfiler estimates with an analytical model that 4-bit TinyLlama is 71 percent smaller and 2.45 times faster \cite{pinnock2025edgeprofiler}, and Belhaouari and Kraidia report 45 percent compression and a 1.68 times speedup from pruning with retraining \cite{belhaouari2025smartpruning}. No published study was found that measures a compressed TinyLlama in vLLM, or that combines quantization, pruning and KV-cache compression on it.

\section{Discussion}
\label{sec:ch5-discussion}

This section answers the five research questions of Chapter~1 with the evidence above.

\textbf{RQ1: how uneven is the sensitivity, and does the score predict the damage?} The sensitivity is very uneven: five of the 22 layers have a score above 10 and two of them above 300, while the other 17 are between 1.4 and 5.7 (Figure~\ref{fig:ch5-sensitivity}). The score is a useful guide for protection. Giving the protected layers 8 bits lowered the perplexity of all eleven weight techniques, and leaving them unpruned lowered the perplexity of all three pruning techniques at every level. The score is not a good guide for how much more to compress the other layers: the same-size pruning plan, which moves the pruning from the sensitive layers onto the robust ones, was worse than uniform pruning in 19 of 20 comparisons. A layer that can be bypassed with little damage cannot, on this evidence, absorb a much larger share of the pruning. Measured layer by layer, the score is only a weak to moderate predictor of the damage of compressing a single layer (Spearman correlations of 0.04 to 0.54, Section~\ref{sec:ch5-checks}).

\textbf{RQ2: does the per-layer plan improve the trade-off?} For weight-only quantization it improves quality at a cost in size and speed: every technique moves closer to FP16, the file grows from 0.81 to 1.05~GB and llama.cpp slows from about 650 to about 550 tokens per second. This is a different point on the trade-off, between the uniform 4-bit and the uniform 8-bit model, and not a gain at equal size. At equal size the layers chosen by the score are a better choice than the least sensitive or random layers, and at equal perplexity the plan is as small as or smaller than these choices; a plan with 8 and 3 bits at the size of the uniform 4-bit model is much worse (Section~\ref{sec:ch5-checks}). For rounded activations the plan was worse than the uniform 8-bit setting for eight of twelve techniques, because it spends fewer activation bits. For the KV cache the plan gives a cache 3 to 4 times smaller at a higher perplexity for five of six techniques, and is better on both counts with InfiniGen. The plan is most useful with simple techniques: it brings round-to-nearest (10.41) close to what the best technique reaches on the uniform plan (SpQR, 10.27).

\textbf{RQ3: quantization only, pruning only, or both?} In the engines studied here quantization is the only stage that reduces the file size and the latency. Pruning costs perplexity and returns nothing in size or speed, because the removed weights are stored as zeros (Section~\ref{sec:ch5-prune}). The searches agree: the final models within both limits prune nothing or very little (ratios of 0 and 0.082), and with structured pruning only, the best try pruned nothing. On this evidence a quantization-only pipeline is preferable for these engines, and pruning becomes worth its quality cost only where the engine or the hardware can exploit the zeros. A KV-cache technique reduces cache memory and does not change the exported model.

\textbf{RQ4: how do the simulated results carry over to the engines?} Three differences were observed. First, quality carries over exactly when the weights are exported without re-rounding (the pruned models agree to within 0.05 percent between Hugging Face and vLLM), but not when the export re-rounds them: the OmniQuant final model moves from 11.50 to 13.00 (from 11.25 to 12.96 without its KV-cache step) and the GPTQ final model from 11.23 to 12.09, while the round-to-nearest final model has 11.12 and 11.10, an agreement that is a coincidence of two effects (without the KV-cache step it moves from 10.44 to 11.10, Section~\ref{sec:ch5-checks}), and one try with rounded activations is better in vLLM (10.98) than in Hugging Face (12.06). Second, speed cannot be read from the simulation: in Hugging Face the weight-only models decode in about the same time as FP16, while in the engines they are faster (about 2.0~ms against 2.54~ms in vLLM, and 1.7 to 2.0 times the FP16 tokens per second in llama.cpp). Third, the kind of quantization decides the kernel: weight-only models are faster than FP16 in vLLM, and models with rounded activations are 13 to 17 times slower. A model should therefore be selected on measurements made in the engine that will serve it.

\textbf{RQ5: can the search find configurations that meet a deployment requirement?} In part. With the perplexity judged in Hugging Face, two searches found models within both limits, and both final models are about 60 percent smaller and about 20 percent faster than FP16 in vLLM. The final model of the weight-only search is within both limits in vLLM as well, and it is the main result of the search: round-to-nearest on the weights only (threshold 0.800, groups of 64), no pruning and SnapKV, 59 percent smaller and 21 percent faster than FP16 in vLLM. With the perplexity judged in vLLM, 40 tries of each of two searches found fast models and low-perplexity models but none that is both. The search also showed what the requirement leaves out: perplexity did not rank the final models by accuracy, since the model with the lowest perplexity lost the most accuracy (20.8 percent) and every final model lost at least 11 percent. The stage ablation shows that this loss comes from the KV-cache step and not from quantization (Section~\ref{sec:ch5-checks}). A requirement that is meant to protect task accuracy needs an accuracy measurement among its limits.

\textbf{Relation to the five gaps.} The single-stage experiments compare 4-bit, plan-based and 8-bit precision under identical conditions (gap one). The pruning and search results give a rule for the pipeline in these engines: quantize, and prune only if the engine can use it (gaps two and three). Every result is reported per engine, and the differences between the simulation and the engines are measured (gap four). The plan and the search choose the bit widths, the pruning ratio and the techniques from a stated requirement (gap five).

\section{Limitations of the Evidence}
\label{sec:ch5-limits}

The results hold for one model, TinyLlama, and one seed; the variation between seeds was not measured, so small differences between rows should not be over-read. Most techniques are simplified re-implementations (Section~\ref{sec:ch1-scope}), and their numbers are not the numbers of the original methods. Standard and Our method differ in size in the quantization and pruning experiments; the comparisons at equal size are the same-size pruning plan and, for four weight techniques in Hugging Face, the plans of Table~\ref{tab:ch5-equal}. The cache memory of the KV-cache techniques is computed from the plan and not measured in an engine, and the KV-cache techniques could not be exported. Latency was measured at a batch of one with a prompt of 128 tokens and 64 generated tokens, on two GPUs. Accuracy was measured for the final models and for the stage ablation only. No TensorRT-LLM measurement is reported. In the searches, the number of tries is small compared with the size of the search space, and a search that found no try within the limits does not show that no such configuration exists.
\clearpage
\chapter{Conclusions and Future Work}
\label{Chapter6}

\section{Summary of the Work}
\label{sec:ch6-summary}

This thesis built a framework that compresses a decoder-only LLM in stages and measures the result in the engines that would deploy it. Stage~0 measures each decoder layer by the rise in calibration perplexity when the layer is bypassed and writes per-layer plans for the weights, the activations, pruning and the KV cache. Stages~1 to 3 apply published quantization, pruning and KV-cache techniques that follow those plans, and each technique is also run on a uniform plan with the same code, so that Standard and Our method differ only in the plan. Stage~4 measures the model in Hugging Face Transformers and, after export, in llama.cpp and vLLM. A constrained Bayesian search chooses one technique per stage and the plan settings under limits on perplexity and decode latency. The six objectives of Chapter~1 were carried out on TinyLlama.

\section{Conclusions}
\label{sec:ch6-conclusions}

The main findings, each for TinyLlama under the conditions of Chapter~\ref{Chapter4}, are the following.

\begin{enumerate}
  \item \textbf{Sensitivity is concentrated in a few layers and is cheap to measure.} Five of the 22 layers account for almost all of the measured sensitivity, with layers 0 and 2 far ahead of the rest. The profile takes about two minutes on one GPU and is reused by every later run.
  \item \textbf{Protecting the sensitive layers improves quality at a cost in size.} With 8 bits in the 11 protected layers, all eleven weight techniques have a lower perplexity than at 4 bits in every layer, in Hugging Face and in llama.cpp, for a file of 1.05~GB instead of 0.81~GB. With the protected layers left unpruned, all three pruning techniques have a lower perplexity at every level. The plan narrows the difference between techniques: round-to-nearest on the plan (10.41) comes close to the best technique on the uniform plan (10.27).
  \item \textbf{The plan does not help where it spends fewer bits than the uniform setting, or at equal size in pruning.} The activation plan (6 bits on average) was worse than uniform 8-bit activations for eight of twelve techniques. The pruning plan that removes the same total as uniform pruning was worse than uniform pruning in 19 of 20 comparisons. The KV-cache plan gives a cache 3 to 4 times smaller at a higher perplexity for five of six techniques. For weight-only quantization at equal size, the layers chosen by the score gave a lower perplexity than the least sensitive layers in all eight comparisons and than every random choice of layers in six of eight, but a plan with 8 and 3 bits at the size of 4 bits in every layer was much worse.
  \item \textbf{In the engines studied, only quantization reduces size and latency.} A 4-bit file is 0.81~GB against 2.20~GB and runs at twice the tokens per second of FP16 in llama.cpp. Pruned models keep the FP16 file size and the FP16 latency in llama.cpp and vLLM, because the removed weights are stored as zeros.
  \item \textbf{Simulated results do not carry over unchanged.} The export can change the perplexity (for the OmniQuant model without its KV-cache step, 11.25 in Hugging Face and 12.96 in vLLM; for round-to-nearest, 10.44 and 11.10), and the kind of quantization decides the speed in vLLM: weight-only models decode in about 2.0~ms per token against 2.54~ms for FP16, and models with rounded activations in 33 to 43~ms.
  \item \textbf{The search finds deployable models when its limits are measured where the model will run.} Two searches gave final models that are about 60 percent smaller and about 20 percent faster than FP16 in vLLM. The main result uses round-to-nearest on the weights only (threshold 0.800, groups of 64, 128 windows of C4), no pruning and SnapKV: 0.898~GB, 2.00~ms per token, perplexity 11.12 in Hugging Face and 11.10 in vLLM, mean accuracy 0.495; it is within both limits in vLLM. The second uses OmniQuant on the weights only (threshold 0.857, per-channel groups, 32 windows of C4), unstructured pruning with a ratio of 0.082 and InfiniGen: 0.838~GB, 1.98~ms per token, perplexity 11.50 in Hugging Face and 13.00 in vLLM, mean accuracy 0.556, so it exceeds the perplexity limit in vLLM. The FP16 model has 2.200~GB, 2.54~ms per token, a perplexity of 10.18 and a mean accuracy of 0.625. When the perplexity limit was checked in vLLM, two searches of 40 tries each found no model within both limits. Against published work on TinyLlama in Hugging Face, the two final models raise perplexity by 9.2 and 13.0 percent, less than pruning half of the weights (38.6 to 44.4 percent), and lose more accuracy than the pruning studies, which use a different set of tasks (Section~\ref{sec:ch5-published}).
  \item \textbf{Perplexity is not enough to protect accuracy.} Every final model lost between 11 and 21 percent of the mean accuracy on the four tasks, and the final model with the lowest perplexity lost the most. The stage ablation shows that the loss comes from the KV-cache step: with quantization alone the mean accuracy is 0.624 for round-to-nearest and 0.613 for OmniQuant against 0.625 for FP16, while SnapKV alone lowers it to 0.497 and InfiniGen alone to 0.562.
\end{enumerate}

\section{Guidance for Practitioners}
\label{sec:ch6-guidance}

The results suggest the following for someone who has to deploy a small LLM of this kind on a single GPU.

\begin{itemize}
  \item \textbf{Precision.} Start from weight-only quantization. If the 4-bit model is not good enough, give 8 bits to the most sensitive layers before changing the technique; a simple technique with a per-layer plan came close to the best technique without one.
  \item \textbf{Activations.} Keep the activations at 16 bits if latency in vLLM matters. Models with rounded activations were small but 13 to 17 times slower than FP16 in vLLM.
  \item \textbf{Pruning.} Do not prune for size or speed unless the engine stores or runs sparse weights. If pruning is used, leave the most sensitive layers unpruned.
  \item \textbf{KV cache.} Choose a KV-cache technique for memory, not for speed, and check its effect on task accuracy.
  \item \textbf{Measurement.} Measure perplexity, size and latency in the engine that will serve the model, and include a task accuracy in the requirement.
\end{itemize}

\section{Future Work}
\label{sec:ch6-future}

Several directions follow from the limitations stated in Section~\ref{sec:ch5-limits}.

\begin{itemize}
  \item \textbf{Other models and seeds.} Repeating the experiments on larger models of the Llama family and with several seeds would show whether the findings transfer and how large the variation between runs is.
  \item \textbf{Comparisons at equal size.} The comparison of bit plans at equal size was made for four weight techniques in Hugging Face. Repeating it in the engines, for the other techniques and with other ways of choosing the layers below 4 bits would show how far it holds.
  \item \textbf{Export-aware quantization.} The Hugging Face and the vLLM perplexity differed for the models whose weights are re-rounded during export. Rounding the weights directly onto the grid of the target engine would avoid the second rounding.
  \item \textbf{Accuracy in the requirement.} Adding a task accuracy as a limit or as an objective of the search would address the loss of accuracy seen in every final model.
  \item \textbf{Engines and sparse formats.} Measurements in TensorRT-LLM, and in an engine that stores pruned weights in a sparse format, would complete the comparison across engines and give pruning a fair test of size and speed.
  \item \textbf{KV cache in the engines.} Implementing the KV-cache techniques inside an engine, and testing them with longer contexts and larger batches where the cache is a larger share of the memory, would allow their effect on memory and speed to be measured.
  \item \textbf{Other sensitivity scores.} The framework implements six other scores that were not used in the experiments; comparing the plans they produce is a direct extension.
\end{itemize}
\clearpage
%\include{Chapters/Chapter5v2}
%\include{Chapters/Chapter6v2}


%----------------------------------------------------------------------------------------
%	THESIS CONTENT - APPENDICES
%----------------------------------------------------------------------------------------

\appendix % Cue to tell LaTeX that the following "chapters" are Appendices

% Include the appendices of the thesis as separate files from the Appendices folder
% Uncomment the lines as you write the Appendices

\clearpage
\clearpage
%\include{Appendices/AppendixB}
%\include{Appendices/AppendixC}

%----------------------------------------------------------------------------------------
%	BIBLIOGRAPHY
%----------------------------------------------------------------------------------------
% or another style you prefer like plain, apalike, etc.
   % without .bib extension
\printbibliography



%----------------------------------------------------------------------------------------

\end{document}
